% Leçon 11 — Expression écrite : le recyclage — Chinois Tle A — Cours signé OnBuch+
% Programme officiel MINESEC. Compiler avec : tectonic ZH11-ecrit-le-recyclage.tex
\newcommand{\DOCMATIERE}{Chinois}
\newcommand{\DOCNIVEAU}{Tle A}
\newcommand{\DOCMODULE}{Module 2 : Environnement et santé}
\newcommand{\DOCLECON}{ZH11}
\newcommand{\DOCTITRE}{Leçon 11 — Expression écrite : le recyclage}
% OnBuch+ — préambule « cours signés » ENRICHI (compilé avec tectonic / XeLaTeX)
% Système visuel « Pop » d'OnBuch+ : fond crème (jamais blanc), bordures encre
% épaisses, ombres « dures » décalées, titres Archivo Black en capitales.
% Version MATHÉMATIQUES : graphes (pgfplots/tikz), boîtes pédagogiques
% (définition, propriété, méthode, attention, le savais-tu, exercice résolu,
% exercice, TP, bilan), page de garde et signature OnBuch+ ancrée en bas.
%
% Macros attendues dans main.tex AVANT \input{preamble} :
%   \DOCMATIERE  \DOCNIVEAU  \DOCMODULE  \DOCLECON (numéro)  \DOCTITRE
\documentclass[11pt]{article}

\usepackage{fontspec}
\usepackage[french]{babel} % français = langue principale ; le chinois via \zh{..} / textebox (xeCJK)
\usepackage{xeCJK}
\usepackage{pifont} % \ding{51} \ding{55} utilisés par les modèles
\usepackage[a4paper,margin=2cm,top=3.4cm,bottom=2.8cm,headheight=1.4cm,footskip=1.3cm]{geometry}
\usepackage{amsmath,amssymb,amsfonts}
\usepackage{mathtools} % \xmapsto, \coloneqq... souvent utilisés par les modèles
\usepackage{cancel} % simplifications barrées (bilans de forces)
% \abs{z} (module, valeur absolue) et \norm{u} (norme), utilisés par les modèles
\providecommand{\abs}{}\let\abs\relax\DeclarePairedDelimiter{\abs}{\lvert}{\rvert}
\providecommand{\norm}{}\let\norm\relax\DeclarePairedDelimiter{\norm}{\lVert}{\rVert}
\usepackage[table]{xcolor} % [table] : fournit aussi \rowcolors (alternance de lignes)
\usepackage{pagecolor}
\usepackage{tikz}
\usetikzlibrary{arrows.meta,positioning,calc,shapes.geometric,shapes.misc,shapes.arrows,decorations.pathmorphing,decorations.pathreplacing,decorations.markings,patterns,shadows,mindmap,trees,fit,backgrounds,matrix,intersections,angles,quotes}
\usepackage{pgfplots}
\usepackage[version=4]{mhchem} % équations nucléaires \ce{^{235}_{92}U}
\usepackage[european,siunitx]{circuitikz} % schémas électriques
\pgfplotsset{compat=1.18}
\usepgfplotslibrary{fillbetween,groupplots} % aires entre courbes (calcul intégral)
\usepackage{enumitem}
\usepackage{fancyhdr}
% [most] charge toutes les bibliothèques tcolorbox (listings, minted, raster,
% documentation...) et gonfle inutilement la mémoire TeX sur un document long
% (~300 boîtes) : on ne charge que ce qui est réellement utilisé.
\usepackage[skins,breakable]{tcolorbox}
\usepackage{graphicx}
\usepackage{colortbl}
\usepackage{booktabs}
\usepackage{tabularx}
\usepackage{diagbox} % case d'en-tête diagonale des tableaux à double entrée
% Colonnes extensibles centrée / gauche / droite (C, L, R), souvent utilisées par les modèles
\newcolumntype{C}{>{\centering\arraybackslash}X}
\newcolumntype{L}{>{\raggedright\arraybackslash}X}
\newcolumntype{R}{>{\raggedleft\arraybackslash}X}
\usepackage{array}
\usepackage{multicol}
\usepackage{multirow}
\usepackage{siunitx}
% parse-numbers=false : accepte aussi \SI{2,5 \times 10^{-3}}{...}
\sisetup{locale=FR,output-decimal-marker={,},group-minimum-digits=4,parse-numbers=false}
\DeclareSIUnit\ohm{\ensuremath{\symup{\Omega}}} % U+2126 absent de la police
\DeclareSIUnit\division{div} % réglages d'oscilloscope (ms/div, V/div)
\DeclareSIUnit\tr{tr} % tours (vitesse de rotation, ex. \SI{1500}{\tr\per\minute})
\DeclareSIUnit\molar{M} % concentration molaire (ex. \SI{3}{\molar}, \SI{1}{\micro\molar})
\DeclareSIUnit\an{an} % année (le modèle confond parfois avec \year, un compteur TeX) — ex. \SI{5}{\kilo\meter\per\an}
\DeclareSIUnit\FCFA{FCFA} % franc CFA (ex. \SI{500}{\FCFA\per\kilo\gram})
\DeclareSIUnit\degreeBrix{\textdegree Brix} % teneur en sucre d'un sirop/jus (réfractométrie)
\DeclareSIUnit\month{mois} % le modèle utilise parfois \month, un compteur TeX, comme unité de durée
\usepackage{qrcode}
% \degree : commande fréquemment utilisée par les modèles pour un angle ou une
% température en dehors de \SI{}{} (non fournie par les paquets chargés).
\providecommand{\degree}{^{\circ}}
% \qed (fin de démonstration), utilisé par les modèles sans amsthm
\providecommand{\qed}{\hfill$\square$}
% \barre : double barre des valeurs interdites dans un tableau de variations (tkz-tab)
\providecommand{\barre}{\ensuremath{\|}}
% environnement proof (amsthm non chargé) utilisé par les modèles
\ifdefined\proof\else\newenvironment{proof}[1][Démonstration]{\par\noindent\textit{#1.}\ }{\qed\par}\fi
\usepackage{lastpage}
\usepackage{hyperref}
\hypersetup{hidelinks}

% ============================================================
% Palette « Pop » OnBuch+ — mêmes hex que la classe P (plus_ui.dart)
% ============================================================
\definecolor{poporange}{HTML}{F59321}
\definecolor{popdark}{HTML}{E07A0C}
\definecolor{popcream}{HTML}{FFF6E8}
\definecolor{popink}{HTML}{1C1714}
\definecolor{poppurple}{HTML}{7A5AE0}
\definecolor{popgreen}{HTML}{2E9E6A}
\definecolor{poppink}{HTML}{E0457B}
\definecolor{popblue}{HTML}{2F7DE1}
\definecolor{popgold}{HTML}{C98A12}
\definecolor{popmuted}{HTML}{8A7F76}
\definecolor{popline}{HTML}{EADFD2}
\definecolor{popgray}{HTML}{8A7F76}
\colorlet{popgrey}{popgray}
% Alias tolérants (le modèle nomme parfois les couleurs différemment)
\colorlet{popwhite}{white}
\colorlet{popblack}{popink}
\colorlet{popred}{poppink}
\colorlet{popyellow}{popgold}
% Teintes claires (fonds de tableaux/encadrés) — le modèle nomme parfois
% les couleurs avec un suffixe L (ex. popblueL) sans les avoir définies.
\colorlet{poporangeL}{poporange!20}
\colorlet{popdarkL}{popdark!20}
\colorlet{poppurpleL}{poppurple!20}
\colorlet{popgreenL}{popgreen!20}
\colorlet{poppinkL}{poppink!20}
\colorlet{popblueL}{popblue!20}
\colorlet{popgoldL}{popgold!20}
\colorlet{popredL}{popred!20}
\colorlet{popgrayL}{popgray!20}
% Teintes claires pour fonds de tableaux / zones
\colorlet{popblueL}{popblue!10!white}
\colorlet{poporangeL}{poporange!14!white}
\colorlet{popgreenL}{popgreen!12!white}
\colorlet{poppurpleL}{poppurple!12!white}
\colorlet{poppinkL}{poppink!12!white}
\colorlet{popgoldL}{popgold!14!white}

\pagecolor{popcream}

% ============================================================
% Typographie
% ============================================================
\defaultfontfeatures{Ligatures=TeX}
\setmainfont{PlusJakartaSans-Medium.ttf}[Path=../../../fonts/,BoldFont=PlusJakartaSans-Bold.ttf]
% Symboles, API et emojis absents de la police principale : repli sur DejaVu Sans (caractères actifs)
\newfontface\symfont{DejaVuSans.ttf}[Path=../../../fonts/]
\makeatletter
\newcommand\symactive[1]{\catcode#1=13 \begingroup\lccode`\~=#1 \lowercase{\endgroup\def~}{{\symfont\char#1}}}
\newcommand\symblank[1]{\catcode#1=13 \begingroup\lccode`\~=#1 \lowercase{\endgroup\def~}{}}
\newcommand\symhyph[1]{\catcode#1=13 \begingroup\lccode`\~=#1 \lowercase{\endgroup\def~}{\mbox{-}}}
\newcount\symc
\newcommand\symrange[2]{\symc=#1 \@whilenum\symc<#2 \do{\expandafter\symactive\expandafter{\the\symc}\advance\symc by 1 }}
\newcommand\symblankrange[2]{\symc=#1 \@whilenum\symc<#2 \do{\expandafter\symblank\expandafter{\the\symc}\advance\symc by 1 }}
\makeatother
\symrange{"0250}{"02FF}\symactive{"2713}\symactive{"2714}\symactive{"2717}\symactive{"2718}\symactive{"2610}\symactive{"2611}\symactive{"2612}
\symrange{"2600}{"27BF}
\catcode"2705=13 \begingroup\lccode`\~="2705 \lowercase{\endgroup\def~}{{\symfont\char"2713}}\symblank{"26BD}\symblank{"26A1}
\symrange{"2460}{"257F}\symactive{"223C}\symactive{"2248}\symactive{"2260}
\symhyph{"2011}\symhyph{"2E1A}\symblankrange{"1F300}{"1FAFF}\symblank{"FE0F}
% Caractères chinois : WenQuanYi Zen Hei (sous-ensemble GB2312 + pinyin), gras/italique simulés
\setCJKmainfont{WenQuanYiZenHei-subset.ttf}[Path=../../../fonts/,AutoFakeBold=3,AutoFakeSlant=0.2]
\xeCJKsetup{CJKspace=false}
% Pinyin : voyelles à caron (ǎ ǐ ǒ ǔ ǖ ǘ ǚ ǜ) absentes de la police principale -> caractères actifs
\newfontface\pyfont{WenQuanYiZenHei-subset.ttf}[Path=../../../fonts/]
\newcommand\pyactive[1]{\catcode#1=13 \begingroup\lccode`\~=#1 \lowercase{\endgroup\def~}{{\pyfont\char#1}}}
\pyactive{"01CD}\pyactive{"01CE}\pyactive{"01CF}\pyactive{"01D0}\pyactive{"01D1}\pyactive{"01D2}\pyactive{"01D3}\pyactive{"01D4}\pyactive{"01D5}\pyactive{"01D6}\pyactive{"01D7}\pyactive{"01D8}\pyactive{"01D9}\pyactive{"01DA}\pyactive{"01DB}\pyactive{"01DC}
% Maths en sans-serif (Fira Math) avec chiffres et lettres droites de Plus
% Jakarta, pour que les formules se fondent dans le texte.
\usepackage[math-style=ISO,bold-style=ISO]{unicode-math}
\setmathfont{FiraMath-Regular.otf}[Path=../../../fonts/]
\setmathfont{PlusJakartaSans-Medium.ttf}[Path=../../../fonts/,range={up/num,up/{Latin,latin},\mathup}]
\usepackage{icomma} % virgule décimale sans espace en mode math
% Caractères mathématiques Unicode absents des polices de texte (Plus Jakarta,
% Archivo Black) : rendus par leur équivalent mathématique, en texte comme en
% formule — sinon ils s'affichent en carré vide (ex. « Arithmétique dans ℤ »).
\usepackage{newunicodechar}
\newunicodechar{ℝ}{\ensuremath{\mathbb{R}}}
\newunicodechar{ℤ}{\ensuremath{\mathbb{Z}}}
\newunicodechar{ℂ}{\ensuremath{\mathbb{C}}}
\newunicodechar{ℕ}{\ensuremath{\mathbb{N}}}
\newunicodechar{ℚ}{\ensuremath{\mathbb{Q}}}
\newunicodechar{𝔻}{\ensuremath{\mathbb{D}}}
\newunicodechar{α}{\ensuremath{\alpha}}
\newunicodechar{β}{\ensuremath{\beta}}
\newunicodechar{γ}{\ensuremath{\gamma}}
\newunicodechar{δ}{\ensuremath{\delta}}
\newunicodechar{ε}{\ensuremath{\varepsilon}}
\newunicodechar{θ}{\ensuremath{\theta}}
\newunicodechar{λ}{\ensuremath{\lambda}}
\newunicodechar{μ}{\ensuremath{\mu}}
\newunicodechar{σ}{\ensuremath{\sigma}}
\newunicodechar{τ}{\ensuremath{\tau}}
\newunicodechar{φ}{\ensuremath{\varphi}}
\newunicodechar{ψ}{\ensuremath{\psi}}
\newunicodechar{ω}{\ensuremath{\omega}}
\newunicodechar{Σ}{\ensuremath{\Sigma}}
% majuscules produites par \MakeUppercase dans les titres (α-aminés -> Α)
\newunicodechar{Α}{\ensuremath{\alpha}}
\newunicodechar{Β}{\ensuremath{\beta}}
\newunicodechar{Γ}{\ensuremath{\Gamma}}
\newunicodechar{Δ}{\ensuremath{\Delta}}
\newunicodechar{Ε}{\ensuremath{\varepsilon}}
\newunicodechar{Θ}{\ensuremath{\Theta}}
\newunicodechar{Λ}{\ensuremath{\Lambda}}
\newunicodechar{Μ}{\ensuremath{\mu}}
\newunicodechar{Τ}{\ensuremath{\tau}}
\newunicodechar{Φ}{\ensuremath{\Phi}}
\newunicodechar{Ψ}{\ensuremath{\Psi}}
\newunicodechar{Ω}{\ensuremath{\Omega}}
\newunicodechar{↦}{\ensuremath{\mapsto}}
\newunicodechar{∈}{\ensuremath{\in}}
\newunicodechar{∉}{\ensuremath{\notin}}
\newunicodechar{∧}{\ensuremath{\wedge}}
\newunicodechar{⇔}{\ensuremath{\Leftrightarrow}}
\newunicodechar{⟺}{\ensuremath{\Longleftrightarrow}}
\newunicodechar{⇒}{\ensuremath{\Rightarrow}}
\newunicodechar{⟹}{\ensuremath{\Longrightarrow}}
\newunicodechar{∘}{\ensuremath{\circ}}
\newunicodechar{∪}{\ensuremath{\cup}}
\newunicodechar{∩}{\ensuremath{\cap}}
\newunicodechar{≡}{\ensuremath{\equiv}}
\newunicodechar{⊂}{\ensuremath{\subset}}
\newunicodechar{⊥}{\ensuremath{\perp}}
\newunicodechar{∃}{\ensuremath{\exists}}
\newunicodechar{∀}{\ensuremath{\forall}}
\newunicodechar{∛}{\ensuremath{\sqrt[3]{\,}}}
\newunicodechar{∜}{\ensuremath{\sqrt[4]{\,}}}
\newunicodechar{⃗}{\ensuremath{{}^{\to}}}
\newunicodechar{ⁿ}{\ifmmode^{n}\else\textsuperscript{\ensuremath{n}}\fi}
\newunicodechar{ˣ}{\ifmmode^{x}\else\textsuperscript{\ensuremath{x}}\fi}
\newunicodechar{ᵇ}{\ifmmode^{b}\else\textsuperscript{\ensuremath{b}}\fi}
\newunicodechar{ʳ}{\ifmmode^{r}\else\textsuperscript{\ensuremath{r}}\fi}
\newunicodechar{ᵘ}{\ifmmode^{u}\else\textsuperscript{\ensuremath{u}}\fi}
\newunicodechar{ʸ}{\ifmmode^{y}\else\textsuperscript{\ensuremath{y}}\fi}
\newunicodechar{ᵃ}{\ifmmode^{a}\else\textsuperscript{\ensuremath{a}}\fi}
\newunicodechar{ᵉ}{\ifmmode^{e}\else\textsuperscript{\ensuremath{e}}\fi}
\newunicodechar{ᵏ}{\ifmmode^{k}\else\textsuperscript{\ensuremath{k}}\fi}
\newunicodechar{ᵖ}{\ifmmode^{p}\else\textsuperscript{\ensuremath{p}}\fi}
\newunicodechar{ᴺ}{\ifmmode^{N}\else\textsuperscript{\ensuremath{N}}\fi}
\newunicodechar{ᶜ}{\ifmmode^{c}\else\textsuperscript{\ensuremath{c}}\fi}
\newunicodechar{ʰ}{\ifmmode^{h}\else\textsuperscript{\ensuremath{h}}\fi}
\newunicodechar{ᵠ}{\ifmmode^{q}\else\textsuperscript{\ensuremath{q}}\fi}
\newunicodechar{ₙ}{\ifmmode_{n}\else\textsubscript{\ensuremath{n}}\fi}
\newunicodechar{ₐ}{\ifmmode_{a}\else\textsubscript{\ensuremath{a}}\fi}
\newunicodechar{ᵢ}{\ifmmode_{i}\else\textsubscript{\ensuremath{i}}\fi}
\newunicodechar{ᵣ}{\ifmmode_{r}\else\textsubscript{\ensuremath{r}}\fi}
\newunicodechar{ₖ}{\ifmmode_{k}\else\textsubscript{\ensuremath{k}}\fi}
\newunicodechar{ᵦ}{\ifmmode_{\beta}\else\textsubscript{\ensuremath{\beta}}\fi}
\newunicodechar{ₓ}{\ifmmode_{x}\else\textsubscript{\ensuremath{x}}\fi}
\newfontfamily\popdisplay{ArchivoBlack-Regular.ttf}[Path=../../../fonts/]
\newfontfamily\poplogo{SpaceGrotesk-Bold.ttf}[Path=../../../fonts/]
\newfontfamily\popsemi{PlusJakartaSans-SemiBold.ttf}[Path=../../../fonts/]
\newfontfamily\popxbold{PlusJakartaSans-ExtraBold.ttf}[Path=../../../fonts/]
\newfontfamily\popxboldls{PlusJakartaSans-ExtraBold.ttf}[Path=../../../fonts/,LetterSpace=3.0]

\setlist[enumerate]{leftmargin=1.7em,itemsep=5pt,topsep=4pt}
\setlist[itemize]{leftmargin=1.7em,itemsep=4pt,topsep=4pt,label={\color{poporange}\textbullet}}
\setlength{\parindent}{0pt}
\setlength{\parskip}{5pt}
\renewcommand{\arraystretch}{1.25}

% pgfplots : style Pop par défaut
\pgfplotsset{
  every axis/.append style={axis line style={popink,line width=1pt},tick style={popink},
    grid style={popline,line width=0.5pt},label style={font=\small},tick label style={font=\footnotesize}},
  popcurve/.style={poporange,line width=1.8pt,no markers,smooth},
}
% Même style disponible dans un \draw TikZ simple (hors pgfplots)
\tikzset{popcurve/.style={poporange,line width=1.8pt}}
\tikzset{popgrid/.style={popline,very thin}}
\colorlet{popcurve}{poporange} % le nom du style est parfois utilisé comme couleur

% ============================================================
% Pill / squiggle
% ============================================================
\newcommand{\poppill}[3]{%
  \tikz[baseline=-0.55ex]\node[draw=popink,line width=1.2pt,rounded corners=2.6pt,%
    fill=#2,inner xsep=6.5pt,inner ysep=3pt]{%
    \popxboldls\fontsize{7}{7}\selectfont\color{#3}\MakeUppercase{#1}};%
}
\newcommand{\popsquiggle}[1]{%
  \tikz[baseline=-0.4ex]{
    \draw[#1,line width=2.6pt,line cap=round]
      plot [smooth] coordinates {(0,0.09) (0.34,0.01) (0.68,0.21) (1.02,0.01) (1.36,0.10)};
  }%
}
% Mot-clé surligné (marqueur orange sous le texte)
\newcommand{\cle}[1]{{\popxbold\color{popdark}#1}}

% ============================================================
% En-tête / pied
% ============================================================
\pagestyle{fancy}
\fancyhf{}
\renewcommand{\headrulewidth}{0pt}
\renewcommand{\footrulewidth}{0pt}
\lhead{\raisebox{-0.15ex}{\poplogo\fontsize{13}{13}\selectfont\color{popink}OnBuch\color{poporange}+}}
\rhead{\poppill{\DOCMATIERE\ \textbullet\ \DOCNIVEAU\ \textbullet\ Leçon \DOCLECON}{white}{popink}}
\lfoot{\popxbold\fontsize{7.6}{7.6}\selectfont\color{popmuted}\MakeUppercase{Cours signé OnBuch+}\ \textbullet\ {\popsemi\DOCTITRE}}
\rfoot{\tikz[baseline=-0.6ex]\node[rounded corners=8.5pt,fill=popink,minimum height=17pt,minimum width=17pt,inner xsep=5pt,inner sep=0]{\popxbold\fontsize{8}{8}\selectfont\color{popcream}\thepage\,/\,\pageref*{LastPage}};}

% ============================================================
% Titres
% ============================================================
\newcommand{\chaptitre}[2]{%
  \begin{tcolorbox}[enhanced,colback=poporange,colframe=popink,boxrule=2.6pt,arc=11pt,
      drop shadow=popink,left=15pt,right=15pt,top=12pt,bottom=11pt,before skip=3pt,after skip=18pt]
    \poppill{#2}{popink}{popcream}\par\vspace{9pt}
    {\raggedright\hyphenpenalty=10000\exhyphenpenalty=10000\popdisplay\fontsize{22}{24}\selectfont\color{popcream}\MakeUppercase{#1}\par}
  \end{tcolorbox}
}
% Section numérotée : pastille orange avec le numéro + titre Archivo
\newcounter{popsec}
\newcommand{\coursec}[1]{%
  \stepcounter{popsec}\setcounter{popsub}{0}%
  \par\vspace{16pt}\noindent
  \begin{minipage}{\linewidth}
  \tikz[baseline=-0.7ex]\node[draw=popink,line width=1.6pt,fill=poporange,rounded corners=4pt,minimum size=22pt,inner sep=0]{\popdisplay\fontsize{12}{12}\selectfont\color{popcream}\thepopsec};%
  \hspace{8pt}{\popdisplay\fontsize{15}{17}\selectfont\color{popink}\MakeUppercase{#1}}\par
  \vspace{4pt}\noindent{\color{poporange}\rule{36pt}{3.6pt}}
  \end{minipage}\par\nopagebreak\vspace{8pt}
}
\newcounter{popsub}
\newcommand{\courssub}[1]{%
  \stepcounter{popsub}%
  \par\vspace{9pt}\noindent{\popxbold\fontsize{11.5}{13}\selectfont\color{popink}{\color{poporange}\thepopsec.\thepopsub}\hspace{6pt}#1}\par\nopagebreak\vspace{4pt}
}

% ============================================================
% Boîtes pédagogiques — même recette : fond blanc, bordure épaisse colorée,
% ombre dure encre, badge plein. Titre optionnel [#1].
% ============================================================
\tcbset{popbase/.style={breakable,enhanced,colback=white,boxrule=2.3pt,arc=9pt,
  shadow={3pt}{-3pt}{0pt}{popink},left=14pt,right=14pt,top=11pt,bottom=11pt,before skip=11pt,after skip=15pt,
  fontupper=\popsemi\fontsize{10.6}{14.5}\selectfont\color{popink},
  fontlower=\popsemi\fontsize{10.6}{14.5}\selectfont\color{popink}}}
% Coche dessinée (le glyphe ✓ n'existe pas dans les polices du document)
\newcommand{\popcheck}[1]{\tikz[baseline=-0.5ex]\draw[#1,line width=1.6pt,line cap=round,line join=round](-0.13,0.0)--(-0.04,-0.09)--(0.13,0.1);}
\providecommand{\coche}{\popcheck{popgreen}}
\providecommand{\resultat}[1]{\par\smallskip\noindent\fcolorbox{popgreen}{popgreenL}{\parbox{\dimexpr\linewidth-2\fboxsep-2\fboxrule\relax}{\textbf{Résultat.} #1}}\par\smallskip}
\newcommand{\popboxhead}[3]{\poppill{#1}{#2}{white}\ifx\relax#3\relax\else\hspace{6pt}{\popxbold\fontsize{10.6}{12}\selectfont\color{popink}#3}\fi\par\vspace{7pt}}

\newtcolorbox{aretenir}[1][]{popbase,colframe=popgreen,before upper={\popboxhead{À retenir}{popgreen}{#1}}}
% Texte en chinois (document de lecture, dialogue, extrait) : cadre
% d'étude. Un titre optionnel (source, auteur) s'affiche dans l'en-tête.
\newtcolorbox{textebox}[1][]{popbase,colframe=popblue,before upper={\popboxhead{文本 · Texte}{popblue}{#1}},after upper={}}
% Marques ✓ / ✗ (forme correcte / fautive) souvent utilisées par les modèles
\providecommand{\cross}{\textcolor{poppink}{\ensuremath{\times}}}
\providecommand{\xmark}{\cross}\providecommand{\crossmark}{\cross}\providecommand{\cmark}{\ensuremath{\checkmark}}
% Ligne de réponse / justification à compléter (inventée par les modèles)
\providecommand{\justification}[1]{\par\noindent\textit{Justification~:} \dotfill\par}
% \textipa (package tipa non chargé) : la police ne couvre pas l'API -> simple passage
\providecommand{\textipa}[1]{#1}
% Soulignement (ulem non chargé) : \uline, \ul
\providecommand{\uline}[1]{\underline{#1}}\providecommand{\ul}[1]{\underline{#1}}
% \glossaire{\item ...} inventée par les modèles : liste de glossaire
\providecommand{\glossaire}[1]{\begin{itemize}#1\end{itemize}}
% \alert (beamer) inventée par les modèles : texte mis en évidence
\providecommand{\alert}[1]{\textbf{\textcolor{poppink}{#1}}}
% variantes ASCII d'ordinaux écrites en macro
\providecommand{\iere}{\textsuperscript{re}}\providecommand{\ieme}{\textsuperscript{e}}\providecommand{\ere}{\textsuperscript{re}}\providecommand{\eme}{\textsuperscript{e}}\providecommand{\er}{\textsuperscript{er}}\providecommand{\ie}{\textsuperscript{e}}
% Ordinaux français écrits en macro par les modèles : 1\ière, 3\ième
\newcommand{\ière}{\textsuperscript{re}}\newcommand{\ième}{\textsuperscript{e}}
% \exemple (inventée par les modèles) : étiquette « Exemple : »
\providecommand{\exemple}{\textbf{Exemple~:}\ }
% \square utilisable aussi hors mode math (cases à cocher)
\AtBeginDocument{\let\squareMath\square\renewcommand{\square}{\ensuremath{\squareMath}}}
% Mot ou phrase en chinois à l'intérieur d'un paragraphe français
\newcommand{\zh}[1]{\ifmmode\text{#1}\else{#1}\fi}
\let\eng\zh \let\esp\zh \let\all\zh % alias tolérés
% flèches d'intonation inventées par les modèles
\providecommand{\textsearrow}{}\renewcommand{\textsearrow}{\ensuremath{\searrow}}\providecommand{\textnearrow}{}\renewcommand{\textnearrow}{\ensuremath{\nearrow}}
% \fr{..} : mot français (inventé par les modèles)
\providecommand{\fr}[1]{\textit{#1}}
% zhbox : encadré de texte chinois inventé par les modèles
\newenvironment{zhbox}{\begin{quote}}{\end{quote}}
% \vs : « versus » inventé par les modèles
\providecommand{\vs}{\textit{vs}\ }
% \blank : case à compléter inventée par les modèles
\providecommand{\blank}[1][]{\underline{\hspace{1.6em}}}
\newtcolorbox{exemplebox}[1][]{popbase,colframe=poporange,before upper={\popboxhead{Exemple}{poporange}{#1}}}
\newtcolorbox{definition}[1][]{popbase,colframe=popblue,before upper={\popboxhead{Définition}{popblue}{#1}}}
\newtcolorbox{propriete}[1][]{popbase,colframe=poppurple,before upper={\popboxhead{Propriété}{poppurple}{#1}}}
\newtcolorbox{methode}[1][]{popbase,colframe=popgold,before upper={\popboxhead{Méthode}{popgold}{#1}}}
\newtcolorbox{attention}[1][]{popbase,colframe=poppink,before upper={\popboxhead{Attention}{poppink}{#1}}}
\newtcolorbox{experience}[1][]{popbase,colframe=popblue,colback=popblueL,before upper={\popboxhead{Expérience}{popblue}{#1}}}
% « Le savais-tu ? » : carte violette pleine (contexte, culture, Cameroun)
\newtcolorbox{savaistu}[1][]{popbase,colback=poppurple,colframe=popink,
  fontupper=\popsemi\fontsize{10.4}{14.2}\selectfont\color{white},
  before upper={\poppill{Le savais-tu\,?}{white}{poppurple}\ifx\relax#1\relax\else\hspace{6pt}{\popxbold\color{white}#1}\fi\par\vspace{7pt}}}
% Exercice résolu : énoncé en haut, \tcblower puis la solution sur fond crème
\newcounter{popexo}\newcounter{popexr}
\newtcolorbox{exoresolu}[1][]{popbase,colframe=poporange,colbacklower=popcream,
  segmentation style={popink,line width=1.2pt,dashed},
  before upper={\stepcounter{popexr}\popboxhead{Exercice résolu \thepopexr}{poporange}{#1}},
  before lower={\poppill{Solution}{popink}{popcream}\par\vspace{6pt}}}
% Exercice (non résolu en place) — la correction est dans la partie Corrigés
\newtcolorbox{exercice}[1][]{popbase,colframe=popink,
  before upper={\stepcounter{popexo}\popboxhead{Exercice \thepopexo}{popink}{#1}}}
% Corrigé
\newtcolorbox{corrige}[1][]{popbase,colframe=popgreen,colback=popgreenL,
  before upper={\popboxhead{Corrigé}{popgreen}{#1}}}
% Objectifs / prérequis / bilan
\newtcolorbox{objectifs}{popbase,colframe=popink,colback=poporangeL,
  before upper={\popboxhead{Objectifs de la leçon}{popink}{}}}
\newtcolorbox{prerequis}{popbase,colframe=popink,colback=popblueL,
  before upper={\popboxhead{Prérequis}{popblue}{}}}
\newtcolorbox{bilan}{popbase,colframe=popink,colback=white,boxrule=2.8pt,
  before upper={\popboxhead{Fiche bilan}{poporange}{L'essentiel à connaître pour l'examen}}}
% Figure encadrée avec légende
\newtcolorbox{popfigure}[1][]{enhanced,breakable=false,colback=white,colframe=popline,boxrule=1.4pt,arc=8pt,
  left=10pt,right=10pt,top=10pt,bottom=8pt,before skip=10pt,after skip=12pt,halign=center,
  lower separated=false,halign lower=center,
  fontlower=\popsemi\fontsize{9}{11.5}\selectfont\color{popmuted},
  #1}
\newcommand{\legende}[1]{\par\vspace{6pt}{\popsemi\fontsize{9}{11.5}\selectfont\color{popmuted}#1}}

% ============================================================
% Signature OnBuch+
% ============================================================
\newcommand{\poplogoinline}[1]{{\poplogo\fontsize{#1}{#1}\selectfont\color{popink}OnBuch\color{poporange}+}}
\newcommand{\signatureonbuch}{%
  \begin{tcolorbox}[enhanced,colback=popink,colframe=popink,boxrule=2.2pt,arc=10pt,
      drop shadow=poporange,left=14pt,right=14pt,top=10pt,bottom=10pt,before skip=0pt,after skip=0pt]
    \tikz[baseline=-0.7ex]\node[circle,fill=popgreen,draw=popcream,line width=1.3pt,minimum size=22pt,inner sep=0]{\popcheck{white}};%
    \hspace{9pt}\begin{minipage}[c]{0.62\linewidth}
      {\popxbold\fontsize{12}{13}\selectfont\color{popcream}Signé\ \textbullet\ {\poplogo OnBuch\color{poporange}+}}\par\vspace{2pt}
      {\raggedright\popsemi\fontsize{8}{10.5}\selectfont\color{popline}\MakeUppercase{Équipe pédagogique OnBuch+}\\\MakeUppercase{Conforme au programme officiel MINESEC}\par}
    \end{minipage}\hfill
    {\poplogo\fontsize{22}{22}\selectfont\color{popcream}OnBuch\color{poporange}+}
  \end{tcolorbox}%
}

% ============================================================
% Page de garde
% #1 titre · #2 matière · #3 niveau · #4 module · #5 n° leçon · #6 durée ·
% #7 description / objectifs (texte ou itemize)
% ============================================================
\newcommand{\pagegarde}[7]{%
  \thispagestyle{empty}
  \newgeometry{a4paper,margin=1.7cm,top=1.6cm,bottom=1.4cm}%
  \begin{tikzpicture}[remember picture,overlay]
    \fill[poporange!18] ([xshift=-3.2cm,yshift=-3.6cm]current page.north east) circle (4.6cm);
    \fill[poppurple!14] ([xshift=2.4cm,yshift=7.6cm]current page.south west) circle (3.8cm);
  \end{tikzpicture}%
  {\poplogo\fontsize{30}{30}\selectfont\color{popink}OnBuch\color{poporange}+}\hfill
  \raisebox{0.6ex}{\poppill{Cours signé \textbullet\ Premium}{popink}{popcream}}\par
  \vspace{1pt}\popsquiggle{poppurple}\par
  \vspace{8pt}
  \begin{tcolorbox}[enhanced,colback=poporange,colframe=popink,boxrule=2.8pt,arc=13pt,
      drop shadow=popink,left=18pt,right=18pt,top=12pt,bottom=13pt,after skip=10pt]
    \poppill{#2 \textbullet\ #3}{popink}{popcream}\hspace{5pt}\poppill{#4}{white}{popink}\par\vspace{8pt}
    {\popdisplay\fontsize{14}{14}\selectfont\color{popink}LEÇON #5}\par\vspace{4pt}
    {\raggedright\hyphenpenalty=10000\exhyphenpenalty=10000\popdisplay\fontsize{25}{27}\selectfont\color{popcream}\MakeUppercase{#1}\par}
    \vspace{8pt}
    {\popxbold\fontsize{9.5}{9.5}\selectfont\color{popink}\MakeUppercase{Durée conseillée : #6}}
  \end{tcolorbox}
  \begin{tcolorbox}[enhanced,colback=white,colframe=popink,boxrule=2.2pt,arc=10pt,
      drop shadow=popink,left=16pt,right=16pt,top=10pt,bottom=10pt,after skip=8pt]
    \poppill{Dans cette leçon}{poppurple}{white}\par\vspace{5pt}
    \popsemi\fontsize{9.3}{12.6}\selectfont\color{popink}#7
  \end{tcolorbox}
  \noindent\begin{minipage}[t]{0.487\linewidth}
    \begin{tcolorbox}[enhanced,colback=popgreen,colframe=popink,boxrule=2.2pt,arc=10pt,
        drop shadow=popink,left=11pt,right=11pt,top=10pt,bottom=10pt,height=3.1cm,valign=center]
      \tcbox[colback=white,colframe=popink,boxrule=1.3pt,arc=5pt,boxsep=0pt,nobeforeafter,
        left=3pt,right=3pt,top=3pt,bottom=3pt,valign=center]{\qrcode[height=1.9cm]{https://play.google.com/store/apps/details?id=com.luvvix.onbuch&hl=fr}}%
      \hspace{8pt}\begin{minipage}[c]{0.5\linewidth}
        {\popdisplay\fontsize{11}{12.5}\selectfont\color{white}\MakeUppercase{Télécharge OnBuch}}\par\vspace{4pt}
        {\raggedright\popsemi\fontsize{8.4}{11}\selectfont\color{white}Gratuit \textbullet\ Google Play\\Tuteur IA, annales, concours, résultats\par}
      \end{minipage}
    \end{tcolorbox}
  \end{minipage}\hfill
  \begin{minipage}[t]{0.487\linewidth}
    \begin{tcolorbox}[enhanced,colback=poppurple,colframe=popink,boxrule=2.2pt,arc=10pt,
        drop shadow=popink,left=11pt,right=11pt,top=10pt,bottom=10pt,height=3.1cm,valign=center]
      \tikz[baseline=-0.7ex]\node[circle,fill=white,draw=popink,line width=1.3pt,minimum size=1.6cm,inner sep=0]{\popdisplay\fontsize{24}{24}\selectfont\color{poppurple}+};%
      \hspace{8pt}\begin{minipage}[c]{0.6\linewidth}
        {\popdisplay\fontsize{11}{12.5}\selectfont\color{white}\MakeUppercase{Passe à OnBuch+}}\par\vspace{4pt}
        {\raggedright\popsemi\fontsize{8.4}{11}\selectfont\color{white}Tous les cours signés, Tuteur IA illimité, fiches PDF — onglet OnBuch+ de l'app\par}
      \end{minipage}
    \end{tcolorbox}
  \end{minipage}
  \vspace{8pt}
  \popcontenu
  \vfill
  \signatureonbuch
  \restoregeometry
  \newpage
}

% Rangée « Ce cours contient » (page de garde)
\newcommand{\poptile}[3]{% #1 couleur #2 gros texte #3 légende
  \begin{tcolorbox}[enhanced,colback=#1,colframe=popink,boxrule=2pt,arc=9pt,shadow={2.5pt}{-2.5pt}{0pt}{popink},
      left=6pt,right=6pt,top=9pt,bottom=9pt,halign=center,valign=center,height=1.75cm,width=\linewidth,nobeforeafter]
    {\popdisplay\fontsize{12.5}{13}\selectfont\color{white}\MakeUppercase{#2}}\par\vspace{3pt}
    {\centering\popsemi\fontsize{7.8}{9.5}\selectfont\color{white}#3\par}
  \end{tcolorbox}}
\newcommand{\popcontenu}{%
  \par\noindent{\popxboldls\fontsize{7.5}{7.5}\selectfont\color{popmuted}CE COURS SIGNÉ CONTIENT}\par\vspace{7pt}
  \noindent\begin{minipage}[t]{0.235\linewidth}\poptile{popblue}{Cours}{complet, illustré, pas à pas}\end{minipage}\hfill
  \begin{minipage}[t]{0.235\linewidth}\poptile{poporange}{Méthodes}{et exercices résolus}\end{minipage}\hfill
  \begin{minipage}[t]{0.235\linewidth}\poptile{poppink}{Exercices}{type examen + corrigés}\end{minipage}\hfill
  \begin{minipage}[t]{0.235\linewidth}\poptile{popgreen}{Fiche bilan}{l'essentiel à retenir}\end{minipage}\par}

% Sommaire
\newtcolorbox{sommaire}{enhanced,colback=white,colframe=popink,boxrule=2.2pt,arc=10pt,
  drop shadow=popink,left=18pt,right=18pt,top=13pt,bottom=13pt,before skip=6pt,after skip=20pt,
  before upper={\poppill{Sommaire}{poppurple}{white}\par\vspace{9pt}\popsemi\fontsize{10.6}{14.5}\selectfont\color{popink}}}

% Signature de fin de cours — poussée tout en bas de la dernière page
\newcommand{\findecours}{%
  \par\vfill\nopagebreak
  \begin{center}\popsquiggle{poporange}\end{center}\vspace{4pt}
  \signatureonbuch
}
\AtBeginDocument{\def\llbracket{\mathopen{[\mkern-3.5mu[}}\def\rrbracket{\mathclose{]\mkern-3.5mu]}}} % intervalles d'entiers (glyphe ⟦ absent de la police)
% Glyphes absents de Fira Math : redéfinis à partir de glyphes disponibles
\AtBeginDocument{%
  \def\setminus{\mathbin{\backslash}}\let\smallsetminus\setminus
  \def\vdots{\mathord{\vbox{\baselineskip3.2pt\lineskiplimit0pt\kern4pt\hbox{.}\hbox{.}\hbox{.}}}}%
  \def\ddots{\mathinner{\mkern1mu\raise6.5pt\hbox{.}\mkern2mu\raise3.6pt\hbox{.}\mkern2mu\raise0.7pt\hbox{.}\mkern1mu}}%
  \def\triangle{\mathord{\scalebox{0.9}{$\symup{\Delta}$}}}%
  \def\nabla{\mathord{\raisebox{\depth}{\scalebox{1}[-1]{$\symup{\Delta}$}}}}%
  \def\checkmark{\popcheck{popdark}}%
  \def\star{\mathbin{\ast}}\def\bigstar{\mathord{\ast}}%
  \def\diamond{\mathbin{\scalebox{0.6}{$\Diamond$}}}%
  \def\bigcup{\mathop{\mathchoice{\raisebox{-0.3em}{\scalebox{1.7}{$\cup$}}}{\scalebox{1.25}{$\cup$}}{\cup}{\cup}}\displaylimits}%
  \def\bigcap{\mathop{\mathchoice{\raisebox{-0.3em}{\scalebox{1.7}{$\cap$}}}{\scalebox{1.25}{$\cap$}}{\cap}{\cap}}\displaylimits}%
  \def\lesseqqgtr{\mathrel{\lessgtr}}\def\bigoplus{\mathop{\oplus}\displaylimits}%
}
\providecommand{\ssi}{si et seulement si} % abréviation fréquente
\AtBeginDocument{\providecommand{\wideparen}{\overparen}} % arc de cercle

%% FIGFIX-BEGIN v5 — correctifs globaux des figures (docs/onbuch-generation/tools/figfix)
\usepackage{adjustbox}\usepackage{etoolbox}\tcbuselibrary{raster}
% toute figure TikZ/pgfplots plus large que la ligne est réduite à la largeur disponible
\BeforeBeginEnvironment{tikzpicture}{\begin{adjustbox}{max width=\linewidth}}
\AfterEndEnvironment{tikzpicture}{\end{adjustbox}}
% légendes pgfplots lisibles par-dessus les courbes
\pgfplotsset{every axis/.append style={legend style={fill=white,fill opacity=0.92,text opacity=1,draw=black!15,inner sep=3pt}}}
% étiquettes d'arêtes (syntaxe edge["..."]) avec fond blanc
\tikzset{every edge quotes/.append style={fill=white,inner sep=1.2pt}}
%% FIGFIX-END
\begin{document}

\pagegarde{Leçon 11 — Expression écrite : le recyclage}{Chinois}{Tle A}{Module 2 : Environnement et santé}{ZH11}{2 h}{Cette leçon d'expression écrite entraîne les élèves à rédiger un texte simple en chinois sur le recyclage, à écrire correctement les caractères du thème et à pratiquer la traduction (thème et version). Elle fournit un modèle de texte commenté, des connecteurs utiles et une grille d'évaluation.
\vspace{4pt}
\begin{multicols}{2}\small
\begin{itemize}
\item À la fin de cette leçon, je sais écrire correctement les caractères clés du thème du recyclage (traits, radicaux, caractères proches).
\item À la fin de cette leçon, je sais distinguer les caractères proches (圾/级, 环/坏, 收/放) pour éviter les erreurs d'écriture.
\item À la fin de cette leçon, je sais utiliser les connecteurs 首先、然后、最后、因为、所以 pour structurer un texte.
\item À la fin de cette leçon, je sais rédiger un texte simple de 5 à 8 phrases sur le recyclage en chinois.
\item À la fin de cette leçon, je sais traduire du français vers le chinois des phrases simples sur l'environnement (thème).
\item À la fin de cette leçon, je sais traduire du chinois vers le français un texte simple sur le recyclage (version).
\end{itemize}\end{multicols}}

\begin{sommaire}
\begin{multicols}{2}
\begin{enumerate}
\item Modèle de texte commenté
\item Structure et connecteurs du texte
\item Écriture correcte des caractères du thème
\item Thème : français → chinois
\item Version : chinois → français
\item Production guidée et grille d'évaluation
\item Activité d'intégration
\item Méthodes et exercices résolus
\item Exercices
\item Corrigés des exercices
\item Fiche bilan
\end{enumerate}
\end{multicols}
\end{sommaire}

\begin{objectifs}
\begin{itemize}
\item À la fin de cette leçon, je sais écrire correctement les caractères clés du thème du recyclage (traits, radicaux, caractères proches).
\item À la fin de cette leçon, je sais distinguer les caractères proches (圾/级, 环/坏, 收/放) pour éviter les erreurs d'écriture.
\item À la fin de cette leçon, je sais utiliser les connecteurs 首先、然后、最后、因为、所以 pour structurer un texte.
\item À la fin de cette leçon, je sais rédiger un texte simple de 5 à 8 phrases sur le recyclage en chinois.
\item À la fin de cette leçon, je sais traduire du français vers le chinois des phrases simples sur l'environnement (thème).
\item À la fin de cette leçon, je sais traduire du chinois vers le français un texte simple sur le recyclage (version).
\item À la fin de cette leçon, je sais employer les structures 应该、可以、把 pour exprimer le devoir et la possibilité dans un texte.
\item À la fin de cette leçon, je sais m'auto-évaluer à l'aide d'une grille (caractères, grammaire, connecteurs, sens).
\end{itemize}
\end{objectifs}

\begin{prerequis}
\begin{itemize}
\item Vocabulaire de l'environnement vu en Module 2 (环境, 污染, 垃圾, 保护)
\item Verbes modaux 应该, 可以, 会 (Première)
\item Structure de base sujet-verbe-objet et particule 了
\item Radicaux de base (土, 木, 氵, 扌) et ordre général des traits
\item Connecteurs simples 和, 但是, 因为...所以...
\end{itemize}
\end{prerequis}

\newpage

\coursec{Leçon 11 — Expression écrite : le recyclage}

\noindent Module 2 : Environnement et santé — Durée indicative : 2 h \quad Série : Tle A

\noindent À Yaoundé comme à Douala, les montagnes de plastique envahissent les rues après la pluie. Au marché Mokolo, une coopérative de jeunes trie les bouteilles pour les revendre. En Chine, le tri sélectif (\zh{垃圾分类}) est devenu obligatoire dans les grandes villes comme Shanghai depuis 2019. Comment raconter tout cela en chinois ? Aujourd'hui, vous allez apprendre à écrire un texte simple sur le recyclage, à maîtriser les caractères du thème et à traduire des phrases dans les deux sens (thème / version).

\courssub{Modèle de texte commenté}

\begin{textebox}[Texte modèle : le recyclage]
我们应该保护环境。首先，我们要把垃圾分类：塑料、纸和玻璃要分开放。然后，旧报纸和瓶子可以回收利用。最后，少用塑料袋，多用布袋。因为地球只有一个，所以保护环境就是保护我们自己。
\end{textebox}

\noindent \textbf{Lecture guidée (pinyin phrase par phrase) :}
\begin{itemize}
    \item \zh{我们应该保护环境。} \emph{Wǒmen yīnggāi bǎohù huánjìng.}
    \item \zh{首先，我们要把垃圾分类：塑料、纸和玻璃要分开放。} \emph{Shǒuxiān, wǒmen yào bǎ lājī fēnlèi: sùliào, zhǐ hé bōlí yào fēnkāi fàng.}
    \item \zh{然后，旧报纸和瓶子可以回收利用。} \emph{Ránhòu, jiù bàozhǐ hé píngzi kěyǐ huíshōu lìyòng.}
    \item \zh{最后，少用塑料袋，多用布袋。} \emph{Zuìhòu, shǎo yòng sùliào dài, duō yòng bùdài.}
    \item \zh{因为地球只有一个，所以保护环境就是保护我们自己。} \emph{Yīnwèi dìqiú zhǐ yǒu yīgè, suǒyǐ bǎohù huánjìng jiùshì bǎohù wǒmen zìjǐ.}
\end{itemize}

\noindent \textbf{Traduction littérale (mot à mot) :}
\begin{quote}
Nous devons protéger environnement. Premièrement, nous voulons \textbf{ba} déchets classer : plastique, papier et verre veulent séparément placer. Ensuite, vieux journaux et bouteilles peuvent recycler utiliser. Finalement, peu utiliser sacs plastique, beaucoup utiliser sacs tissu. Parce que terre seulement avoir un, donc protéger environnement c'est protéger nous-mêmes.
\end{quote}

\noindent \textbf{Traduction naturelle (français correct) :}
\begin{quote}
Nous devons protéger l'environnement. D'abord, il faut trier les déchets : le plastique, le papier et le verre doivent être mis séparément. Ensuite, les vieux journaux et les bouteilles peuvent être recyclés. Enfin, utilisons moins de sacs en plastique et plus de sacs en tissu. Car la Terre n'est qu'une seule, donc protéger l'environnement, c'est nous protéger nous-mêmes.
\end{quote}

\courssub{Structure et connecteurs du texte}

\noindent Le texte suit une structure en \textbf{trois temps} classique de l'argumentation simple : Introduction $\rightarrow$ Actions séquentielles $\rightarrow$ Conclusion causale.

\begin{center}
\begin{tikzpicture}[
    node distance=1.1cm and 1.2cm,
    box/.style={rectangle, draw=popblue, thick, fill=popblueL, rounded corners, minimum width=3.5cm, minimum height=1.0cm, align=center, font=\small},
    boxg/.style={rectangle, draw=popgreen, thick, fill=popgreenL, rounded corners, minimum width=3.5cm, minimum height=1.0cm, align=center, font=\small},
    arrow/.style={-Stealth, thick, popdark}
]
\node[box, align=center] (intro){Introduction\\ \zh{我们应该保护环境}\\ Thèse / Problème};
\node[box, below=of intro, align=center] (a1){\textbf{首先} (D'abord)\\ \zh{把垃圾分类}\\ Action 1 : Tri};
\node[box, below=of a1, align=center] (a2){\textbf{然后} (Ensuite)\\ \zh{回收利用}\\ Action 2 : Recyclage};
\node[box, below=of a2, align=center] (a3){\textbf{最后} (Enfin)\\ \zh{少用塑料袋, 多用布袋}\\ Action 3 : Réduction};
\node[boxg, below=of a3, align=center] (concl){Conclusion\\ \zh{因为...所以...}\\ Justification cause-effet};
\draw[arrow] (intro) -- (a1);
\draw[arrow] (a1) -- (a2);
\draw[arrow] (a2) -- (a3);
\draw[arrow] (a3) -- (concl);
\end{tikzpicture}
\legende{Schéma de la structure du texte modèle}
\end{center}

\begin{popfigure}
\begin{center}\tcbox[colback=popink,colframe=popink,coltext=white,arc=9pt,boxsep=4pt,left=6pt,right=6pt]{\bfseries\small 回收 huíshōu}\end{center}
\vspace{-2pt}
\begin{tcbraster}[raster columns=2,raster equal height=rows,raster column skip=6pt,raster row skip=6pt,enhanced,blankest]
\begin{tcolorbox}[enhanced,colback=poporange!70,colframe=poporange!70,coltext=white,boxrule=0.9pt,arc=3pt,left=4pt,right=4pt,top=3pt,bottom=3pt,boxsep=1pt,halign=left]\bfseries\scriptsize 分类 fēnlèi\end{tcolorbox}
\begin{tcolorbox}[enhanced,colback=popgreen!70,colframe=popgreen!70,coltext=white,boxrule=0.9pt,arc=3pt,left=4pt,right=4pt,top=3pt,bottom=3pt,boxsep=1pt,halign=left]\bfseries\scriptsize 利用 lìyòng\end{tcolorbox}
\begin{tcolorbox}[enhanced,colback=poppurple!70,colframe=poppurple!70,coltext=white,boxrule=0.9pt,arc=3pt,left=4pt,right=4pt,top=3pt,bottom=3pt,boxsep=1pt,halign=left]\bfseries\scriptsize 减少 jiǎnshǎo\end{tcolorbox}
\begin{tcolorbox}[enhanced,colback=poppink!70,colframe=poppink!70,coltext=white,boxrule=0.9pt,arc=3pt,left=4pt,right=4pt,top=3pt,bottom=3pt,boxsep=1pt,halign=left]\bfseries\scriptsize 保护 bǎohù\end{tcolorbox}
\end{tcbraster}

\legende{Carte mentale lexicale autour de \zh{回收} (recyclage)}
\end{popfigure}

\noindent \textbf{Tableau des connecteurs logiques (indispensables pour l'écrit) :}

\begin{center}
\begin{tabularx}{\linewidth}{|X|X|X|X|}
\hline
\rowcolor{popblueL} \textbf{Connecteur} & \textbf{Pinyin} & \textbf{Fonction} & \textbf{Exemple du texte} \\ \hline
首先 & shǒuxiān & Première étape (temporel/logique) & \zh{首先，我们要把垃圾分类} \\ \hline
然后 & ránhòu & Étape suivante, suite & \zh{然后，旧报纸和瓶子可以回收利用} \\ \hline
最后 & zuìhòu & Étape finale, conclusion d'énumération & \zh{最后，少用塑料袋，多用布袋} \\ \hline
因为...所以... & yīnwèi... suǒyǐ... & Cause (因为) $\rightarrow$ Conséquence (所以) & \zh{因为地球只有一个，所以保护环境就是保护我们自己} \\ \hline
\end{tabularx}
\end{center}

\begin{attention}[Règles de ponctuation pour les connecteurs]
\begin{itemize}
    \item \zh{首先}，\zh{然后}，\zh{最后} sont suivis d'une virgule chinoise \zh{，} quand ils sont en tête de phrase.
    \item La proposition introduite par \zh{因为} est \textbf{toujours} suivie d'une virgule \zh{，} avant \zh{所以}. Ex : \zh{因为地球只有一个，所以...}
    \item \zh{所以} n'est \textbf{jamais} placé en début de phrase avant \zh{因为}. Ordre fixe : Cause $\rightarrow$ Conséquence.
\end{itemize}
\end{attention}

\courssub{Vocabulaire clé du thème (HSK 3--4)}

\begin{center}
\begin{tabularx}{\linewidth}{|X|X|X|X|}
\hline
\rowcolor{popblueL} \textbf{Caractères} & \textbf{Pinyin} & \textbf{Nature} & \textbf{Traduction} \\ \hline
环境 & huánjìng & nom & environnement \\ \hline
垃圾 & lājī & nom & déchets, ordures \\ \hline
分类 & fēnlèi & verbe / nom & trier / tri, classification \\ \hline
塑料 & sùliào & nom & plastique \\ \hline
纸 & zhǐ & nom & papier \\ \hline
玻璃 & bōlí & nom & verre \\ \hline
分开放 & fēnkāi fàng & locution verbale & mettre séparément, mettre à part \\ \hline
旧 & jiù & adjectif & vieux, usagé, ancien \\ \hline
报纸 & bàozhǐ & nom & journal (quotidien) \\ \hline
瓶子 & píngzi & nom & bouteille \\ \hline
回收 & huíshōu & verbe & recycler, récupérer \\ \hline
利用 & lìyòng & verbe / nom & utiliser, exploiter / utilisation \\ \hline
回收利用 & huíshōu lìyòng & verbe & recycler (et réutiliser) \\ \hline
塑料袋 & sùliào dài & nom & sac en plastique \\ \hline
布袋 & bùdài & nom & sac en tissu \\ \hline
地球 & dìqiú & nom & Terre (planète) \\ \hline
保护 & bǎohù & verbe & protéger \\ \hline
自己 & zìjǐ & pronom réfléchi & soi-même, nous-mêmes \\ \hline
\end{tabularx}
\end{center}

\courssub{Écriture correcte des caractères du thème}

\noindent Maîtriser l'ordre des traits (\zh{笔顺} bǐshùn) et les radicaux (\zh{部首} bùshǒu) est indispensable pour écrire vite et bien, et pour utiliser le dictionnaire.

\begin{propriete}[Points d'écriture : caractères du thème]
\noindent \textbf{1. 环 (huán) -- Environnement :} Radical \zh{王} (jade/roi, 4 traits). Structure : haut-bas. Ordre : \zh{王} (horizontal, horizontal, vertical, horizontal) $\rightarrow$ \zh{彐} (horizontal, crochet vertical, horizontal) $\rightarrow$ fermeture bas. \textit{Attention : ne pas confondre avec \zh{还} (huán/hái, radical \zh{辶} marche).}
\vspace{0.3cm}

\noindent \textbf{2. 垃 (lā) -- Déchet :} Radical \zh{土} (terre, 3 traits). Phonétique \zh{立} (debout). Structure : gauche-droite. Ordre : \zh{土} (horizontal, vertical, horizontal) $\rightarrow$ \zh{立} (horizontal, vertical, horizontal, horizontal).
\vspace{0.3cm}

\noindent \textbf{3. 类 (lèi) -- Catégorie / Trier :} Radical \zh{米} (riz, 6 traits). Structure : haut-bas. Ordre : \zh{米} $\rightarrow$ \zh{大} (grand) $\rightarrow$ trait final long (\zh{乀}). \textit{Caractère proche : \zh{米} (mǐ) vs \zh{类} (lèi).}
\vspace{0.3cm}

\noindent \textbf{4. 塑 (sù) -- Plastique :} Radical \zh{土} (terre). Phonétique \zh{溯} (remonter). Structure : gauche-droite. \textit{Ne pas confondre avec \zh{俗} (sú, coutume, radical \zh{人}).}
\vspace{0.3cm}

\noindent \textbf{5. 玻 (bō) / 璃 (lí) -- Verre :} Tous deux radical \zh{王} (jade). \zh{玻} : phonétique \zh{皮} (peau). \zh{璃} : phonétique \zh{离} (quitter). S'écrivent toujours ensemble \zh{玻璃}.
\vspace{0.3cm}

\noindent \textbf{6. 回 (huí) -- Retourner / Recycler :} Radical \zh{囗} (enclos, 3 traits). Structure : enclos. Ordre : \zh{囗} (vertical, haut horizontal + angle, bas horizontal) $\rightarrow$ intérieur \zh{口} (bouche). \textit{Piège : ne pas fermer l'enclos avant d'écrire l'intérieur.}
\vspace{0.3cm}

\noindent \textbf{7. 利 (lì) -- Avantage / Tranchant / Utiliser :} Radical \zh{刂} (couteau, 2 traits, à droite). Haut : \zh{禾} (riz). Ordre : \zh{禾} $\rightarrow$ \zh{刂} (vertical, vertical court). \textit{Caractère proche : \zh{丽} (lì, beau, radical \zh{一} un / \zh{鹿} cerf simplifié). Différence : \zh{刂} (couteau) vs \zh{一/鹿}.}
\vspace{0.3cm}

\noindent \textbf{8. 护 (hù) -- Protéger :} Radical \zh{扌} (main, 3 traits). Phonétique \zh{胡}. Structure : gauche-droite. Ordre : \zh{扌} (horizontal, crochet vertical, point) $\rightarrow$ \zh{胡} (lune \zh{月} + ancien \zh{古}).
\end{propriete}

\begin{exercice}[Entraînement à l'écriture -- 10 min]
Écrivez chaque caractère ci-dessous \textbf{5 fois} en respectant l'ordre des traits, puis écrivez le mot complet 2 fois.
\begin{enumerate}
    \item \zh{环境} (huánjìng) \quad \item \zh{垃圾} (lājī) \quad \item \zh{分类} (fēnlèi) \quad \item \zh{塑料} (sùliào)
    \item \zh{玻璃} (bōlí) \quad \item \zh{回收} (huíshōu) \quad \item \zh{利用} (lìyòng) \quad \item \zh{保护} (bǎohù)
\end{enumerate}
\textbf{Autodictée :} Fermez le cahier, dictez-vous les 8 mots en français, écrivez les caractères + pinyin + tons.
\end{exercice}

\courssub{Structures grammaticales réemployées (Révision HSK 3)}

\begin{propriete}[Structure \texorpdfstring{\zh{应该 + V}}{应该 + V} : Obligation morale / Conseil / « Il faut »]
\noindent \textbf{Forme :} Sujet + \zh{应该} (yīnggāi) + Verbe \quad | \quad Négation : Sujet + \zh{不应该} + Verbe

\noindent \textbf{Emploi :} Exprime un devoir moral, une recommandation forte, un « il faut » impersonnel ou personnel. Invariable (pas de conjugaison). Plus fort que \zh{要} (vouloir/devoir futur), moins contraignant que \zh{必须} (bìxū, obligation stricte/règlement).

\noindent \textbf{Exemples :}
\begin{itemize}
    \item \zh{我们应该保护环境。} \emph{Wǒmen yīnggāi bǎohù huánjìng.} — Nous devons / Il faut que nous protégions l'environnement.
    \item \zh{学生不应该乱扔垃圾。} \emph{Xuéshēng bù yīnggāi luàn rēng lājī.} — Les élèves ne doivent pas jeter les déchets n'importe où.
    \item \zh{你应该多喝水。} \emph{Nǐ yīnggāi duō hē shuǐ.} — Tu devrais boire plus d'eau.
\end{itemize}

\noindent \textbf{Comparaison FR/CN :} FR : « devoir » conjugué + infinitif / « il faut que » + subjonctif. CN : \zh{应该} invariable + V. Pas de subjonctif.
\end{propriete}

\begin{propriete}[Structure \texorpdfstring{\zh{可以 + V}}{可以 + V} : Possibilité / Permission / Capacité objective]
\noindent \textbf{Forme :} Sujet + \zh{可以} (kěyǐ) + Verbe \quad | \quad Négation : \zh{不可以} (interdiction) / \zh{不能} (impossibilité/incapacité).

\noindent \textbf{Emploi :} 1) Permission (« avoir le droit de »). 2) Possibilité objective (« il est possible de »). 3) Capacité/fonction (« pouvoir » au sens capacité technique).

\noindent \textbf{Exemples :}
\begin{itemize}
    \item \zh{瓶子可以回收利用。} \emph{Píngzi kěyǐ huíshōu lìyòng.} — Les bouteilles peuvent être recyclées (possibilité objective).
    \item \zh{这里可以扔垃圾吗？} \emph{Zhèlǐ kěyǐ rēng lājī ma?} — Peut-on jeter des déchets ici ? (permission).
    \item \zh{这个软件可以翻译中文。} \emph{Zhège ruǎnjiàn kěyǐ fānyì Zhōngwén.} — Ce logiciel peut traduire le chinois (fonction).
\end{itemize}

\noindent \textbf{Distinction cruciale (piège francophone) :}
\begin{center}
\begin{tabularx}{\linewidth}{|X|X|X|}
\hline
\rowcolor{popblueL} \textbf{Mot} & \textbf{Sens principal} & \textbf{Exemple} \\ \hline
\zh{可以} (kěyǐ) & Permission / Possibilité objective & \zh{我可以去吗？} (Puis-je y aller ?) \\ \hline
\zh{会} (huì) & Compétence acquise, savoir-faire & \zh{我会说中文。} (Je sais parler chinois.) \\ \hline
\zh{能} (néng) & Capacité physique / Circonstancielle & \zh{我今天不能去。} (Je ne peux pas y aller aujourd'hui.) \\ \hline
\end{tabularx}
\end{center}
\end{propriete}

\begin{propriete}[Structure \texorpdfstring{\zh{把 + Objet + Verbe + Complément}}{把 + Objet + Verbe + Complément} : Disposition / Manipulation de l'objet]
\noindent \textbf{Forme :} Sujet + \zh{把} (bǎ) + Objet \textbf{défini/connu/spécifique} + Verbe + \textbf{Complément obligatoire} (résultat, direction, \zh{了}, \zh{完}, \zh{好}, \zh{错}, \zh{在+Lieu}, \zh{给+Quelqu'un}...)

\noindent \textbf{Emploi :} Met l'accent sur l'action effectuée \textbf{sur l'objet} et son devenir/résultat. L'objet est mis en avant (topicalisé). Le verbe \textbf{ne peut pas rester nu}.

\noindent \textbf{Exemples du texte \& variations :}
\begin{itemize}
    \item \zh{我们要把垃圾分类。} \emph{Wǒmen yào bǎ lājī fēnlèi.} — Nous devons trier les déchets. (\zh{分类} verbe composé V-O agissant comme résultat).
    \item \zh{请把瓶子放进蓝色的桶。} \emph{Qǐng bǎ píngzi fàng jìn lánsè de tǒng.} — Mettez les bouteilles dans la poubelle bleue (complément de direction \zh{进...}).
    \item \zh{我已经把旧报纸卖\underline{了}。} \emph{Wǒ yǐjīng bǎ jiù bàozhǐ mài \underline{le}.} — J'ai déjà vendu les vieux journaux (aspect accompli \zh{了}).
    \item \zh{请把垃圾分\underline{好}类。} \emph{Qǐng bǎ lājī fēn \underline{hǎo} lèi.} — Triez bien les déchets (complément de résultat \zh{好} inséré dans le verbe composé).
\end{itemize}

\noindent \textbf{PIÈGES FRANCOPHONES MAJEURS :}
\begin{itemize}
    \item \textbf{Objet indéfini interdit :} \zh{*我把一个苹果吃了} $\rightarrow$ \zh{我吃了一个苹果} (pas de \zh{把} si objet nouveau/indéfini).
    \item \textbf{Verbe nu interdit :} \zh{*我把书看} $\rightarrow$ \zh{我把书看\underline{完了}} / \zh{我把书看\underline{了}}.
    \item \textbf{Négation :} \zh{不} / \zh{没} se place \textbf{avant} \zh{把}. Ex : \zh{我没把垃圾分类} (Je n'ai pas trié les déchets).
\end{itemize}
\end{propriete}

\courssub{Méthode : Analyser et réutiliser un texte modèle}

\begin{methode}[Comment analyser et réutiliser un texte modèle (5 étapes)]
\begin{enumerate}
    \item \textbf{Lire à voix haute} en suivant le pinyin : travaillez les tons, repérez les groupes rythmiques (pause après \zh{首先}，\zh{然后}，\zh{最后}).
    \item \textbf{Traduire phrase par phrase} : d'abord littéral (mot à mot), puis naturel (français correct). Notez les écarts de structure (ex: \zh{把}, \zh{因为...所以...}).
    \item \textbf{Souligner les connecteurs logiques} : \zh{首先 / 然后 / 最后} (séquence), \zh{因为 / 所以} (cause-conséquence). Ce sont les « charnières » du texte.
    \item \textbf{Encadrer les structures grammaticales cibles} : \zh{应该 + V}, \zh{可以 + V}, \zh{把 + O + V + Compl}. Repérez leur place dans la phrase.
    \item \textbf{Construire votre « banque de phrases »} : pour chaque structure, écrivez 2 phrases personnelles (ex: \zh{我应该少用塑料袋}, \zh{纸可以回收利用}, \zh{我要把瓶子洗干净}).
    \item \textbf{Produire votre propre texte} en réutilisant le même squelette : Introduction $\rightarrow$ 3 actions ($\zh{首先/然后/最后}$) $\rightarrow$ Conclusion (\zh{因为...所以...}).
\end{enumerate}
\end{methode}

\courssub{Exemples traduits et variations (Banque de phrases)}

\begin{exemplebox}[Variations sur les structures cibles -- À mémoriser]
\begin{tabularx}{\linewidth}{|X|X|}
\hline
\rowcolor{popblueL} \textbf{Chinois (caractères + pinyin)} & \textbf{Français} \\ \hline
\zh{我们应该保护环境。} \emph{Wǒmen yīnggāi bǎohù huánjìng.} & Nous devons protéger l'environnement. \\ \hline
\zh{每个人都应该分类垃圾。} \emph{Měigè rén dōu yīnggāi fēnlèi lājī.} & Chacun devrait trier ses déchets. \\ \hline
\zh{瓶子可以回收利用。} \emph{Píngzi kěyǐ huíshōu lìyòng.} & Les bouteilles peuvent être recyclées. \\ \hline
\zh{这些塑料袋可以再利用。} \emph{Zhèxiē sùliào dài kěyǐ zài lìyòng.} & Ces sacs en plastique peuvent être réutilisés. \\ \hline
\zh{请把纸放进蓝色的箱子。} \emph{Qǐng bǎ zhǐ fàng jìn lánsè de xiāngzi.} & Mettez le papier dans la boîte bleue. \\ \hline
\zh{我已经把旧衣服捐了。} \emph{Wǒ yǐjīng bǎ jiù yīfu juān le.} & J'ai déjà donné mes vieux vêtements. \\ \hline
\zh{因为污染严重，所以我们要行动。} \emph{Yīnwèi wūrǎn yánzhòng, suǒyǐ wǒmen yào xíngdòng.} & Parce que la pollution est grave, il faut agir. \\ \hline
\end{tabularx}
\end{exemplebox}

\courssub{Thème : Français → Chinois (Entraînement guidé)}

\noindent Traduisez les phrases suivantes en chinois (caractères + pinyin). Utilisez le vocabulaire et les structures de la leçon.

\begin{exercice}[Thème -- Phrases isolées]
\begin{enumerate}
    \item Nous devons trier les déchets. (Structure \zh{应该} + \zh{把})
    \item Les vieux journaux peuvent être recyclés. (Structure \zh{可以})
    \item Mettez les bouteilles en plastique dans la poubelle bleue. (Structure \zh{把} + direction)
    \item Parce que la Terre n'est qu'une, nous devons la protéger. (Structure \zh{因为...所以...} + \zh{保护})
    \item D'abord, on trie ; ensuite, on recycle ; enfin, on réduit. (Connecteurs \zh{首先/然后/最后})
\end{enumerate}
\end{exercice}

\begin{corrige}[Thème -- Proposition de correction]
\begin{enumerate}
    \item \zh{我们应该把垃圾分类。} \emph{Wǒmen yīnggāi bǎ lājī fēnlèi.}
    \item \zh{旧报纸可以回收利用。} \emph{Jiù bàozhǐ kěyǐ huíshōu lìyòng.}
    \item \zh{请把塑料瓶放进蓝色的桶。} \emph{Qǐng bǎ sùliào píng fàng jìn lánsè de tǒng.}
    \item \zh{因为地球只有一个，所以我们应该保护它。} \emph{Yīnwèi dìqiú zhǐ yǒu yīgè, suǒyǐ wǒmen yīnggāi bǎohù tā.}
    \item \zh{首先分类，然后回收，最后减少。} \emph{Shǒuxiān fēnlèi, ránhòu huíshōu, zuìhòu jiǎnshǎo.}
\end{enumerate}
\end{corrige}

\courssub{Version : Chinois → Français (Entraînement guidé)}

\noindent Traduisez les phrases suivantes en français naturel.

\begin{exercice}[Version -- Phrases isolées]
\begin{enumerate}
    \item \zh{大家都应该少用塑料袋。}
    \item \zh{玻璃和金属可以回收利用。}
    \item \zh{请把湿垃圾放进棕色的桶。}
    \item \zh{因为环境污染严重，所以我们要保护环境。}
    \item \zh{首先洗干净瓶子，然后分类放好。}
\end{enumerate}
\end{exercice}

\begin{corrige}[Version -- Proposition de correction]
\begin{enumerate}
    \item Tout le monde devrait utiliser moins de sacs en plastique.
    \item Le verre et le métal peuvent être recyclés.
    \item Mettez les déchets humides dans la poubelle brune.
    \item Parce que la pollution de l'environnement est grave, nous devons protéger l'environnement.
    \item D'abord, lavez bien les bouteilles ; ensuite, triez-les et rangez-les bien.
\end{enumerate}
\end{corrige}

\courssub{Exercice résolu : Recomposition et traduction}

\begin{exoresolu}[Recomposition de texte et traduction]
\textbf{Énoncé :} Remettez les segments dans l'ordre logique pour former un texte cohérent, puis traduisez en français.
\begin{enumerate}
    \item \zh{最后，我们要少用塑料袋。}
    \item \zh{因为地球是我们的家。}
    \item \zh{首先，把垃圾分类。}
    \item \zh{我们应该保护环境。}
    \item \zh{然后，回收利用旧纸和瓶子。}
    \item \zh{所以保护环境很重要。}
\end{enumerate}

\tcblower
\textbf{Solution détaillée :}

\noindent \textbf{Ordre logique :} ④ $\rightarrow$ ③ $\rightarrow$ ⑤ $\rightarrow$ ① $\rightarrow$ ② $\rightarrow$ ⑥

\noindent \textbf{Texte reconstitué :}
\zh{我们应该保护环境。首先，把垃圾分类。然后，回收利用旧纸和瓶子。最后，我们要少用塑料袋。因为地球是我们的家，所以保护环境很重要。}

\noindent \textbf{Pinyin :}
\emph{Wǒmen yīnggāi bǎohù huánjìng. Shǒuxiān, bǎ lājī fēnlèi. Ránhòu, huíshōu lìyòng jiù zhǐ hé píngzi. Zuìhòu, wǒmen yào shǎo yòng sùliào dài. Yīnwèi dìqiú shì wǒmen de jiā, suǒyǐ bǎohù huánjìng hěn zhòngyào.}

\noindent \textbf{Traduction :}
Nous devons protéger l'environnement. D'abord, trions les déchets. Ensuite, recyclons le vieux papier et les bouteilles. Enfin, nous devons utiliser moins de sacs en plastique. Car la Terre est notre maison, donc protéger l'environnement est très important.

\noindent \textbf{Analyse des connecteurs :}
\begin{itemize}
    \item \zh{首先 / 然后 / 最后} : enchaînement chronologique des actions.
    \item \zh{因为..., 所以...} : relation cause-effet. Virgule obligatoire après la proposition \zh{因为}.
\end{itemize}
\end{exoresolu}

\courssub{Point de vigilance : Erreurs fréquentes des francophones}

\begin{attention}[Pièges à éviter -- Synthèse]
\begin{itemize}
    \item \textbf{\zh{把} + Verbe nu :} \zh{*我把垃圾分类} (incomplet si pas accompli) $\rightarrow$ \zh{我把垃圾分好类了} (résultat) / \zh{我把垃圾分类了} (accompli). \textbf{Règle :} Verbe + Complément (résultat, direction, \zh{了}, \zh{完}, \zh{好}...).
    \item \textbf{Objet indéfini avec \zh{把} :} \zh{*我把一个瓶子扔了} $\rightarrow$ \zh{我扔了一个瓶子} (pas de \zh{把}).
    \item \textbf{Confusion \zh{可以} / \zh{会} / \zh{能} :} \zh{*我可以中文} $\rightarrow$ \zh{我会中文} (compétence) / \zh{我可以说中文}. \zh{可以} = permission / possibilité objective.
    \item \textbf{Ordre \zh{因为...所以...} :} Jamais \zh{所以...因为...}. Cause TOUJOURS avant conséquence. Virgule après clause \zh{因为}.
    \item \textbf{Tons critiques :} \zh{分类} \emph{fēn\underline{lèi}} (1+4), pas *fēnlěi. \zh{利用} \emph{lì\underline{yòng}} (4+4). \zh{垃圾} \emph{lā\underline{jī}} (1+1). \zh{环境} \emph{huán\underline{jìng}} (2+4).
    \item \textbf{Caractères proches (pièges visuels) :}
    \begin{itemize}
        \item \zh{利} (lì, avantage, radical \zh{刂} couteau) $\neq$ \zh{丽} (lì, beau, radical \zh{一/鹿}).
        \item \zh{回} (huí, retourner, radical \zh{囗} enclos) $\neq$ \zh{囗} (wéi, enclos seul) / \zh{国} (guó, pays).
        \item \zh{塑} (sù, plastique, radical \zh{土}) $\neq$ \zh{俗} (sú, coutume, radical \zh{人}).
        \item \zh{环} (huán, anneau/env., radical \zh{王}) $\neq$ \zh{还} (huán/hái, radical \zh{辶}).
    \end{itemize}
\end{itemize}
\end{attention}

\courssub{Le saviez-vous ? : Le tri obligatoire à Shanghai}

\begin{savaistu}[Shanghai 2019 : Le modèle chinois du tri obligatoire]
Depuis le 1\up{er} juillet 2019, Shanghai applique le règlement le plus strict de Chine (\zh{上海市生活垃圾管理条例}). Quatre catégories obligatoires, code couleur universel :
\begin{itemize}
    \item \zh{可回收物} (kě huíshōu wù) — \textbf{Bleu} : Papier, plastique, verre, métal, textile.
    \item \zh{有害垃圾} (yǒuhài lājī) — \textbf{Rouge} : Piles, ampoules, médicaments, peinture.
    \item \zh{湿垃圾} (shī lājī) — \textbf{Brun} : Épluchures, restes repas, fleurs (compostable).
    \item \zh{干垃圾} (gān lājī) — \textbf{Gris} : Mouchoirs, cendres, sacs souillés (incinération/enfouissement).
\end{itemize}
Amendes : Particuliers 50--200 \zh{元} (yuán) ; Entreprises jusqu'à 50 000 \zh{元}. Des « superviseurs de tri » (\zh{分类督导员}) veillent aux points de dépôt à heures fixes. Pékin, Shenzhen, Chengdu ont suivi. Au Cameroun, des initiatives locales existent (coopératives Mokolo, Ngousso, ramassage sélectif à l'université), mais pas encore de cadre national contraignant comparable. La loi-cadre sur l'environnement (1996) et l'arrêté sur les plastiques non biodégradables (2012/2014) posent les bases, l'application reste le défi.
\end{savaistu}

\courssub{Production guidée : Votre paragraphe sur le recyclage}

\begin{exercice}[Expression écrite -- Production finale (20 min)]
Écrivez en chinois (caractères + pinyin) un court texte de \textbf{5 à 6 phrases} sur le recyclage chez vous / à l'école / dans votre quartier.
\textbf{Contraintes obligatoires (grille d'évaluation ci-dessous) :}
\begin{itemize}
    \item [$\square$] 1 phrase d'introduction avec \zh{应该 + V}.
    \item [$\square$] 3 actions séquentielles avec \zh{首先}，\zh{然后}，\zh{最后}.
    \item [$\square$] Au moins 1 phrase avec \zh{可以 + V} (possibilité/permission).
    \item [$\square$] Au moins 1 phrase avec \zh{把 + O + V + Complément} (objet spécifique).
    \item [$\square$] Conclusion avec \zh{因为..., 所以...} (virgule après \zh{因为}).
    \item [$\square$] Minimum \textbf{8 mots} du vocabulaire du tableau (Section 3).
    \item [$\square$] Caractères corrects, pinyin avec tons, ponctuation chinoise.
\end{itemize}
\end{exercice}

\courssub{Grille d'évaluation de la production écrite (Sur 20 pts)}

\begin{center}
\begin{tabularx}{\linewidth}{|X|c|X|X|}
\hline
\rowcolor{popblueL} \textbf{Critère} & \textbf{Points} & \textbf{Indicateurs de réussite} \\ \hline
Vocabulaire (8+ mots thème, exacts) & 4 & Mots du tableau utilisés à bon escient, orthographe caractères correcte. \\ \hline
Structure \zh{应该} & 2 & \zh{应该/不应该} + V bien formé, sens obligation/conseil respecté. \\ \hline
Structure \zh{可以} & 2 & \zh{可以/不可以} + V bien formé, distinction \zh{会/能} respectée. \\ \hline
Structure \zh{把} & 3 & Objet défini + \zh{把} + V + Complément (résultat/direction/le/guo). Pas de verbe nu. \\ \hline
Connecteurs \zh{首先/然后/最后} & 3 & Trois étapes claires, ponctuation correcte (virgule après connecteur). \\ \hline
Structure \zh{因为...所以...} & 2 & Cause $\rightarrow$ Conséquence, virgule après clause \zh{因为}, sens logique. \\ \hline
Pinyin & Tons & 2 & Tons justes sur tout le texte, espaces entre mots, majuscules noms propres. \\ \hline
Caractères & Ponctuation & 2 & Traits corrects, radicaux visibles, ponctuation pleine chasse (\zh{，。}). \\ \hline
\rowcolor{popblueL} \textbf{TOTAL} & \textbf{20} &  \\ \hline
\end{tabularx}
\end{center}

\begin{corrige}[Exercice Production -- Proposition de correction type]
\textbf{Exemple de production attendue (niveau excellence) :}
\zh{我们应该保护环境。首先，我们要把垃圾分好类：塑料、纸、玻璃分开放。然后，旧报纸和塑料瓶可以回收利用。最后，少用塑料袋，多用布袋子。因为环境保护关系到我们的健康，所以人人都要行动。}

\noindent \textbf{Pinyin :}
\emph{Wǒmen yīnggāi bǎohù huánjìng. Shǒuxiān, wǒmen yào bǎ lājī fēn hǎo lèi: sùliào, zhǐ, bōlí fēnkāi fàng. Ránhòu, jiù bàozhǐ hé sùliào píng kěyǐ huíshōu lìyòng. Zuìhòu, shǎo yòng sùliào dài, duō yòng bùdàizi. Yīnwèi huánjìng bǎohù guāndào wǒmen de jiànkāng, suǒyǐ rénrén dōu yào xíngdòng.}

\noindent \textbf{Traduction :}
Nous devons protéger l'environnement. D'abord, nous devons bien trier les déchets : mettre le plastique, le papier, le verre séparément. Ensuite, les vieux journaux et les bouteilles en plastique peuvent être recyclés. Enfin, utilisons moins de sacs en plastique, plus de sacs en tissu. Car la protection de l'environnement concerne notre santé, donc tout le monde doit agir.

\noindent \textbf{Vérification grille :}
\begin{itemize}
    \item Vocabulaire : 环境, 垃圾, 分类, 塑料, 纸, 玻璃, 报纸, 瓶子, 回收利用, 塑料袋, 布袋 (11 mots) ✓
    \item \zh{应该保护} ✓ | \zh{可以回收利用} ✓ | \zh{把垃圾分好类} (把+O+V+Résultat) ✓
    \item \zh{首先 / 然后 / 最后} ✓ | \zh{因为..., 所以...} (virgule) ✓
    \item Tons : \emph{fēn hǎo lèi, huíshōu lìyòng, guāndào, jiànkāng} ✓
\end{itemize}
\end{corrige}

\courssub{Devoirs / Approfondissement personnel}
\begin{enumerate}
    \item Réviser l'écriture des 16 mots du tableau vocabulaire (dictée caractères + pinyin + sens).
    \item Refaire les exercices Thème/Version (Sections 8 \& 9) sans regarder la correction.
    \item Préparer un exposé oral de 1 minute : « Le tri des déchets dans mon quartier / mon lycée » en réutilisant les 3 structures grammaticales et les connecteurs. Enregistrez-vous sur OnBuch+.
    \item Recherche : Comparez les couleurs des poubelles au Cameroun (si existantes) et en Chine. Notez 3 différences / similitudes en chinois.
\end{enumerate}

\coursec{Structure et connecteurs du texte}

Dans la section précédente, nous avons lu et commenté un texte modèle sur le recyclage. Nous avons remarqué que ce texte n'est pas une simple suite de phrases : les idées sont \cle{enchaînées}, les actions sont \cle{ordonnées}, et certaines structures reviennent régulièrement, comme \zh{把} ou \zh{因为……所以……}. Dans cette deuxième section, nous allons ouvrir la « boîte à outils » du rédacteur chinois : d'abord les connecteurs qui organisent le texte, ensuite les verbes modaux qui expriment le devoir et la possibilité, et enfin la fameuse construction en \zh{把}, indispensable pour parler du \emph{traitement} des déchets.

\courssub{Observer avant de comprendre : un mini-texte repère}

Lisons d'abord ce court paragraphe. Soulignez mentalement les mots qui relient les phrases entre elles.

\begin{textebox}[Mini-texte repère : 回收垃圾，从我做起]
\zh{首先，我们应该把垃圾分类。然后，可以把纸、塑料和玻璃放进不同的箱子。最后，把旧瓶子送到回收站。因为地球只有一个，所以我们要保护环境。如果每个人都做一点小事，我们的城市就会更干净。}\\
\emph{Shǒuxiān, wǒmen yīnggāi bǎ lājī fēnlèi. Ránhòu, kěyǐ bǎ zhǐ, sùliào hé bōli fàng jìn bùtóng de xiāngzi. Zuìhòu, bǎ jiù píngzi sòng dào huíshōu zhàn. Yīnwèi dìqiú zhǐyǒu yí ge, suǒyǐ wǒmen yào bǎohù huánjìng. Rúguǒ měi ge rén dōu zuò yìdiǎn xiǎoshì, wǒmen de chéngshì jiù huì gèng gānjìng.}\\
« D'abord, nous devons trier les déchets. Ensuite, on peut mettre le papier, le plastique et le verre dans des boîtes différentes. Enfin, on apporte les vieilles bouteilles au centre de recyclage. Parce qu'il n'y a qu'une seule Terre, nous devons protéger l'environnement. Si chacun fait un petit geste, notre ville sera plus propre. »
\end{textebox}

On observe trois familles de mots-outils :
\begin{itemize}
\item des \cle{connecteurs de séquence} qui ordonnent les actions : \zh{首先} (d'abord), \zh{然后} (ensuite), \zh{最后} (enfin) ;
\item des \cle{connecteurs logiques} qui relient les idées : \zh{因为……所以……} (parce que… donc…), \zh{如果……就……} (si… alors…), \zh{但是} (mais) ;
\item des \cle{verbes modaux} qui expriment le devoir ou la possibilité : \zh{应该} (devoir), \zh{要} (devoir, vouloir), \zh{可以} (pouvoir, permission), \zh{不能} (ne pas pouvoir, interdiction).
\end{itemize}

\courssub{Les connecteurs de séquence : ordonner les actions}

\begin{propriete}[Connecteurs de séquence]
Pour décrire une \emph{procédure} (comment trier, comment recycler), le chinois utilise trois connecteurs placés \textbf{en tête de phrase}, suivis d'une virgule chinoise \zh{，} :
\begin{itemize}
\item \zh{首先} \emph{shǒuxiān} = d'abord, premièrement ;
\item \zh{然后} \emph{ránhòu} = ensuite, puis ;
\item \zh{最后} \emph{zuìhòu} = enfin, pour finir.
\end{itemize}
Schéma : \textbf{首先，+ phrase 1。然后，+ phrase 2。最后，+ phrase 3。}
\end{propriete}

\begin{exemplebox}[Une procédure complète]
\zh{首先，我们把垃圾分类。然后，把瓶子洗干净。最后，把它们送到回收站。}\\
\emph{Shǒuxiān, wǒmen bǎ lājī fēnlèi. Ránhòu, bǎ píngzi xǐ gānjìng. Zuìhòu, bǎ tāmen sòng dào huíshōu zhàn.}\\
« D'abord, nous trions les déchets. Ensuite, nous lavons les bouteilles. Enfin, nous les apportons au centre de recyclage. »
\end{exemplebox}

\begin{attention}[然后 ≠ « puis » partout]
\zh{然后} s'emploie surtout pour enchaîner des \emph{actions successives}. Pour ajouter une simple information (« et aussi »), on préfère \zh{还有} ou \zh{而且}. Ne mettez pas \zh{然后} devant chaque phrase : un texte avec cinq \zh{然后} est aussi lourd en chinois qu'en français !
\end{attention}

\courssub{Les connecteurs logiques : relier les idées}

\begin{propriete}[因为……所以…… : la cause et la conséquence]
Structure : \textbf{因为 + cause，所以 + conséquence。}\\
En chinois, il est \textbf{correct et même recommandé} d'utiliser les deux mots ensemble dans la même phrase. C'est une différence majeure avec le français, où « parce que… donc… » est une faute.
\begin{itemize}
\item \zh{因为} \emph{yīnwèi} = parce que (introduit la cause) ;
\item \zh{所以} \emph{suǒyǐ} = donc, c'est pourquoi (introduit la conséquence).
\end{itemize}
On peut utiliser \zh{因为} seul ou \zh{所以} seul, mais jamais les deux ne se contredisent : ils se complètent.
\end{propriete}

\begin{exemplebox}[因为……所以…… en action]
\zh{因为我们只有一个地球，所以我们要保护它。}\\
\emph{Yīnwèi wǒmen zhǐyǒu yí ge dìqiú, suǒyǐ wǒmen yào bǎohù tā.}\\
« Parce que nous n'avons qu'une seule Terre, nous devons la protéger. »\\[4pt]
\zh{因为塑料污染河流，所以很多人开始回收塑料。}\\
\emph{Yīnwèi sùliào wūrǎn héliú, suǒyǐ hěn duō rén kāishǐ huíshōu sùliào.}\\
« Parce que le plastique pollue les rivières, beaucoup de gens commencent à recycler le plastique. »
\end{exemplebox}

\begin{attention}[Le piège n°1 du francophone]
En français, on écrit « Parce qu'il pleut, \emph{donc} je reste » : c'est fautif. En chinois, c'est l'inverse : \zh{因为下雨，所以我待在家里。} est parfaitement correct. \textbf{Ne traduisez jamais mot à mot} : quand vous voyez \zh{因为……所以……}, traduisez par « parce que…, … » ou « …, donc… », mais pas les deux en français.
\end{attention}

\begin{propriete}[但是 et 如果……就……]
\begin{itemize}
\item \zh{但是} \emph{dànshì} = mais. Structure : \textbf{phrase 1，但是 + phrase 2。} Il marque l'opposition, comme en français.
\item \zh{如果……就……} \emph{rúguǒ… jiù…} = si… alors… Structure : \textbf{如果 + condition，(sujet) + 就 + résultat。} Le mot \zh{就} se place \emph{avant le verbe} de la seconde proposition et ne se traduit pas séparément.
\end{itemize}
\end{propriete}

\begin{exemplebox}[Opposition et condition]
\zh{回收垃圾有点麻烦，但是它对环境很好。}\\
\emph{Huíshōu lājī yǒudiǎn máfan, dànshì tā duì huánjìng hěn hǎo.}\\
« Recycler les déchets est un peu compliqué, mais c'est très bon pour l'environnement. »\\[4pt]
\zh{如果每个人都回收，我们的城市就会更美丽。}\\
\emph{Rúguǒ měi ge rén dōu huíshōu, wǒmen de chéngshì jiù huì gèng měilì.}\\
« Si tout le monde recycle, notre ville sera plus belle. »\\[4pt]
\zh{如果你不把垃圾分类，回收站就不能处理它们。}\\
\emph{Rúguǒ nǐ bù bǎ lājī fēnlèi, huíshōu zhàn jiù bù néng chǔlǐ tāmen.}\\
« Si tu ne tries pas les déchets, le centre de recyclage ne pourra pas les traiter. »
\end{exemplebox}

\begin{center}
\begin{tabularx}{\linewidth}{|X|X|X|}
\hline
\rowcolor{popblueL} Fonction & Connecteur & Exemple (caractères + pinyin) \\
\hline
Début d'une procédure & \zh{首先} shǒuxiān & \zh{首先，我们应该把垃圾分类。} Shǒuxiān, wǒmen yīnggāi bǎ lājī fēnlèi. \\
\hline
Étape suivante & \zh{然后} ránhòu & \zh{然后，把纸放进蓝箱子。} Ránhòu, bǎ zhǐ fàng jìn lán xiāngzi. \\
\hline
Dernière étape & \zh{最后} zuìhòu & \zh{最后，把瓶子送到回收站。} Zuìhòu, bǎ píngzi sòng dào huíshōu zhàn. \\
\hline
Cause + conséquence & \zh{因为……所以……} yīnwèi… suǒyǐ… & \zh{因为地球只有一个，所以我们要保护它。} Yīnwèi dìqiú zhǐyǒu yí ge, suǒyǐ wǒmen yào bǎohù tā. \\
\hline
Opposition & \zh{但是} dànshì & \zh{回收有点难，但是很重要。} Huíshōu yǒudiǎn nán, dànshì hěn zhòngyào. \\
\hline
Condition & \zh{如果……就……} rúguǒ… jiù… & \zh{如果大家一起努力，环境就会变好。} Rúguǒ dàjiā yìqǐ nǔlì, huánjìng jiù huì biàn hǎo. \\
\hline
\end{tabularx}
\end{center}

\courssub{Les verbes modaux : devoir et possibilité}

Pour écrire un texte qui \emph{conseille} ou \emph{encourage} (genre typique du texte sur le recyclage), le chinois utilise des verbes modaux placés \textbf{directement avant le verbe principal} (ou avant \zh{把} si la structure \zh{把} est utilisée).

\begin{definition}[Les quatre modaux du texte argumentatif]
\begin{itemize}
\item \zh{应该} \emph{yīnggāi} = devoir (conseil moral, ce qui est souhaitable) : « on devrait » ;
\item \zh{要} \emph{yào} = devoir, il faut (nécessité plus forte, ou intention) ;
\item \zh{可以} \emph{kěyǐ} = pouvoir, il est permis de (possibilité, permission) ;
\item \zh{不能} \emph{bù néng} = ne pas pouvoir, il est interdit de (interdiction, impossibilité).
\end{itemize}
Schéma : \textbf{Sujet + modal + (把 + objet +) verbe (+ complément)。}
\end{definition}

\begin{exemplebox}[Nuances de sens]
\zh{我们应该把垃圾分类。} \emph{Wǒmen yīnggāi bǎ lājī fēnlèi.} « Nous devrions trier les déchets. » (conseil)\\[3pt]
\zh{我们要保护环境。} \emph{Wǒmen yào bǎohù huánjìng.} « Nous devons protéger l'environnement. » (nécessité, engagement)\\[3pt]
\zh{你可以把旧书送给别人。} \emph{Nǐ kěyǐ bǎ jiù shū sòng gěi biérén.} « Tu peux donner tes vieux livres à d'autres. » (possibilité)\\[3pt]
\zh{我们不能把垃圾扔在河里。} \emph{Wǒmen bù néng bǎ lājī rēng zài hé li.} « Nous ne pouvons pas jeter les déchets dans la rivière. » (interdiction)
\end{exemplebox}

\begin{attention}[La négation des modaux]
La négation se place \textbf{avant le modal}, pas avant le verbe principal : \zh{不应该} (ne devrait pas), \zh{不要} (ne pas devoir, défense de), \zh{不能} (ne pas pouvoir). Ne dites jamais \zh{应该不} pour « ne devrait pas » : la forme correcte est \zh{不应该}.
\end{attention}

\courssub{La construction 把 : traiter un objet}

Pour parler du recyclage, on doit constamment dire ce qu'on \emph{fait} des objets : trier les déchets, laver les bouteilles, jeter le papier. Le français utilise souvent la voix passive (« les déchets \emph{sont triés} ») ; le chinois préfère une construction active spéciale avec \zh{把}.

\begin{propriete}[La règle de 把]
Structure : \textbf{Sujet + 把 + Objet + Verbe + complément (ou 了)。}\\[4pt]
Deux conditions obligatoires :
\begin{enumerate}
\item \textbf{L'objet doit être défini ou connu.} On déplace avant le verbe un objet précis, identifié par le contexte (\zh{这个}, \zh{这些}, possessif, ou objet déjà mentionné). On ne dit pas \zh{把一个瓶子} pour « une bouteille quelconque » sans contexte.
\item \textbf{Le verbe ne peut pas rester seul.} Il doit être suivi d'un complément (résultat, direction, lieu) ou de \zh{了}. On dit \zh{把垃圾分类} (le verbe \zh{分类} est déjà un résultat composé), \zh{把瓶子洗干净} (verbe + résultat \zh{干净}), \zh{把书放在桌子上} (verbe + lieu). Jamais \zh{把垃圾扔} tout seul.
\end{enumerate}
Sens : le sujet \emph{fait subir une action} à l'objet ; l'objet est « manipulé, traité, déplacé, transformé ».
\end{propriete}

\begin{popfigure}
\begin{tikzpicture}[
box/.style={draw=popblue, fill=popblueL, rounded corners=3pt, align=center, minimum height=0.9cm, inner sep=5pt},
arr/.style={-{Stealth[length=3mm]}, thick, popdark}]
\node[box, align=center] (s) at (0,0){\textbf{Sujet}\\ \zh{我们} wǒmen};
\node[box, fill=poporangeL, draw=poporange, align=center] (ba) at (3.2,0){\zh{把}\\ bǎ};
\node[box, fill=poppurpleL, draw=poppurple, align=center] (o) at (5.6,0){\textbf{Objet}\\ \zh{垃圾} lājī};
\node[box, fill=popgreenL, draw=popgreen, align=center] (v) at (9.2,0){\textbf{Verbe + complément}\\ \zh{分类} fēnlèi};
\draw[arr] (s) -- (ba);
\draw[arr] (ba) -- (o);
\draw[arr] (o) -- (v);
\node[align=center, popmuted] at (4.6,-1.4) {我们 + 把 + 垃圾 + 分类。= Nous trions les déchets.};
\end{tikzpicture}
\legende{Schéma de la construction 把 : l'objet passe avant le verbe.}
\end{popfigure}

\begin{exemplebox}[把 dans le thème du recyclage]
\zh{请把垃圾分类。} \emph{Qǐng bǎ lājī fēnlèi.} « Trie les déchets, s'il te plaît. »\\[3pt]
\zh{妈妈把旧报纸卖给了回收的人。} \emph{Māma bǎ jiù bàozhǐ mài gěi le huíshōu de rén.} « Maman a vendu les vieux journaux au récupérateur. »\\[3pt]
\zh{请把瓶子放进绿色的箱子。} \emph{Qǐng bǎ píngzi fàng jìn lǜsè de xiāngzi.} « Mets les bouteilles dans la caisse verte, s'il te plaît. »\\[3pt]
\zh{他把水喝完了。} \emph{Tā bǎ shuǐ hē wán le.} « Il a fini de boire l'eau (la bouteille est vide, on peut la recycler !). »
\end{exemplebox}

\begin{savaistu}[把 ou 被 ?]
Le français dit : « Les déchets \emph{sont triés} par les habitants. » Le chinois a deux stratégies : la construction active \zh{把} : \zh{居民把垃圾分类。} (les habitants trient les déchets), ou la construction passive \zh{被} : \zh{垃圾被居民分类。} (les déchets sont triés par les habitants). Dans un texte de conseils comme le nôtre, on préfère \zh{把}, plus dynamique, car on s'adresse au lecteur : \zh{请你把垃圾分类！}
\end{savaistu}

\courssub{Méthode : construire un paragraphe de conseils}

\begin{methode}[La recette du paragraphe sur le recyclage]
Pour écrire un paragraphe simple et bien structuré, suivez ces cinq étapes :
\begin{enumerate}
\item \textbf{Ouvrir avec un connecteur de séquence} : \zh{首先，} / \zh{然后，} / \zh{最后，}
\item \textbf{Ajouter un modal} selon l'intention : \zh{我们应该…} (conseil), \zh{我们可以…} (possibilité), \zh{我们不能…} (interdiction).
\item \textbf{Utiliser 把 + verbe d'action} pour parler du traitement : \zh{把垃圾分类}, \zh{把瓶子送到回收站}.
\item \textbf{Justifier avec 因为……所以……} ou proposer une condition avec \zh{如果……就……}.
\item \textbf{Conclure avec la formule type} : \zh{保护环境就是保护我们自己。} \emph{Bǎohù huánjìng jiùshì bǎohù wǒmen zìjǐ.} « Protéger l'environnement, c'est nous protéger nous-mêmes. »
\end{enumerate}
Modèle complet : \zh{首先，我们应该把垃圾分类。然后，我们可以把纸和塑料送到回收站。如果我们每个人都这样做，我们的城市就会更干净。保护环境就是保护我们自己！}
\end{methode}

\begin{exoresolu}[Assembler un paragraphe]
Énoncé : remettez les phrases suivantes dans l'ordre logique et ajoutez les connecteurs manquants (\zh{首先}, \zh{然后}, \zh{最后}, \zh{所以}) pour obtenir un paragraphe cohérent.\\
a) \zh{把旧瓶子送到回收站。} b) \zh{我们应该把垃圾分类。} c) \zh{我们的城市会更干净。} d) \zh{把纸和塑料放进不同的箱子。}
\tcblower
Solution détaillée : la procédure va du tri (b) au rangement (d) puis au transport (a) ; la phrase (c) est une conséquence, introduite par \zh{所以}.\\
\zh{首先，我们应该把垃圾分类。然后，把纸和塑料放进不同的箱子。最后，把旧瓶子送到回收站，所以我们的城市会更干净。}\\
\emph{Shǒuxiān, wǒmen yīnggāi bǎ lājī fēnlèi. Ránhòu, bǎ zhǐ hé sùliào fàng jìn bùtóng de xiāngzi. Zuìhòu, bǎ jiù píngzi sòng dào huíshōu zhàn, suǒyǐ wǒmen de chéngshì huì gèng gānjìng.}\\
« D'abord, nous devrions trier les déchets. Ensuite, mettre le papier et le plastique dans des caisses différentes. Enfin, apporter les vieilles bouteilles au centre de recyclage ; ainsi notre ville sera plus propre. »
\end{exoresolu}

\begin{exercice}[À vous de jouer]
Construisez trois phrases complètes avec les éléments donnés, en respectant la structure \zh{把} et en ajoutant un modal :
\begin{enumerate}
\item nous / devoir / déchets / trier (\zh{应该}, \zh{把}, \zh{垃圾}, \zh{分类})
\item tu / pouvoir / vieux journaux / vendre (\zh{可以}, \zh{把}, \zh{旧报纸}, \zh{卖})
\item on / ne pas pouvoir / déchets / jeter dans la rivière (\zh{不能}, \zh{把}, \zh{垃圾}, \zh{扔在河里})
\end{enumerate}
\end{exercice}

\begin{corrige}[Exercice « À vous de jouer »]
\begin{enumerate}
\item \zh{我们应该把垃圾分类。} \emph{Wǒmen yīnggāi bǎ lājī fēnlèi.} « Nous devrions trier les déchets. »
\item \zh{你可以把旧报纸卖掉。} \emph{Nǐ kěyǐ bǎ jiù bàozhǐ mài diào.} « Tu peux vendre les vieux journaux. » (le complément \zh{掉} marque l'achèvement : le verbe ne reste pas seul.)
\item \zh{我们不能把垃圾扔在河里。} \emph{Wǒmen bù néng bǎ lājī rēng zài hé li.} « Nous ne pouvons pas jeter les déchets dans la rivière. »
\end{enumerate}
Points de vigilance : le modal se place \emph{avant} \zh{把} ; l'objet après \zh{把} est défini ; le verbe est toujours accompagné d'un complément ou d'un élément de résultat.
\end{corrige}

\begin{aretenir}[L'essentiel de la section]
\begin{itemize}
\item Connecteurs de séquence : \zh{首先} → \zh{然后} → \zh{最后}, en tête de phrase, suivis de \zh{，}.
\item Connecteurs logiques : \zh{因为……所以……} (les deux ensemble, c'est correct !), \zh{但是}, \zh{如果……就……} (\zh{就} avant le verbe du résultat).
\item Modaux : \zh{应该} (conseil), \zh{要} (nécessité), \zh{可以} (possibilité), \zh{不能} (interdiction) ; négation avant le modal.
\item Structure \zh{把} : Sujet + \zh{把} + objet défini + verbe + complément ; le verbe ne reste jamais seul.
\item Conclusion type : \zh{保护环境就是保护我们自己。}
\end{itemize}
\end{aretenir}

\coursec{Écriture correcte des caractères du thème}

Dans la section précédente, nous avons vu comment organiser un texte sur le recyclage grâce aux connecteurs \zh{首先}、\zh{然后}、\zh{最后} et aux structures simples. Mais un texte bien structuré perd toute sa valeur si les caractères sont mal écrits : un trait de travers, un radical confondu, et le sens change complètement. Cette section est donc consacrée à l'\cle{écriture correcte} des dix caractères clés de notre thème : observer leurs composants, respecter l'ordre des traits, et ne pas les confondre avec leurs « jumeaux » dangereux.

\courssub{Observation : les dix caractères du thème}

Avant d'écrire, observons. Voici les dix caractères que vous devez maîtriser pour rédiger votre texte sur le recyclage, avec leur nature grammaticale et leur composition.

\begin{center}
\begin{tabularx}{\linewidth}{|X|X|X|X|}
\hline
\rowcolor{popblueL} Caractères & Pinyin & Nature & Traduction \\
\hline
垃 & lā & Radical 土 + Phonétique 立 (Morphème lié) & déchet (dans 垃圾) \\
\hline
圾 & jī & Radical 土 + Phonétique 及 (Morphème lié) & déchet (dans 垃圾) \\
\hline
回 & huí & Radical 囗 + Phonétique 口 (Verbe) & retourner, revenir \\
\hline
收 & shōu & Composant 丩 + Radical 攵 (Verbe) & recevoir, ramasser \\
\hline
利 & lì & Radical 禾 + Phonétique 刂 (Adj./Nom) & profit, avantage \\
\hline
用 & yòng & Caractère unique (Verbe) & utiliser \\
\hline
环 & huán & Radical 王 + Phonétique 不 (Nom) & anneau ; entourage \\
\hline
境 & jìng & Radical 土 + Phonétique 竟 (Nom) & territoire, milieu \\
\hline
塑 & sù & Phonétique 朔 + Radical 土 (Verbe) & modeler (plastique) \\
\hline
料 & liào & Radical 米 + Phonétique 斗 (Nom) & matériau, matière \\
\hline
\end{tabularx}
\end{center}

Observez attentivement : beaucoup de ces caractères se composent d'un \cle{radical} (qui donne une indication de sens) et d'une \cle{partie phonétique} (qui donne une indication de prononciation). C'est la clé de la mémorisation.

\begin{definition}[回收 huíshōu]
\zh{回收} (huíshōu) est un verbe qui signifie « recycler, récupérer ». Il est composé de \zh{回} (huí, « retourner, revenir ») et de \zh{收} (shōu, « recevoir, ramasser ») : littéralement, « faire revenir ce qu'on a ramassé », c'est-à-dire récupérer pour réutiliser. Exemple : \zh{我们回收塑料瓶。} (Wǒmen huíshōu sùliào píng.) « Nous recyclons les bouteilles en plastique. »
\end{definition}

\begin{exemplebox}[Trois expressions à écrire sans faute]
\zh{垃圾分类} (lājī fēnlèi) : « le tri des déchets ».\\
\zh{回收利用} (huíshōu lìyòng) : « recycler et réutiliser ».\\
\zh{塑料袋} (sùliào dài) : « le sac en plastique ».
\end{exemplebox}

\courssub{Les radicaux : la clé du sens}

\begin{propriete}[Les radicaux de notre thème]
Un \cle{radical} (ou clé) est le composant d'un caractère qui indique sa famille de sens.
\begin{itemize}
\item \zh{土} (tǔ, « terre ») : on le retrouve dans \zh{垃}, \zh{圾}, \zh{境} et \zh{塑}. Logique : les déchets et l'environnement sont liés à la terre.
\item \zh{王} (wáng, « jade ») : dans \zh{环}. À l'origine, 环 désignait un anneau de jade, d'où l'idée de « ce qui entoure » → environnement.
\item \zh{囗} (wéi, « enceinte ») : dans \zh{回}. Attention, ce n'est pas 口 (kǒu, « bouche ») : c'est un grand cadre qui enferme.
\item \zh{禾} (hé, « céréale ») : dans \zh{利}. La céréale coupée avec le couteau 刂 évoque le profit de la récolte.
\end{itemize}
\end{propriete}

\begin{methode}[Comment mémoriser un caractère en trois étapes]
\begin{enumerate}
\item \textbf{Identifier le radical} : il donne le sens général (ex. 土 = terre dans 境).
\item \textbf{Identifier la partie phonétique} : elle donne une piste pour le son (ex. 竟 jìng dans 境 jìng ; 及 jí dans 圾 jī).
\item \textbf{Écrire 5 fois en comptant les traits} à voix haute, en respectant l'ordre des traits. Compter les traits oblige le cerveau à enregistrer la structure.
\end{enumerate}
\end{methode}

\courssub{L'ordre des traits : les quatre règles d'or}

\begin{propriete}[Règles générales d'écriture]
\begin{enumerate}
\item De \textbf{haut en bas} : 立 dans 垃 s'écrit du point supérieur vers la base.
\item De \textbf{gauche à droite} : dans 收, on écrit d'abord la partie gauche 丩, puis la partie droite (variante de 攵).
\item L'\textbf{extérieur avant l'intérieur} : dans 回, on trace d'abord le cadre extérieur (sans le fermer), puis le petit 口 intérieur.
\item La \textbf{fermeture en dernier} : le trait horizontal du bas qui ferme une enceinte (囗) se trace toujours en dernier.
\end{enumerate}
\end{propriete}

Le caractère \zh{回} (6 traits) illustre parfaitement les règles 3 et 4 : on ne ferme jamais l'enceinte avant d'avoir écrit ce qu'elle contient.

\begin{popfigure}
\begin{center}
\begin{tikzpicture}[scale=0.85, every node/.style={font=\small}]
% Stroke 1
\begin{scope}
\draw[popline, dashed] (0,0) rectangle (2,2);
\draw[popblue, very thick] (0.4,1.7) -- (0.4,0.3);
\node[popblue, align=center] at (1,-0.5) {1 : vertical\\gauche};
\end{scope}
% Stroke 2
\begin{scope}[xshift=2.6cm]
\draw[popline, dashed] (0,0) rectangle (2,2);
\draw[popblue, very thick] (0.4,1.7) -- (1.6,1.7) -- (1.6,0.3);
\draw[popblue!40, thin] (0.4,1.7) -- (0.4,0.3);
\node[popblue, align=center] at (1,-0.5) {2 : haut +\\côté droit};
\end{scope}
% Stroke 3
\begin{scope}[xshift=5.2cm]
\draw[popline, dashed] (0,0) rectangle (2,2);
\draw[popblue!40, thin] (0.4,1.7) -- (0.4,0.3);
\draw[popblue!40, thin] (0.4,1.7) -- (1.6,1.7) -- (1.6,0.3);
\draw[poporange, very thick] (0.75,1.25) -- (0.75,0.75);
\node[poporange, align=center] at (1,-0.5) {3 : vertical du\\petit 口};
\end{scope}
% Stroke 4
\begin{scope}[xshift=7.8cm]
\draw[popline, dashed] (0,0) rectangle (2,2);
\draw[popblue!40, thin] (0.4,1.7) -- (0.4,0.3);
\draw[popblue!40, thin] (0.4,1.7) -- (1.6,1.7) -- (1.6,0.3);
\draw[poporange!40, thin] (0.75,1.25) -- (0.75,0.75);
\draw[poporange, very thick] (0.75,1.25) -- (1.25,1.25) -- (1.25,0.75);
\node[poporange, align=center] at (1,-0.5) {4 : haut + droit\\du petit 口};
\end{scope}
% Stroke 5
\begin{scope}[xshift=10.4cm]
\draw[popline, dashed] (0,0) rectangle (2,2);
\draw[popblue!40, thin] (0.4,1.7) -- (0.4,0.3);
\draw[popblue!40, thin] (0.4,1.7) -- (1.6,1.7) -- (1.6,0.3);
\draw[poporange!40, thin] (0.75,1.25) -- (0.75,0.75);
\draw[poporange!40, thin] (0.75,1.25) -- (1.25,1.25) -- (1.25,0.75);
\draw[popgreen, very thick] (0.75,0.75) -- (1.25,0.75);
\node[popgreen, align=center] at (1,-0.5) {5 : on ferme\\le petit 口};
\end{scope}
% Stroke 6
\begin{scope}[xshift=13.0cm]
\draw[popline, dashed] (0,0) rectangle (2,2);
\draw[popblue!40, thin] (0.4,1.7) -- (0.4,0.3);
\draw[popblue!40, thin] (0.4,1.7) -- (1.6,1.7) -- (1.6,0.3);
\draw[poporange!40, thin] (0.75,1.25) -- (0.75,0.75);
\draw[poporange!40, thin] (0.75,1.25) -- (1.25,1.25) -- (1.25,0.75);
\draw[poporange!40, thin] (0.75,0.75) -- (1.25,0.75);
\draw[poppurple, very thick] (0.4,0.3) -- (1.6,0.3);
\node[poppurple, align=center] at (1,-0.5) {6 : on ferme\\l'enceinte};
\end{scope}
\end{tikzpicture}
\end{center}
\legende{Ordre des 6 traits de 回 : l'extérieur d'abord, l'intérieur ensuite, la fermeture en tout dernier.}
\end{popfigure}

Pour \zh{收} (6 traits), la règle dominante est « de gauche à droite » : on écrit d'abord la partie gauche 丩 (3 traits : vertical, puis les deux traits horizontaux), puis la partie droite (variante simplifiée du radical 攵, 3 traits : point, horizontal, trait final croisé \zh{乀}).

\courssub{Caractères proches : les pièges à éviter}

\begin{attention}[环 n'est pas 坏 !]
\zh{环} (huán, « environnement ») s'écrit avec le radical \zh{王} (jade) à gauche : 王 + 不. Si vous écrivez \zh{土} (terre) à la place, vous obtenez \zh{坏} (huài, « mauvais, cassé ») ! Écrire 保护环境 (bǎohù huánjìng, « protéger l'environnement ») avec 坏 donnerait « protéger le mauvais territoire » : le sens change complètement. De même, ne confondez pas 圾 (jī, déchet, radical 土) avec 级 (jí, niveau, radical 纟 soie), ni 收 (shōu, recevoir) avec 放 (fàng, poser, libérer).
\end{attention}

\begin{center}
\begin{tabularx}{\linewidth}{|X|X|X|}
\hline
\rowcolor{popblueL} Paire & Pinyin et composition & Sens \\
\hline
圾 / 级 & jī : 土 + 及 / jí : 纟 + 及 & déchet (垃圾) / niveau, classe (年级) \\
\hline
环 / 坏 & huán : 王 + 不 / huài : 土 + 不 & anneau, environnement / mauvais, cassé \\
\hline
收 / 放 & shōu : 丩 + 攵(又) / fàng : 方 + 攵 & recevoir, ramasser / poser, libérer \\
\hline
\end{tabularx}
\end{center}

Astuce de mémorisation : les deux membres de chaque paire partagent souvent la même partie phonétique (及, 不, 攵) ; c'est le \cle{radical} qui tranche le sens. Demandez-vous toujours : « quel est le radical ? »

\begin{savaistu}[Un caractère vieux de plus de trois mille ans]
\zh{回} est l'un des caractères les plus anciens de l'écriture chinoise. Sur les os oraculaires de la dynastie Shāng, sa forme ancienne dessinait un \cle{tourbillon d'eau} qui tourne et revient sur lui-même : d'où le sens de « retourner ». Aujourd'hui encore, dans 回收 (recycler), l'idée est la même : ce qui revient pour être réutilisé.
\end{savaistu}

\courssub{Exercice résolu et copie guidée}

\begin{exoresolu}[Compter les traits et choisir le bon caractère]
Énoncé. 1) Combien de traits comptent 回, 收 et 环 ? 2) Complétez avec le bon caractère (环 ou 坏) : 我们要保护环\_\_。/ 这个塑料袋\_\_了，不能用了。
\tcblower
Solution détaillée. 1) 回 compte \textbf{6 traits} (vertical gauche ; haut + droit ; puis le petit 口 en 3 traits ; fermeture finale). 收 compte \textbf{6 traits} (3 à gauche pour 丩, 3 à droite pour la variante de 攵). 环 compte \textbf{8 traits} (4 pour 王 + 4 pour 不). 2) Première phrase : 保护环\textbf{境} (bǎohù huánjìng) : « protéger l'environnement » → il faut 环 (王 + 不). Deuxième phrase : « ce sac en plastique est cassé » → il faut \textbf{坏} (huài, 土 + 不). Le radical 王 renvoie au jade/à l'anneau, le radical 土 à la dégradation : le sens de la phrase décide.
\end{exoresolu}

\begin{exercice}[Copie guidée]
Sur votre cahier, écrivez \textbf{3 lignes} de chaque mot : 回收, 环境, 垃圾. Pour chaque caractère, comptez les traits à voix haute et notez le total à côté (回 : 6 ; 收 : 6 ; 环 : 8 ; 境 : 14 ; 垃 : 8 ; 圾 : 6). Soulignez le radical de chaque caractère et vérifiez que vous avez fermé 回 en dernier.
\end{exercice}

\begin{corrige}[Copie guidée]
Vérification : 回收 = 6 + 6 = 12 traits ; 环境 = 8 + 14 = 22 traits ; 垃圾 = 8 + 6 = 14 traits. Radicaux à souligner : 囗 dans 回, 攵 dans 收, 王 dans 环, 土 dans 境, 土 dans 垃, 土 dans 圾. Erreur fréquente : oublier le dernier trait de fermeture de 回 ou écrire 境 sans le 土 de gauche.
\end{corrige}

\begin{aretenir}[L'essentiel de la section]
\begin{itemize}
\item Chaque caractère = radical (sens) + partie phonétique (son) : analysez avant d'écrire.
\item Ordre des traits : haut → bas, gauche → droite, extérieur → intérieur, fermeture en dernier (回 : 6 traits).
\item 环 (王, environnement) ≠ 坏 (土, mauvais) ; 圾 (土, déchet) ≠ 级 (纟, niveau) ; 收 (recevoir) ≠ 放 (poser).
\item Écrire 5 fois en comptant les traits est la méthode la plus sûre pour mémoriser durablement.
\end{itemize}
\end{aretenir}

\coursec{Thème : français → chinois}

Après avoir appris à écrire correctement les caractères du thème du recyclage (radicaux, ordre des traits, caractères proches), nous passons à l'étape suivante : les utiliser dans des phrases complètes. Le \cle{thème} est l'exercice qui consiste à traduire du français vers le chinois. C'est l'épreuve la plus exigeante du programme de Terminale, car elle mobilise à la fois le vocabulaire, la grammaire et l'écriture des caractères. Bonne nouvelle : avec une méthode rigoureuse, le thème devient un exercice presque mécanique. C'est cette méthode que nous allons construire ensemble.

\courssub{Écriture correcte des caractères clés du thème}

Avant de traduire, rappelons les points de vigilance graphique pour les mots du thème « recyclage ». Un caractère mal formé (radical erroné, trait manquant, ordre incorrect) pénalise la note, même si la phrase est grammaticale.

\begin{aretenir}[Caractères à surveiller : radicaux, traits, confusions]
\begin{itemize}
\item \zh{垃圾} (lājī, déchets) : \zh{垃} (radical \zh{土} tǔ, terre, en bas à gauche) + \zh{圾} (radical \zh{土} aussi). Ordre : haut → bas, gauche → droite. \textbf{Piège :} ne pas confondre \zh{圾} avec \zh{及} (jí) ni \zh{纳} (nà).
\item \zh{分类} (fēnlèi, trier/classer) : \zh{分} (radical \zh{刀} dāo, couteau, 8 traits : horizontal, vertical, pli, points) + \zh{类} (radical \zh{米} mǐ, riz, 18 traits : écrire le toit 米 puis \zh{大} et le trait final). \textbf{Piège :} \zh{类} s'écrit avec \zh{米} en haut, pas \zh{未} (wèi).
\item \zh{塑料} (sùliào, plastique) : \zh{塑} (radical \zh{土} tǔ, 13 traits) + \zh{料} (radical \zh{米} mǐ, 10 traits : \zh{米} + \zh{斗} dǒu). \textbf{Piège :} \zh{料} ≠ \zh{科} (kē, radical \zh{禾}).
\item \zh{回收} (huíshōu, recycler) : \zh{回} (radical \zh{囗} wéi, enclos, 6 traits : l'enclos se ferme en dernier) + \zh{收} (radical \zh{手} shǒu/main à gauche \zh{扌}, 6 traits : horizontal, crochet vertical, horizontal, horizontal, crochet, point). \textbf{Piège majeur :} \zh{收} (recevoir) vs \zh{改} (gǎi, changer, radical \zh{攵} pū) — le trait final diffère (point pour \zh{收}, point oblique pour \zh{改}).
\item \zh{保护} (bǎohù, protéger) : \zh{保} (radical \zh{亻} rén, homme, 9 traits) + \zh{护} (radical \zh{扌} shǒu, 7 traits : \zh{扌} + \zh{户} hù). \textbf{Piège :} \zh{护} (protéger) vs \zh{互} (hù, mutuel, pas de main \zh{扌}).
\item \zh{环境} (huánjìng, environnement) : \zh{环} (radical \zh{王} wáng, jade, 8 traits) + \zh{境} (radical \zh{土} tǔ, 14 traits : \zh{土} + \zh{竟} jìng). \textbf{Piège :} \zh{境} ≠ \zh{镜} (jìng, miroir, radical \zh{钅} jīn).
\item \zh{布袋} (bùdài, sac en tissu) : \zh{布} (radical \zh{巾} jīn, tissu, 5 traits) + \zh{袋} (radical \zh{衣} yī/vêtement \zh{衤} à gauche, 11 traits). \textbf{Piège :} \zh{袋} (sac) vs \zh{代} (dài, génération, radical \zh{人}).
\end{itemize}
\textbf{Rappel ordre des traits universel :} horizontal avant vertical (十) ; gauche avant droite (人) ; haut avant bas (三) ; extérieur avant intérieur (同) ; fermer l'enclos en dernier (回) ; trait central avant les ailes (小).
\end{aretenir}

\courssub{Observer avant de traduire : quatre phrases modèles}

Avant d'énoncer la moindre règle, observons quatre phrases françaises et leur traduction chinoise. Lisez-les attentivement et cherchez ce qui change — et ce qui ne change pas — entre les deux langues.

\begin{exemplebox}[Quatre phrases à observer]
\textbf{Phrase 1.} « Nous devons trier les déchets. »\\
我们应该把垃圾分类。\\
\emph{Wǒmen yīnggāi bǎ lājī fēnlèi.}

\textbf{Phrase 2.} « Les bouteilles en plastique peuvent être recyclées. »\\
塑料瓶可以回收利用。\\
\emph{Sùliàopíng kěyǐ huíshōu lìyòng.}

\textbf{Phrase 3.} « Parce que la Terre est notre maison, nous devons la protéger. »\\
因为地球是我们的家，所以我们要保护它。\\
\emph{Yīnwèi dìqiú shì wǒmen de jiā, suǒyǐ wǒmen yào bǎohù tā.}

\textbf{Phrase 4.} « D'abord utilise un sac en tissu, ensuite jette moins de plastique. »\\
首先用布袋，然后少用塑料。\\
\emph{Shǒuxiān yòng bùdài, ránhòu shǎo yòng sùliào.}
\end{exemplebox}

Que remarquez-vous ?

\begin{itemize}
\item Le sujet vient toujours en premier, comme en français : \zh{我们} (nous), \zh{塑料瓶} (les bouteilles en plastique).
\item Le verbe modal (\zh{应该} « devoir », \zh{可以} « pouvoir », \zh{要} « devoir/vouloir ») se place \textbf{avant} le verbe principal, exactement comme en français (« nous \emph{devons} trier »).
\item Le chinois n'a pas besoin d'article (« les », « un ») ni de conjugaison : le verbe ne change jamais de forme.
\item La voix passive française « être recyclées » disparaît : le chinois dit simplement \zh{可以回收利用} « peuvent recycler et réutiliser », littéralement « peuvent recycler ».
\item Les connecteurs se correspondent : \zh{因为……所以……} = « parce que... donc... », \zh{首先……然后……} = « d'abord... ensuite... ».
\end{itemize}

\courssub{La méthode de traduction en quatre étapes}

Traduire, ce n'est pas remplacer mot à mot : c'est d'abord comprendre le sens, puis reconstruire la phrase selon l'architecture chinoise.

\begin{methode}[Réussir le thème en quatre étapes]
\begin{enumerate}
\item \textbf{Traduire le sens, pas les mots.} Lisez la phrase française entière et demandez-vous : que veut dire cette phrase ? Ne cherchez pas l'équivalent de chaque mot français (articles, auxiliaires « être/avoir ») : beaucoup n'ont pas de correspondant en chinois.
\item \textbf{Repérer les quatre blocs : Sujet / Modal / Verbe / Objet.} Soulignez dans la phrase française : qui fait l'action (S), le verbe modal s'il y en a un (devoir, pouvoir, vouloir — M), le verbe principal (V), et ce sur quoi porte l'action (O).
\item \textbf{Appliquer l'ordre chinois : S + M + V + O.} Placez le modal \textbf{avant} le verbe principal, jamais après. Ajoutez les connecteurs (\zh{因为……所以……}, \zh{首先……然后……}) et les compléments de temps ou de lieu juste après le sujet.
\item \textbf{Vérifier chaque caractère.} Relisez votre phrase caractère par caractère : bon radical ? traits complets ? pas de confusion avec un caractère proche (\zh{塑料} et non une invention ; attention à \zh{收} dans \zh{回收} vs \zh{改}) ?
\end{enumerate}
\end{methode}

\begin{popfigure}
\begin{tikzpicture}[
node distance=0.55cm,
etape/.style={rectangle, rounded corners=3pt, draw=popblue, fill=popblueL, text width=4.6cm, align=center, font=\small, inner sep=5pt},
fleche/.style={-{Stealth[length=2.5mm]}, thick, poporange}
]
\node[etape, align=center] (e1){\textbf{1. Phrase française}\\ « Nous devons trier les déchets. »};
\node[etape, below=of e1, fill=poporangeL, draw=poporange, align=center] (e2){\textbf{2. Repérage des blocs}\\ S = nous \quad M = devons\\ V = trier \quad O = les déchets};
\node[etape, below=of e2, fill=popgreenL, draw=popgreen, align=center] (e3){\textbf{3. Ordre chinois S + M + V + O}\\ 我们 + 应该 + 把垃圾分类};
\node[etape, below=of e3, fill=poppurpleL, draw=poppurple, align=center] (e4){\textbf{4. Vérification des caractères}\\ 我们应该把垃圾分类。};
\draw[fleche] (e1) -- (e2);
\draw[fleche] (e2) -- (e3);
\draw[fleche] (e3) -- (e4);
\end{tikzpicture}
\legende{Organigramme de la méthode de thème : de la phrase française à la phrase chinoise vérifiée.}
\end{popfigure}

\courssub{La règle d'or : la place des verbes modaux}

Les phrases 1, 2 et 3 contiennent toutes un verbe modal. C'est le point de grammaire central de cette section.

\begin{propriete}[Place des modaux : 应该 / 可以 / 要]
En chinois, le verbe modal se place \textbf{toujours AVANT le verbe principal, jamais après}. La structure est :\\
\centerline{\textbf{Sujet + Modal (应该 / 可以 / 要) + Verbe + (Objet)}}
\begin{itemize}
\item \zh{应该} (yīnggāi) = devoir (conseil, obligation morale) : \zh{我们应该保护环境。} \emph{Wǒmen yīnggāi bǎohù huánjìng.} « Nous devons protéger l'environnement. »
\item \zh{可以} (kěyǐ) = pouvoir (possibilité, permission) : \zh{旧报纸可以回收。} \emph{Jiù bàozhǐ kěyǐ huíshōu.} « Les vieux journaux peuvent être recyclés. »
\item \zh{要} (yào) = devoir, vouloir, il faut : \zh{我们要保护它。} \emph{Wǒmen yào bǎohù tā.} « Nous devons la protéger. »
\end{itemize}
\textbf{Négation :} la négation \zh{不} se place devant le modal : \zh{我们不应该浪费水。} \emph{Wǒmen bù yīnggāi làngfèi shuǐ.} « Nous ne devons pas gaspiller l'eau. »
\end{propriete}

\begin{exemplebox}[Exemples traduits et commentés]
\zh{我们应该少用塑料袋。}\\
\emph{Wǒmen yīnggāi shǎo yòng sùliàodài.}\\
« Nous devrions moins utiliser de sacs plastiques. »\\
S = \zh{我们}, M = \zh{应该}, V = \zh{少用} (peu utiliser), O = \zh{塑料袋}. Notez \zh{少} « peu » placé devant le verbe \zh{用} : « moins utiliser » se dit littéralement « peu utiliser ».

\zh{旧报纸可以回收。}\\
\emph{Jiù bàozhǐ kěyǐ huíshōu.}\\
« Les vieux journaux peuvent être recyclés. »\\
S = \zh{旧报纸} (vieux journaux ; l'adjectif \zh{旧} précède le nom sans \zh{的}), M = \zh{可以}, V = \zh{回收}. Le français utilise un passif, le chinois non.
\end{exemplebox}

\courssub{Le grand piège : la voix passive française}

\begin{attention}[« Être recyclé » ne se traduit pas par 是回收 !]
En français, « être » est partout : « les bouteilles \emph{sont} recyclées », « le papier \emph{est} trié ». En chinois, \zh{是} ne s'emploie \textbf{jamais devant un verbe} : il ne relie qu'un sujet à un nom ou à une qualité nominale (\zh{地球是我们的家} « la Terre est notre maison »).

Pour rendre le passif français, deux stratégies :
\begin{enumerate}
\item \textbf{La plus simple (à privilégier au lycée) : modal + verbe actif.} « Peut être recyclé » → \zh{可以回收} ; « doit être protégée » → \zh{应该保护}. Le sens passif est porté par le contexte : ce sont les bouteilles qui subissent le recyclage, c'est évident.
\item \textbf{La structure passive avec 被 (bèi), niveau HSK 4 :} \zh{塑料瓶可以被回收。} \emph{Sùliàopíng kěyǐ bèi huíshōu.} « Les bouteilles en plastique peuvent être recyclées. » Correcte, mais plus lourde ; le chinois préfère souvent l'éviter.
\end{enumerate}
\textbf{Erreurs typiques de francophones :} \zh{塑料瓶是回收} (faux : \zh{是} devant un verbe), \zh{回收可以塑料瓶} (faux : le modal ne se place jamais après le verbe ni avant le sujet).
\end{attention}

\courssub{Analyse guidée des quatre phrases modèles}

Reprenons chaque phrase avec la méthode en quatre étapes.

\textbf{Phrase 1.} « Nous devons trier les déchets. »\\
Repérage : S = nous, M = devons, V = trier, O = les déchets. Le chinois utilise ici la structure avec \zh{把} (bǎ), qui avance l'objet avant le verbe quand on insiste sur ce qu'on fait subir à l'objet : S + 把 + O + V. D'où : \zh{我们应该把垃圾分类。} \emph{Wǒmen yīnggāi bǎ lājī fēnlèi.} Littéralement : « Nous devons, les déchets, classer. » La version sans \zh{把}, \zh{我们应该分类垃圾}, est compréhensible mais moins naturelle : retenez \zh{把垃圾分类} comme une expression toute faite, « trier les déchets ».

\textbf{Phrase 2.} « Les bouteilles en plastique peuvent être recyclées. »\\
Repérage : S = les bouteilles en plastique (\zh{塑料瓶}), M = peuvent (\zh{可以}), V = être recyclées → actif chinois \zh{回收利用} « recycler et réutiliser ». Ordre S + M + V : \zh{塑料瓶可以回收利用。} Le passif français s'efface : c'est le piège signalé plus haut.

\textbf{Phrase 3.} « Parce que la Terre est notre maison, nous devons la protéger. »\\
Le français peut dire « parce que... » seul ; le chinois préfère la paire obligatoire \zh{因为……所以……} « parce que... donc... ». Attention : contrairement au français, on peut (et on doit souvent) garder les deux connecteurs dans la même phrase. « La protéger » : le pronom « la » se rend par \zh{它} (tā, « il/elle » pour les choses), placé après le verbe comme un objet ordinaire : \zh{保护它}. D'où : \zh{因为地球是我们的家，所以我们要保护它。}

\textbf{Phrase 4.} « D'abord utilise un sac en tissu, ensuite jette moins de plastique. »\\
L'impératif français se rend par la phrase sans sujet, exactement comme en français : \zh{首先用布袋，然后少用塑料。} « Sac en tissu » = \zh{布袋} (tissu + sac : le complément du nom « en tissu » précède le nom, sans préposition). « Jette moins de plastique » se dit naturellement \zh{少用塑料} « utilise moins de plastique » : le chinois préfère ici la formulation positive « réduire l'usage ». C'est l'illustration parfaite de l'étape 1 de la méthode : traduire le sens, pas les mots.

\begin{center}
\begin{tabularx}{\linewidth}{|X|X|X|}
\hline
\rowcolor{popblueL} \textbf{Français} & \textbf{Chinois (structure)} & \textbf{Point de grammaire} \\
\hline
Nous devons trier les déchets. & 我们应该把垃圾分类。 & Modal 应该 + structure 把 + O + V \\
\hline
Les bouteilles en plastique peuvent être recyclées. & 塑料瓶可以回收利用。 & Passif français → 可以 + verbe actif \\
\hline
Parce que la Terre est notre maison, nous devons la protéger. & 因为地球是我们的家，所以我们要保护它。 & Paire 因为……所以…… ; pronom objet 它 \\
\hline
D'abord utilise un sac en tissu, ensuite jette moins de plastique. & 首先用布袋，然后少用塑料。 & Impératif sans sujet ; connecteurs 首先 / 然后 \\
\hline
\end{tabularx}
\end{center}

\begin{center}
\begin{tabularx}{\linewidth}{|X|X|X|X|}
\hline
\rowcolor{popblueL} \textbf{Caractères} & \textbf{Pinyin} & \textbf{Nature} & \textbf{Traduction} \\
\hline
应该 & yīnggāi & verbe modal & devoir (conseil) \\
\hline
可以 & kěyǐ & verbe modal & pouvoir \\
\hline
要 & yào & verbe modal & devoir, vouloir, il faut \\
\hline
把 & bǎ & préposition & (marque de l'objet avancé) \\
\hline
垃圾 & lājī & nom & déchets, ordures \\
\hline
分类 & fēnlèi & verbe & trier, classer \\
\hline
塑料瓶 & sùliàopíng & nom & bouteille en plastique \\
\hline
回收利用 & huíshōu lìyòng & verbe & recycler et réutiliser \\
\hline
地球 & dìqiú & nom & la Terre \\
\hline
保护 & bǎohù & verbe & protéger \\
\hline
它 & tā & pronom & il, elle (chose) \\
\hline
布袋 & bùdài & nom & sac en tissu \\
\hline
少用 & shǎo yòng & verbe & utiliser moins \\
\hline
\end{tabularx}
\end{center}

\courssub{Exercice résolu et entraînement individuel}

\begin{exoresolu}[Thème guidé : « Nous devons protéger l'environnement de Yaoundé. »]
Traduisez en chinois : « Nous devons protéger l'environnement de Yaoundé. »
\tcblower
\textbf{Étape 1 — Sens :} une obligation (devoir) portant sur une action (protéger) et un objet (l'environnement de Yaoundé).\\
\textbf{Étape 2 — Repérage :} S = nous → \zh{我们} ; M = devons → \zh{应该} ; V = protéger → \zh{保护} ; O = l'environnement de Yaoundé → \zh{雅温得的环境} (le complément du nom « de Yaoundé » précède le nom avec \zh{的}).\\
\textbf{Étape 3 — Ordre chinois S + M + V + O :} \zh{我们 + 应该 + 保护 + 雅温得的环境}。\\
\textbf{Étape 4 — Vérification :} \zh{环} a le radical du jade \zh{王} à gauche ; \zh{境} a le radical de la terre \zh{土} ; \zh{保} et \zh{护} sont complets.\\
\textbf{Traduction finale :} 我们应该保护雅温得的环境。\\
\emph{Wǒmen yīnggāi bǎohù Yǎwēndé de huánjìng.}
\end{exoresolu}

\begin{exercice}[Entraînement individuel — 10 minutes]
Traduisez les trois phrases suivantes en chinois, sur votre cahier, en appliquant les quatre étapes de la méthode. La correction se fera collectivement au tableau.
\begin{enumerate}
\item « Nous devons utiliser moins de sacs en plastique. »
\item « Le papier et le verre peuvent être recyclés. »
\item « Parce que les déchets polluent la ville, nous devons les trier. »
\end{enumerate}
\end{exercice}

\begin{corrige}[Entraînement individuel]
\textbf{1.} S = \zh{我们}, M = \zh{应该}, V = \zh{少用}, O = \zh{塑料袋}.\\
我们应该少用塑料袋。 \emph{Wǒmen yīnggāi shǎo yòng sùliàodài.}\\
\textbf{2.} S = \zh{纸和玻璃} (le papier et le verre), M = \zh{可以}, V = \zh{回收} ; le passif français « être recyclés » disparaît.\\
纸和玻璃可以回收。 \emph{Zhǐ hé bōli kěyǐ huíshōu.}\\
\textbf{3.} Paire \zh{因为……所以……} ; « les trier » = \zh{把它们分类} (structure 把 + pronom objet).\\
因为垃圾污染城市，所以我们要把它们分类。 \emph{Yīnwèi lājī wūrǎn chéngshì, suǒyǐ wǒmen yào bǎ tāmen fēnlèi.}
\end{corrige}

\begin{aretenir}[L'essentiel du thème]
\begin{itemize}
\item Méthode en quatre étapes : sens → repérage S/M/V/O → ordre chinois → vérification des caractères.
\item Ordre de base : \textbf{Sujet + Modal + Verbe + Objet}. Le modal (\zh{应该}, \zh{可以}, \zh{要}) précède toujours le verbe.
\item Le passif français « être + participe » se rend par \zh{可以/应该} + verbe actif, jamais par \zh{是} + verbe.
\item Les paires de connecteurs \zh{因为……所以……} et \zh{首先……然后……} s'emploient ensemble.
\item Dernière relecture obligatoire : un caractère mal écrit (radical faux, trait manquant) coûte des points même si la grammaire est juste.
\end{itemize}
\end{aretenir}

Vous êtes maintenant capables de construire des phrases chinoises à partir du français. Dans la section suivante, nous ferons le chemin inverse : la \cle{version}, c'est-à-dire la traduction du chinois vers le français, qui demande d'autres réflexes.

\coursec{Version : chinois → français}

Après avoir appris à traduire du français vers le chinois (le \emph{thème}), nous faisons maintenant le chemin inverse : la \cle{version}, c'est-à-dire la traduction du chinois vers le français. C'est une compétence essentielle : elle prouve que vous avez \emph{compris} le texte chinois dans son ensemble, et pas seulement reconnu quelques mots isolés. Nous allons travailler sur un texte court consacré au tri des déchets et au recyclage en Chine, puis apprendre à restituer ce texte dans un français correct, naturel et élégant.

\courssub{Observer d'abord : le texte support}

Lisez le texte ci-dessous en entier, \emph{avant} de chercher à traduire quoi que ce soit. En version, la première lecture sert toujours à saisir le sens global : de quoi parle le texte ? Qui fait quoi ? Quelle est l'idée finale ?

\begin{textebox}[Le tri des déchets en Chine]
在中国，很多城市都有垃圾分类。人们把垃圾分成四种。旧东西可以回收再利用。这样既保护环境，又节约资源。\\
\emph{Zài Zhōngguó, hěn duō chéngshì dōu yǒu lājī fēnlèi. Rénmen bǎ lājī fēn chéng sì zhǒng. Jiù dōngxi kěyǐ huíshōu zài lìyòng. Zhèyàng jì bǎohù huánjìng, yòu jiéyuē zīyuán.}\\
« En Chine, de nombreuses villes pratiquent le tri des déchets. Les gens divisent les déchets en quatre catégories. Les vieux objets peuvent être recyclés et réutilisés. Ainsi, cela protège à la fois l'environnement et économise les ressources. »
\end{textebox}

\textbf{Glossaire des caractères difficiles}\par

\begin{center}
\begin{tabularx}{\linewidth}{|X|X|X|X|}
\hline
\rowcolor{popblueL} Caractères & Pinyin & Nature & Traduction \\
\hline
种 & zhǒng & nom / classif. & sorte, catégorie \\
\hline
节约 & jiéyuē & verbe & économiser, épargner \\
\hline
资源 & zīyuán & nom & ressources \\
\hline
回收 & huíshōu & verbe & recycler, récupérer \\
\hline
再利用 & zài lìyòng & verbe & réutiliser \\
\hline
旧 & jiù & adjectif & vieux, usagé \\
\hline
既 & jì & conjonction & à la fois (avec 又) \\
\hline
\end{tabularx}
\end{center}

\begin{attention}[种 : deux prononciations à ne pas confondre]
Le caractère \zh{种} possède deux lectures. Il se prononce \emph{zhǒng} (3\ieme{} ton) quand il signifie « sorte, espèce » : \zh{四种} \emph{sì zhǒng} « quatre sortes ». Il se prononce \emph{zhòng} (4\ieme{} ton) quand il signifie « planter, cultiver » : \zh{种树} \emph{zhòng shù} « planter des arbres ». Dans notre texte, \zh{四种} se lit donc \emph{sì zhǒng}. Erreur fréquente des francophones : lire mécaniquement \emph{zhòng} partout parce que cette lecture a été vue en premier dans le mot \zh{种花} « planter des fleurs ».
\end{attention}

\courssub{La méthode de la version, étape par étape}

Traduire du chinois vers le français ne consiste pas à remplacer mot à mot. Le chinois et le français n'ont ni les mêmes articles, ni les mêmes accords, ni le même goût pour la répétition. Voici la démarche complète.

\begin{methode}[Réussir une version en cinq étapes]
\begin{enumerate}
\item \textbf{Lire le texte entier} (au moins deux fois) et formuler en une phrase française le sujet général : ici, « le tri des déchets en Chine et ses bienfaits ».
\item \textbf{Repérer les connecteurs} : \zh{这样} « ainsi, de cette façon », \zh{既……又……} « à la fois... et... », \zh{把} (construction de disposition), \zh{可以} « pouvoir ». Ils donnent la charpente logique du texte.
\item \textbf{Découper en blocs de sens} : ne traduisez jamais caractère par caractère. Un bloc = une unité de sens (sujet + action, ou verbe + complément).
\item \textbf{Traduire d'abord le sens global} de chaque bloc, même maladroitement, en français « brut ».
\item \textbf{Polir le français} : ajouter les articles (\emph{le, la, les}), accorder les participes passés, choisir un vocabulaire naturel, éviter les répétitions. Un bon français vaut mieux qu'un français « collé » au chinois.
\end{enumerate}
\end{methode}

\courssub{Traduction guidée phrase par phrase}

\textbf{Phrase 1 :} \zh{在中国，很多城市都有垃圾分类。}\par
\emph{Zài Zhōngguó, hěn duō chéngshì dōu yǒu lājī fēnlèi.}\par
Blocs de sens : \zh{在中国} « en Chine » + \zh{很多城市} « de nombreuses villes » + \zh{都有} « ont toutes / pratiquent » + \zh{垃圾分类} « le tri des déchets ».\par
Traduction brute : « En Chine, beaucoup de villes ont le tri des ordures. »\par
Traduction polie : « En Chine, de nombreuses villes pratiquent le tri des déchets. » (On préfère \emph{pratiquent} à \emph{ont}, plus naturel en français, et \emph{déchets} à \emph{ordures} dans un texte écrit.)

\textbf{Phrase 2 :} \zh{人们把垃圾分成四种。}\par
\emph{Rénmen bǎ lājī fēn chéng sì zhǒng.}\par
C'est une construction en \zh{把} : Sujet + \zh{把} + objet + verbe + résultat. Littéralement : « Les gens, les déchets, (ils) divisent en quatre sortes. »\par
Traduction polie : « Les gens divisent les déchets en quatre catégories. » Le français remet l'objet \emph{après} le verbe : ne conservez jamais l'ordre chinois \zh{把} en français !

\textbf{Phrase 3 :} \zh{旧东西可以回收再利用。}\par
\emph{Jiù dōngxi kěyǐ huíshōu zài lìyòng.}\par
\zh{旧东西} « les vieux objets » + \zh{可以} « peuvent » + \zh{回收} « être recyclés » + \zh{再利用} « être réutilisés ». Le chinois n'a pas de passif marqué ici ; le français, lui, préfère le passif : « Les vieux objets peuvent être recyclés et réutilisés. »

\textbf{Phrase 4 :} \zh{这样既保护环境，又节约资源。}\par
\emph{Zhèyàng jì bǎohù huánjìng, yòu jiéyuē zīyuán.}\par
Traduction : « Ainsi, cela protège à la fois l'environnement et économise les ressources. » Cette phrase contient la structure clé de la leçon.

\begin{propriete}[La structure 既……又…… (jì...yòu...)]
La structure \zh{既} A \zh{又} B signifie « à la fois A et B ». Elle relie deux éléments de même nature : deux verbes ou deux adjectifs, et souligne que les deux faits sont vrais en même temps.\par
Schéma : Sujet + \zh{既} + verbe/adjectif 1，\zh{又} + verbe/adjectif 2.\par
Emploi : registre plutôt écrit ou soigné ; à l'oral, on utilise souvent \zh{又……又……}. Elle ne relie pas deux phrases complètes avec deux sujets différents.
\end{propriete}

\begin{exemplebox}[既……又…… en action]
\zh{这样既保护环境，又节约资源。} \emph{Zhèyàng jì bǎohù huánjìng, yòu jiéyuē zīyuán.} « Ainsi, cela protège à la fois l'environnement et économise les ressources. »\\
\zh{回收既简单又重要。} \emph{Huíshōu jì jiǎndān yòu zhòngyào.} « Le recyclage est à la fois simple et important. »\\
\zh{这个城市既大又干净。} \emph{Zhège chéngshì jì dà yòu gānjìng.} « Cette ville est à la fois grande et propre. »
\end{exemplebox}

\courssub{Tableau de correspondance : bloc chinois → bloc français}

La figure suivante visualise le découpage en blocs de sens et la transformation vers le français : remarquez comment l'ordre change dans la phrase en \zh{把}.

\begin{popfigure}
\begin{tikzpicture}[
  zh/.style={rectangle, rounded corners=2pt, draw=popblue, fill=popblueL, align=center, font=\small, minimum height=0.9cm, text width=3.6cm},
  fr/.style={rectangle, rounded corners=2pt, draw=popgreen, fill=popgreenL, align=center, font=\small, minimum height=0.9cm, text width=4.6cm},
  arr/.style={-{Stealth[length=2.5mm]}, thick, poporange}
]
\node[zh, align=center] (z1) at (0,0){在中国，很多城市\\都有垃圾分类。};
\node[fr, right=2.2cm of z1] (f1) {En Chine, de nombreuses villes pratiquent le tri des déchets.};
\draw[arr] (z1) -- (f1);
\node[zh, below=0.5cm of z1, align=center] (z2){人们把垃圾\\分成四种。};
\node[fr, right=2.2cm of z2] (f2) {Les gens divisent les déchets en quatre catégories.};
\draw[arr] (z2) -- (f2);
\node[zh, below=0.5cm of z2, align=center] (z3){旧东西可以\\回收再利用。};
\node[fr, right=2.2cm of z3] (f3) {Les vieux objets peuvent être recyclés et réutilisés.};
\draw[arr] (z3) -- (f3);
\node[zh, below=0.5cm of z3, align=center] (z4){这样既保护环境，\\又节约资源。};
\node[fr, right=2.2cm of z4] (f4) {Ainsi, cela protège à la fois l'environnement et économise les ressources.};
\draw[arr] (z4) -- (f4);
\end{tikzpicture}
\legende{Du bloc chinois au bloc français : quatre lignes, quatre transformations.}
\end{popfigure}

\courssub{Comparer les deux langues pour mieux traduire}

\begin{center}
\begin{tabularx}{\linewidth}{|X|X|X|}
\hline
\rowcolor{popblueL} Point de langue & En chinois & En français \\
\hline
Articles & absents (\zh{垃圾} seul) & obligatoires : \emph{les déchets} \\
\hline
Passif & souvent non marqué (\zh{可以回收}) & passif avec \emph{être} : \emph{peuvent être recyclés} \\
\hline
Objet avec 把 & avant le verbe : \zh{把垃圾分成} & après le verbe : \emph{divisent les déchets} \\
\hline
Pluriel & non marqué (\zh{四种} sans -s) & marqué : \emph{quatre catégories} \\
\hline
Connecteur double & \zh{既……又……} encadre les deux termes & \emph{à la fois... et...} \\
\hline
\end{tabularx}
\end{center}

\begin{attention}[Pièges de version pour les francophones]
\begin{itemize}
\item Ne traduisez pas \zh{都} par « tout » : dans \zh{都有}, il renforce simplement l'idée de généralité ; « pratiquent » ou « ont toutes » suffit.
\item \zh{东西} (dōngxi) signifie « chose, objet » ; ne le décomposez jamais en « est-ouest » !
\item \zh{再} dans \zh{再利用} exprime la répétition (« ré- ») : \emph{réutiliser}, et non « utiliser encore une fois » mot à mot.
\item Évitez le style télégraphique : « En Chine beaucoup villes ont tri déchets » n'est pas du français. Ajoutez articles et liaisons.
\end{itemize}
\end{attention}

\begin{exoresolu}[Version guidée]
Énoncé : traduisez en français soigné la phrase \zh{回收旧东西既保护环境，又节约资源。}\par
\tcblower
Solution détaillée. Étape 1 : sens global — le recyclage a deux avantages. Étape 2 : connecteurs — \zh{既……又……} = « à la fois... et... ». Étape 3 : blocs — \zh{回收旧东西} « recycler les vieux objets » ; \zh{保护环境} « protéger l'environnement » ; \zh{节约资源} « économiser les ressources ». Étape 4 : traduction brute — « Recycler vieilles choses à la fois protège environnement, et économise ressources. » Étape 5 : polissage — « Recycler les vieux objets permet à la fois de protéger l'environnement et d'économiser les ressources. » Notez l'ajout des articles, l'infinitif français après \emph{permet de}, et l'accord harmonieux du style.
\end{exoresolu}

\begin{exercice}[À vous de traduire]
Traduisez en français : 1) \zh{我们学校也有垃圾分类。} 2) \zh{人们把旧瓶子回收再利用。} 3) \zh{这个城市既干净又美丽。}
\end{exercice}

\begin{corrige}[Exercice « À vous de traduire »]
1) « Notre école aussi pratique le tri des déchets. » (\zh{也} = « aussi » ; ne l'oubliez pas.) 2) « Les gens recyclent et réutilisent les vieilles bouteilles. » (construction en \zh{把} : l'objet repasse après le verbe en français.) 3) « Cette ville est à la fois propre et belle. » (\zh{既……又……} relie ici deux adjectifs.)
\end{corrige}

\begin{experience}[Restitution orale : comparer et améliorer]
Un élève lit à voix haute sa traduction complète du texte support. La classe écoute, puis compare avec sa propre version : qui a choisi « ordures », qui a choisi « déchets » ? Qui a écrit « beaucoup de villes », qui a écrit « de nombreuses villes » ? Ensemble, proposez une version collective améliorée au tableau, en veillant au style : articles corrects, vocabulaire soutenu, phrases fluides. Le professeur souligne enfin les choix qui rendent le français plus naturel.
\end{experience}

\begin{aretenir}[L'essentiel de la version]
Lire le texte entier, repérer les connecteurs (\zh{这样}, \zh{既……又……}), découper en blocs de sens, traduire le sens global, puis polir le français. La structure \zh{既}A\zh{又}B = « à la fois A et B ». Attention à \zh{种} : \emph{zhǒng} (sorte) / \emph{zhòng} (planter). Une bonne version est un texte français qui se lit naturellement, pas un calque du chinois.
\end{aretenir}

\coursec{Production guidée et grille d'évaluation}

Après avoir traduit des phrases du chinois vers le français (version) et du français vers le chinois (thème), tu disposes maintenant de tout le matériel linguistique nécessaire : le vocabulaire du recyclage, les structures grammaticales, les connecteurs chronologiques. Il est temps de passer à l'étape finale et la plus importante : \cle{produire toi-même un texte simple en chinois}. C'est exactement ce que l'on te demandera au Baccalauréat. Cette section te guide pas à pas, du plan jusqu'à la correction.

\courssub{Le sujet de production}

\begin{definition}[Sujet officiel de la production]
Rédige un texte de \textbf{5 à 8 phrases} en chinois sur le thème : « \zh{在我的小区} » (\emph{Zài wǒ de xiǎoqū.}) ou « \zh{在我们学校} » (\emph{Zài wǒmen xuéxiào.}) — « Le recyclage dans mon quartier / dans notre école ».
\begin{itemize}
\item \textbf{Plan imposé} : introduction (1 phrase) $\rightarrow$ \zh{首先} $\rightarrow$ \zh{然后} $\rightarrow$ \zh{最后} $\rightarrow$ conclusion avec \zh{因为……所以……}.
\item \textbf{Banque de mots obligatoires} : \zh{回收}, \zh{分类}, \zh{塑料}, \zh{保护}, \zh{环境}, \zh{应该}, \zh{可以}.
\item \textbf{Durée} : 20 minutes, rédaction individuelle sur papier, caractères soignés.
\end{itemize}
\end{definition}

\begin{center}
\begin{tabularx}{\linewidth}{|X|X|X|X|}
\hline
\rowcolor{popblueL} Caractères & Pinyin & Nature & Traduction \\
\hline
回收 & huíshōu & verbe & recycler, récupérer \\
\hline
分类 & fēnlèi & verbe / nom & trier ; le tri \\
\hline
塑料 & sùliào & nom & le plastique \\
\hline
保护 & bǎohù & verbe & protéger \\
\hline
环境 & huánjìng & nom & l'environnement \\
\hline
应该 & yīnggāi & modal & devoir (il faut) \\
\hline
可以 & kěyǐ & modal & pouvoir (permission, possibilité) \\
\hline
首先 & shǒuxiān & connecteur & d'abord, premièrement \\
\hline
然后 & ránhòu & connecteur & ensuite, puis \\
\hline
最后 & zuìhòu & connecteur & enfin, finalement \\
\hline
\end{tabularx}
\end{center}

\courssub{Observer avant d'écrire : un modèle de copie réussie}

Avant de rédiger, lis ce modèle de production d'élève. Observe comment le plan est visible, comment les connecteurs ouvrent chaque étape et comment la conclusion utilise \zh{因为……所以……}.

\begin{textebox}[Modèle de production : 我们学校的回收]
在我们学校，回收很重要。\\
\emph{Zài wǒmen xuéxiào, huíshōu hěn zhòngyào.}\\
« Dans notre école, le recyclage est très important. »\\[4pt]
首先，我们把垃圾分类：塑料、纸和瓶子。\\
\emph{Shǒuxiān, wǒmen bǎ lājī fēnlèi : sùliào, zhǐ hé píngzi.}\\
« D'abord, nous trions les déchets : le plastique, le papier et les bouteilles. »\\[4pt]
然后，我们可以把塑料瓶卖给回收公司。\\
\emph{Ránhòu, wǒmen kěyǐ bǎ sùliào píng mài gěi huíshōu gōngsī.}\\
« Ensuite, nous pouvons vendre les bouteilles en plastique à une entreprise de recyclage. »\\[4pt]
最后，老师告诉我们为什么要保护环境。\\
\emph{Zuìhòu, lǎoshī gàosu wǒmen wèishénme yào bǎohù huánjìng.}\\
« Enfin, le professeur nous explique pourquoi il faut protéger l'environnement. »\\[4pt]
因为干净的学校让我们健康，所以我们应该每天回收。\\
\emph{Yīnwèi gānjìng de xuéxiào ràng wǒmen jiànkāng, suǒyǐ wǒmen yīnggāi měitiān huíshōu.}\\
« Parce qu'une école propre nous garde en bonne santé, nous devons donc recycler chaque jour. »
\end{textebox}

\textbf{Observations guidées :}
\begin{enumerate}
\item Combien de phrases compte ce texte ? (5 phrases : le minimum demandé.)
\item Souligne les trois connecteurs chronologiques : ils ouvrent les phrases 2, 3 et 4.
\item Retrouve les sept mots obligatoires : \zh{回收} (2 fois), \zh{分类}, \zh{塑料} (2 fois), \zh{保护}, \zh{环境}, \zh{应该}, \zh{可以}. Tous sont présents.
\item La dernière phrase commence par \zh{因为} et contient \zh{所以} : c'est la conclusion imposée.
\end{enumerate}

\courssub{La frise du plan de rédaction}

Le plan imposé est une \emph{progression chronologique} : on présente le sujet, on raconte les étapes du recyclage dans l'ordre, puis on conclut avec une cause et une conséquence. Mémorise cette frise ; elle est valable pour presque tous les sujets d'expression écrite du programme.

\begin{popfigure}
\begin{tikzpicture}[
  etape/.style={rectangle, rounded corners=4pt, draw=popblue, fill=popblueL, thick, minimum width=2.4cm, minimum height=1.1cm, align=center, font=\small},
  etapeconc/.style={rectangle, rounded corners=4pt, draw=poppurple, fill=poppurpleL, thick, minimum width=2.4cm, minimum height=1.1cm, align=center, font=\small},
  fleche/.style={-{Stealth[length=3mm]}, thick, poporange}
]
\node[etape, align=center] (intro) at (0,0){\zh{引言}\\ Introduction\\ (1 phrase)};
\node[etape, right=0.7cm of intro, align=center] (shou){\zh{首先}\\ D'abord\\ (le tri)};
\node[etape, right=0.7cm of shou, align=center] (ran){\zh{然后}\\ Ensuite\\ (la récupération)};
\node[etape, right=0.7cm of ran, align=center] (zui){\zh{最后}\\ Enfin\\ (le résultat)};
\node[etapeconc, right=0.7cm of zui, align=center] (conc){\zh{因为…所以…}\\ Conclusion\\ (cause $\to$ conséquence)};
\draw[fleche] (intro) -- (shou);
\draw[fleche] (shou) -- (ran);
\draw[fleche] (ran) -- (zui);
\draw[fleche] (zui) -- (conc);
\end{tikzpicture}
\legende{Frise du plan de rédaction : Introduction $\rightarrow$ 首先 $\rightarrow$ 然后 $\rightarrow$ 最后 $\rightarrow$ Conclusion}
\end{popfigure}

\begin{methode}[Comment rédiger en 20 minutes : la méthode en 4 étapes]
\begin{enumerate}
\item \textbf{Écrire d'abord le plan en connecteurs (3 minutes).} Sur ton brouillon, écris verticalement : \zh{引言} / \zh{首先} / \zh{然后} / \zh{最后} / \zh{因为……所以……}. Ton squelette est prêt : tu ne te demanderas plus « qu'est-ce que j'écris ensuite ? ».
\item \textbf{Remplir avec les structures apprises (10 minutes).} Pour chaque étape, utilise une structure vue en classe : \zh{我们 + 把 + nom + 分类} (structure en \zh{把}), \zh{我们可以 + verbe} (modal + verbe), \zh{我们应该 + verbe}. Place les mots obligatoires un par un et coche-les au brouillon.
\item \textbf{Rédiger au propre avec des caractères soignés (5 minutes).} Écris lentement, trait par trait : horizontal avant vertical, haut avant bas, gauche avant droite. Un caractère lisible vaut mieux qu'un caractère rapide.
\item \textbf{Relire en vérifiant les caractères (2 minutes).} Vérifie les radicaux (\zh{环} = clé du jade 王, pas 王 manquant ; \zh{塑} ne s'écrit pas sans 土 en bas), les tons des modaux et la place de \zh{所以} dans la conclusion.
\end{enumerate}
\end{methode}

\begin{exemplebox}[Phrases de conclusion modèles à réutiliser]
\zh{保护环境就是保护我们自己。}\\
\emph{Bǎohù huánjìng jiùshì bǎohù wǒmen zìjǐ.}\\
« Protéger l'environnement, c'est nous protéger nous-mêmes. »\\[4pt]
\zh{因为塑料污染环境，所以我们应该回收它。}\\
\emph{Yīnwèi sùliào wūrǎn huánjìng, suǒyǐ wǒmen yīnggāi huíshōu tā.}\\
« Parce que le plastique pollue l'environnement, nous devons donc le recycler. »\\[4pt]
\zh{回收很小，但是意义很大。}\\
\emph{Huíshōu hěn xiǎo, dànshì yìyì hěn dà.}\\
« Le recyclage est un petit geste, mais son sens est grand. »
\end{exemplebox}

\courssub{Les critères d'une bonne copie}

\begin{propriete}[Critères d'une bonne copie d'expression écrite au Bac]
Une copie réussie respecte quatre exigences :
\begin{enumerate}
\item \textbf{Un plan visible} : le correcteur doit voir immédiatement l'introduction, les étapes et la conclusion. Les connecteurs \zh{首先}, \zh{然后}, \zh{最后} servent de « panneaux » : \textbf{2 connecteurs minimum} sont exigés.
\item \textbf{Des modaux correctement placés} : en chinois, le modal se place \emph{avant le verbe principal}, jamais après : \zh{我们应该保护环境} (\emph{Wǒmen yīnggāi bǎohù huánjìng.}) et non \zh{我们保护环境应该}. En français, l'ordre est identique : « nous devons protéger » ; profitez-en pour faire le parallèle !
\item \textbf{Des caractères lisibles} : traits complets, radicaux corrects, taille régulière.
\item \textbf{Le respect du sujet} : chaque phrase parle du recyclage dans ton quartier ou ton école ; aucune phrase hors sujet.
\end{enumerate}
\end{propriete}

\begin{attention}[Au Bac : un caractère mal écrit est compté faux]
Au Baccalauréat, un caractère avec un \textbf{trait manquant} ou un \textbf{mauvais radical} est compté \emph{faux}, même si le sens est devinable. Pièges fréquents des francophones :
\begin{itemize}
\item \zh{环} (huán, environnement) : clé du jade 王 à gauche + \zh{不} à droite. Ne pas écrire 坏 (huài, mauvais) : un seul trait change tout !
\item \zh{塑} (sù, plastique) : ne pas oublier la clé de la terre 土 en bas.
\item \zh{类} (lèi, catégorie) : \zh{米} en haut + \zh{大} en bas ; ne pas confondre avec \zh{粪} ni oublier le point final du 大.
\item \zh{该} (gāi, dans 应该) : clé de la parole 讠à gauche ; ne pas écrire 孩 (hái, enfant).
\end{itemize}
Quand tu hésites sur un caractère, choisis un synonyme plus simple que tu maîtrises : mieux vaut \zh{干净} (propre) bien écrit qu'un caractère savant raté.
\end{attention}

\courssub{La grille d'évaluation sur 20 points}

Ta copie sera notée sur 20 points selon cinq critères. Étudie cette grille : elle te dit exactement où se gagnent les points.

\begin{popfigure}
\begin{tikzpicture}
\fill[popblue] (0,0) rectangle (10.5,0.8);
\node[text=white, font=\bfseries\small, anchor=west] at (0.2,0.4) {Critère};
\node[text=white, font=\bfseries\small, anchor=east] at (10.3,0.4) {Points};
\fill[popblueL] (0,-0.8) rectangle (10.5,0);
\node[font=\small, anchor=west] at (0.2,-0.4) {Caractères corrects (traits, radicaux, lisibilité)};
\node[font=\small\bfseries, anchor=east] at (10.3,-0.4) {5};
\fill[white] (0,-1.6) rectangle (10.5,-0.8);
\node[font=\small, anchor=west] at (0.2,-1.2) {Grammaire et structures (把, modaux, 因为…所以…)};
\node[font=\small\bfseries, anchor=east] at (10.3,-1.2) {5};
\fill[poporangeL] (0,-2.4) rectangle (10.5,-1.6);
\node[font=\small, anchor=west] at (0.2,-2.0) {Connecteurs (首先, 然后, 最后 — 2 minimum)};
\node[font=\small\bfseries, anchor=east] at (10.3,-2.0) {4};
\fill[white] (0,-3.2) rectangle (10.5,-2.4);
\node[font=\small, anchor=west] at (0.2,-2.8) {Respect du sujet et du sens (mots obligatoires employés)};
\node[font=\small\bfseries, anchor=east] at (10.3,-2.8) {4};
\fill[popgreenL] (0,-4.0) rectangle (10.5,-3.2);
\node[font=\small, anchor=west] at (0.2,-3.6) {Présentation (copie propre, ponctuation chinoise)};
\node[font=\small\bfseries, anchor=east] at (10.3,-3.6) {2};
\draw[popdark, thick] (0,0.8) rectangle (10.5,-4.0);
\draw[popdark] (0,0) -- (10.5,0);
\foreach \y in {-0.8,-1.6,-2.4,-3.2} \draw[popline] (0,\y) -- (10.5,\y);
\node[font=\small\bfseries, anchor=east, text=popdark] at (10.3,-4.5) {Total : 20};
\end{tikzpicture}
\legende{Grille d'évaluation de la production écrite (sur 20 points)}
\end{popfigure}

\begin{center}
\begin{tabularx}{\linewidth}{|X|X|X|}
\hline
\rowcolor{popblueL} Critère & Ce que le correcteur regarde & Erreur qui fait perdre des points \\
\hline
Caractères (5) & traits complets, bons radicaux & trait manquant, caractère inventé \\
\hline
Grammaire (5) & ordre des mots, modaux avant le verbe, 把 bien construit & \zh{我们保护环境应该} (modal mal placé) \\
\hline
Connecteurs (4) & 首先, 然后, 最后 présents et bien utilisés & texte sans aucun connecteur \\
\hline
Sujet et sens (4) & les 7 mots obligatoires, phrases sur le sujet & phrase hors sujet, mot obligatoire absent \\
\hline
Présentation (2) & copie propre, ponctuation 。， & ratures, ponctuation française \\
\hline
\end{tabularx}
\end{center}

\courssub{Exercice résolu : évaluer une copie avec la grille}

\begin{exoresolu}[Corriger la copie de Marie comme un examinateur]
\textbf{Énoncé.} Voici la copie de Marie (5 phrases). Note-la sur 20 avec la grille, puis propose une correction.\\[3pt]
\zh{在我的小区，回收很重要。首先，我们把垃圾分类。然后，塑料瓶子可以回收。最后，我们保护环境应该。因为回收好，所以我们每天做。}\\[3pt]
\emph{Zài wǒ de xiǎoqū, huíshōu hěn zhòngyào. Shǒuxiān, wǒmen bǎ lājī fēnlèi. Ránhòu, sùliào píngzi kěyǐ huíshōu. Zuìhòu, wǒmen bǎohù huánjìng yīnggāi. Yīnwèi huíshōu hǎo, suǒyǐ wǒmen měitiān zuò.}
\tcblower
\textbf{Solution détaillée, critère par critère :}
\begin{itemize}
\item \textbf{Caractères (5/5)} : tous les caractères sont corrects et lisibles.
\item \textbf{Grammaire (3/5)} : la phrase 4 \zh{我们保护环境应该} est fausse : le modal \zh{应该} doit précéder le verbe. Correction : \zh{我们应该保护环境。} (\emph{Wǒmen yīnggāi bǎohù huánjìng.}) « Nous devons protéger l'environnement. »
\item \textbf{Connecteurs (4/4)} : \zh{首先}, \zh{然后}, \zh{最后} sont tous présents et bien placés.
\item \textbf{Sujet et sens (3/4)} : six mots obligatoires sur sept ; \zh{环境} est présent, mais la phrase 5 \zh{因为回收好} est vague : « parce que recycler c'est bien » manque de précision. Mieux : \zh{因为回收可以保护环境} (« parce que le recyclage peut protéger l'environnement »).
\item \textbf{Présentation (2/2)} : copie propre, ponctuation chinoise correcte.
\end{itemize}
\textbf{Note finale : 5 + 3 + 4 + 3 + 2 = 17/20.} La copie corrigée :\\
\zh{在我的小区，回收很重要。首先，我们把垃圾分类。然后，塑料瓶子可以回收。最后，我们应该保护环境。因为回收可以保护环境，所以我们应该每天回收。}
\end{exoresolu}

\courssub{Déroulement de la séance : rédaction et correction croisée}

\begin{experience}[Rédaction individuelle puis correction croisée (20 min + 15 min)]
\begin{enumerate}
\item \textbf{Rédaction individuelle (20 minutes).} Chaque élève rédige sur papier, en silence, avec des caractères soignés, en suivant le plan imposé et en utilisant les sept mots obligatoires. Le dictionnaire de la banque de mots est autorisé, pas le voisin !
\item \textbf{Correction croisée entre pairs (10 minutes).} Échange ta copie avec ton voisin. Avec la grille, note chaque critère au crayon et écris un commentaire gentil et précis : « ton \zh{应该} est mal placé », « il manque le mot \zh{分类} », « tes caractères sont très lisibles ».
\item \textbf{Correction collective (5 minutes).} Le professeur projette une copie (anonyme) au tableau ; la classe la corrige ensemble, critère par critère, et propose une version améliorée.
\end{enumerate}
\end{experience}

\begin{exercice}[À toi de jouer : ta production personnelle]
Rédige maintenant ton propre texte « \zh{在我的小区} » ou « \zh{在我们学校} » en 5 à 8 phrases. Avant de commencer, écris ton plan en connecteurs au brouillon. Après avoir fini, auto-évalue-toi avec la grille : combien de points te donnes-tu sur 20, et pourquoi ?
\end{exercice}

\begin{corrige}[Exercice : éléments de correction attendus]
Il n'existe pas une seule bonne réponse, mais toute bonne copie contient :
\begin{enumerate}
\item Une introduction d'une phrase, par exemple : \zh{在我的小区，回收越来越重要。} (\emph{Zài wǒ de xiǎoqū, huíshōu yuèláiyuè zhòngyào.}) « Dans mon quartier, le recyclage est de plus en plus important. »
\item \zh{首先} + le tri : \zh{首先，人们把垃圾分类。}
\item \zh{然后} + une action avec \zh{可以} : \zh{然后，我们可以把塑料卖给回收公司。}
\item \zh{最后} + un résultat : \zh{最后，我们的小区变得更干净。}
\item Une conclusion en \zh{因为……所以……} avec \zh{应该} : \zh{因为保护环境就是保护我们自己，所以我们应该每天回收。}
\end{enumerate}
Vérifie : les 7 mots obligatoires sont-ils tous présents ? Les modaux \zh{应该} et \zh{可以} précèdent-ils bien le verbe ? Si oui, ta copie vise 18/20 ou plus.
\end{corrige}

\begin{savaistu}[Le recyclage au Cameroun : un exemple à citer dans ta copie]
Au Cameroun, des entreprises comme \textbf{Namé Recycling} à Douala collectent les déchets plastiques et les transforment en pavés et en objets utiles. Citer cet exemple dans ta copie enrichit ton contenu et montre que tu relies le chinois à la réalité de ton pays. Tu peux écrire par exemple : \zh{在杜阿拉，有公司把塑料变成砖。} (\emph{Zài Dù'ālā, yǒu gōngsī bǎ sùliào biànchéng zhuān.}) « À Douala, il y a des entreprises qui transforment le plastique en pavés. » En Chine aussi, le tri des déchets (\zh{垃圾分类}, \emph{lājī fēnlèi}) est devenu obligatoire dans les grandes villes comme Shanghai : les habitants trient en quatre catégories. Deux pays, un même geste !
\end{savaistu}

\begin{aretenir}[L'essentiel pour le jour de l'examen]
\begin{itemize}
\item \textbf{Plan d'abord, texte ensuite} : écris ton squelette en connecteurs (\zh{首先 → 然后 → 最后 → 因为…所以…}) avant de rédiger.
\item \textbf{Au Bac, la lisibilité des caractères et l'usage d'au moins 2 connecteurs rapportent des points} : soigne chaque trait et fais « briller » tes connecteurs.
\item \textbf{Modaux avant le verbe} : \zh{应该 + verbe}, \zh{可以 + verbe}, jamais après.
\item \textbf{Un caractère mal écrit est faux}, même si le sens se devine : en cas de doute, utilise un mot plus simple que tu maîtrises.
\item \textbf{Termine toujours par une conclusion forte}, par exemple : \zh{保护环境就是保护我们自己。}
\end{itemize}
\end{aretenir}

\newpage

\coursec{Activité d'intégration}

\coursec{Leçon 11 — Expression écrite : le recyclage}

\courssub{Modèle de texte commenté}

\begin{textebox}[Message du club environnement — Yaoundé, Semaine verte]
同学们好！下个星期是我们学校的“绿色周”。\\
Tóngxuémen hǎo! Xià ge xīngqī shì wǒmen xuéxiào de “lǜsè zhōu”.\\
« Chers camarades, la semaine prochaine, c'est la “Semaine verte” de notre école. »\\[4pt]
我们应该保护学校的环境：请大家把垃圾分类，把瓶子放进蓝色的箱子，把纸放进黄色的箱子。\\
Wǒmen yīnggāi bǎohù xuéxiào de huánjìng: qǐng dàjiā bǎ lājī fēnlèi, bǎ píngzi fàng jìn lánsè de xiāngzi, bǎ zhǐ fàng jìn huángsè de xiāngzi.\\
« Nous devons protéger l'environnement de l'école : merci à tous de trier les déchets, de mettre les bouteilles dans la caisse bleue et le papier dans la caisse jaune. »\\[4pt]
每个班可以做一张海报。最好的海报会贴在食堂旁边！\\
Měi ge bān kěyǐ zuò yì zhāng hǎibào. Zuì hǎo de hǎibào huì tiē zài shítáng pángbiān!\\
« Chaque classe peut faire une affiche. Les meilleures affiches seront collées à côté de la cantine ! »
\end{textebox}

\begin{aretenir}[Analyse du texte modèle]
\begin{itemize}
\item \textbf{Contexte :} Communication officielle (président du club $\to$ élèves), registre standard, ton mobilisateur.
\item \textbf{Structures repérées :}
      \begin{itemize}
      \item \zh{应该} + V (obligation morale collective) : \zh{我们应该保护\ldots}
      \item \zh{把} + Objet + V + Complément (trois occurrences) : \zh{把垃圾分类}、\zh{把瓶子放进\ldots}、\zh{把纸放进\ldots}
      \item \zh{可以} + V (possibilité / autorisation) : \zh{每个班可以做\ldots}
      \item Passif par \zh{会} + V (futur passif) : \zh{最好的海报会贴在\ldots}
      \end{itemize}
\item \textbf{Connecteurs :} \zh{：} (deux-points explicatifs), pas de \zh{首先/然后/最后} ici (texte court), mais structure \zh{原因 : 我们应该\ldots ; Consignes : 请大家\ldots}.
\item \textbf{Vocabulaire clé :} \zh{绿色周} (semaine verte), \zh{环境}, \zh{垃圾}, \zh{分类}, \zh{瓶子}, \zh{纸}, \zh{箱子}, \zh{海报}, \zh{食堂}.
\end{itemize}
\end{aretenir}

\courssub{Structure et connecteurs du texte}

\begin{propriete}[Verbes modaux : \zh{应该} vs \zh{可以}]
\begin{center}
\begin{tabularx}{\linewidth}{|X|X|X|}
\hline
\rowcolor{popblueL} Structure & Exemple (caractères + pinyin) & Traduction \\
\hline
Sujet + \cle{应该} + Verbe & \zh{我们应该保护环境。}\newline Wǒmen yīnggāi bǎohù huánjìng. & Nous devons / Il faut que nous protégions l'environnement (devoir moral, conseil). \\
\hline
Sujet + \cle{可以} + Verbe & \zh{我们可以回收瓶子。}\newline Wǒmen kěyǐ huíshōu píngzi. & Nous pouvons / Il est possible de recycler les bouteilles (capacité, possibilité, autorisation). \\
\hline
\end{tabularx}
\end{center}
\emph{Attention :} \zh{应该} exprime une nécessité subjective (ce qui est bien/juste) ; \zh{可以} exprime une faisabilité ou une permission. Ne les confondez pas.
\end{propriete}

\begin{propriete}[Construction en \zh{把} (disposition de l'objet)]
\emph{Schéma :} Sujet + \cle{把} + Objet + Verbe + \textbf{Complément obligatoire} (résultat, direction, lieu, aspect).
\begin{center}
\begin{tabularx}{\linewidth}{|X|X|X|}
\hline
\rowcolor{popblueL} Structure & Exemple & Traduction \\
\hline
S + 把 + O + V + 方向/结果 & \zh{请把瓶子放进蓝色的箱子。}\newline Qǐng bǎ píngzi fàng jìn lánsè de xiāngzi. & Mets les bouteilles dans la caisse bleue. \\
\hline
S + 把 + O + V + 完成/状态 & \zh{请把垃圾分类。}\newline Qǐng bǎ lājī fēnlèi. & Trie les déchets. (\zh{分类} verbe-résultat) \\
\hline
\end{tabularx}
\end{center}
\end{propriete}

\begin{attention}[Pièges \zh{把}]
\begin{itemize}
\item \textbf{Verbe jamais seul :} \zh{*把瓶子放} (incorrect) $\to$ \zh{把瓶子放进去/放好/放在箱子里}.
\item \textbf{Objet défini/connu :} \zh{把一本书看完} (peu naturel) $\to$ \zh{把这本书看完} (objet spécifique).
\item \textbf{Négation :} \zh{不} ou \zh{别} devant \zh{把} : \zh{别把纸扔了。} (Ne jette pas le papier).
\item \textbf{Aspect \zh{了} :} Dans une consigne/impératif, pas de \zh{了} final. \zh{请把垃圾分类。} (Pas \zh{分类了}).
\end{itemize}
\end{attention}


\begin{propriete}[Connecteurs d'enchaînement et corrélation cause-effet]
\begin{center}
\begin{tabularx}{\linewidth}{|X|X|X|}
\hline
\rowcolor{popblueL} Fonction & Connecteur chinois & Exemple court \\
\hline
Étape 1 (initiale) & \cle{首先} (shǒuxiān) & \zh{首先，把垃圾分类。} \\
\hline
Étape 2 (intermédiaire) & \cle{然后} (ránhòu) & \zh{然后，把瓶子放进蓝箱。} \\
\hline
Étape 3 (finale) & \cle{最后} (zuìhòu) & \zh{最后，我们可以回收纸。} \\
\hline
Cause $\to$ Conséquence & \cle{因为} \ldots \cle{所以} \ldots (yīnwèi \ldots suǒyǐ \ldots) & \zh{因为环境重要，所以我们行动。} \\
\hline
\end{tabularx}
\end{center}
\emph{Note :} En chinois, \zh{因为} et \zh{soi} coexistent dans la même phrase (correct), contrairement au français (*\emph{Parce que... donc...}).
\end{propriete}

\begin{center}
\begin{tabularx}{\linewidth}{|X|X|X|X|}
\hline
\rowcolor{popblueL} Caractères & Pinyin & Nature & Traduction \\
\hline
回收 & huíshōu & verbe & recycler, récupérer \\
\hline
垃圾 & lājī & nom & déchets, ordures \\
\hline
分类 & fēnlèi & verbe & trier, classer \\
\hline
瓶子 & píngzi & nom & bouteille \\
\hline
塑料 & sùliào & nom & plastique \\
\hline
纸 & zhǐ & nom & papier \\
\hline
环境 & huánjìng & nom & environnement \\
\hline
保护 & bǎohù & verbe & protéger \\
\hline
校园 & xiàoyuán & nom & campus, cour de l'école \\
\hline
海报 & hǎibào & nom & affiche, poster \\
\hline
干净 & gānjìng & adjectif & propre \\
\hline
箱子 & xiāngzi & nom & caisse, bac, boîte \\
\hline
蓝色 & lánsè & adj/nom & bleu \\
\hline
黄色 & huángsè & adj/nom & jaune \\
\hline
旧 & jiù & adjectif & vieux, usagé \\
\hline
报纸 & bàozhǐ & nom & journal \\
\hline
回收站 & huíshōuzhàn & nom & centre de recyclage \\
\hline
健康 & jiànkāng & adjectif & en bonne santé \\
\hline
行动 & xíngdòng & verbe/nom & agir, action \\
\hline
美化 & měihuà & verbe & embellir \\
\hline
\end{tabularx}
\end{center}

\courssub{Écriture correcte des caractères du thème}

\begin{methode}[Ordre des traits et radicaux — Clés de mémorisation]
Écris chaque caractère 5 fois en respectant l'ordre : \textbf{haut $\to$ bas, gauche $\to$ droite, horizontal $\to$ vertical, trait central $\to$ côtés, fermer en dernier}.
\begin{enumerate}
\item \zh{垃} (lā) : 7 traits. Clé \zh{土} (tǔ, terre) à gauche $\to$ déchets = matière terreuse. \emph{Ordre :} 一 (horiz.), 丨 (vert.), 一 (horiz. base) pour \zh{土} ; puis \zh{立} (lì) à droite (丨, 一, 一, 丨, 一).
\item \zh{圾} (jī) : 11 traits. Clé \zh{土} à gauche. Partie droite \zh{及} (jí) : 丨, 丿, 丶, 丨, (trait brisé), 一, 一. \textbf{Piège :} ne pas confondre avec \zh{己} (jǐ, soi-même) qui a 3 traits (丨, (trait brisé), 丿).
\item \zh{瓶} (píng) : 10 traits. Clé \zh{瓦} (wǎ, tuile) à \textbf{droite} (4 traits : (trait brisé), 丿, 丨, 一). Partie gauche \zh{并} (bìng) : 一, 丨, 一, 丨, 一, 丶. \emph{Astuce :} une bouteille en tuile/terre cuite.
\item \zh{境} (jìng) : 14 traits. Clé \zh{土} (relevée \zh{阝}) à gauche. Partie droite \zh{镜} (jìng, miroir) sans le trait vertical final du \zh{金} ? Non, \zh{境} = \zh{阝} + \zh{镜} (金 + 立 + 日). \textbf{Distinction :} \zh{镜} (miroir, clé \zh{金}) vs \zh{境} (frontière/environnement, clé \zh{土}). Un trait de plus (\zh{土} vs \zh{金}) change le sens.
\item \zh{回} (huí) : 6 traits. \zh{囗} (box) + \zh{口} (bouche) dedans. Ordre : 丨, (trait brisé), 一 (fermeture haut), 丨, (trait brisé), 一 (intérieur).
\item \zh{收} (shōu) : 6 traits. Clé \zh{攵} (pū, main tenant bâton) à droite (4 traits : 丿, 丶, 丿, 丶). Gauche \zh{丩} (jiū) : 丨, (trait brisé), 丿.
\item \zh{分} (fēn) : 4 traits. 丿, 丶, 丿, 丶 (huit divisé en deux).
\item \zh{类} (lèi) : 18 traits. Clé \zh{米} (mǐ, riz) en haut. Bas \zh{大} + \zh{页} (yè, tête/page). \emph{Sens :} grains de riz classés par grosses catégories.
\end{enumerate}
\end{methode}

\begin{experience}[Atelier écriture : « Caractères au tableau »]
En binôme : l'un dicte le pinyin (ex. \emph{lājī}), l'autre écrit au tableau en annonçant les traits à voix haute (\emph{« clé terre, trait horizontal, trait vertical... »}). Inversez les rôles pour les 8 caractères ci-dessus. Vérification par l'enseignant (ordre, proportion, clé).
\end{experience}

\courssub{Thème : français $\to$ chinois}

\begin{exercice}[Thème n°1 — Phrases isolées]
Traduis en chinois (caractères + pinyin) :
\begin{enumerate}
\item Nous devons protéger l'environnement de l'école.
\item S'il te plaît, trie les déchets.
\item Mets les bouteilles en plastique dans la caisse bleue.
\item Nous pouvons vendre les vieux journaux au centre de recyclage.
\item Parce que le campus est propre, nous sommes en bonne santé.
\end{enumerate}
\end{exercice}

\begin{corrige}[Thème n°1]
\begin{enumerate}
\item \zh{我们应该保护学校的环境。} (Wǒmen yīnggāi bǎohù xuéxiào de huánjìng.)
\item \zh{请把垃圾分类。} (Qǐng bǎ lājī fēnlèi.) \emph{Note : impératif poli, pas de 了.}
\item \zh{请把塑料瓶子放进蓝色的箱子。} (Qǐng bǎ sùliào píngzi fàng jìn lánsè de xiāngzi.) \emph{Classificateur implicite ou \zh{个} si nombre.}
\item \zh{我们可以把旧报纸卖给回收站。} (Wǒmen kěyǐ bǎ jiù bàozhǐ mài gěi huíshōuzhàn.)
\item \zh{因为校园干净，所以我们很健康。} (Yīnwèi xiàoyuán gānjìng, suǒyǐ wǒmen hěn jiànkāng.) \emph{Ou : \zh{因为校园干净，所以我们更健康。}}
\end{enumerate}
\end{corrige}

\begin{exercice}[Thème n°2 — Mini-paragraphe (niveau examen)]
Traduis le texte suivant en chinois (environ 60-80 caractères) :
« Demain, c'est la Semaine verte. Notre classe doit faire une affiche. D'abord, nous devons écrire le titre. Ensuite, nous devons dessiner des poubelles bleues et jaunes. Enfin, nous collons l'affiche près de la cantine. Parce que nous aimons notre école, nous agissons ensemble. »
\end{exercice}

\begin{corrige}[Thème n°2]
\zh{明天是绿色周。我们班应该做一张海报。首先，写标题。然后，画蓝色和黄色的箱子。最后，把海报贴在食堂旁边。因为我们爱学校，所以一起行动吧！}\\
Míngtiān shì lǜsè zhōu. Wǒmen bān yīnggāi zuò yì zhāng hǎibào. Shǒuxiān, xiě biāotí. Ránhòu, huà lánsè hé huángsè de xiāngzi. Zuìhòu, bǎ hǎibào tiē zài shítáng pángbiān. Yīnwèi wǒmen ài xuéxiào, suǒyǐ yìqǐ xíngdòng ba!
\end{corrige}

\courssub{Version : chinois $\to$ français}

\begin{exercice}[Version n°1 — Phrases]
Traduis en français :
\begin{enumerate}
\item \zh{同学们，请把纸和塑料分开。}
\item \zh{我们可以把旧瓶子卖给回收站。}
\item \zh{因为保护环境很重要，所以我们应该行动。}
\item \zh{首先分类，然后放进箱子，最后校园干净。}
\item \zh{这张海报很漂亮，会贴在教室里。}
\end{enumerate}
\end{exercice}

\begin{corrige}[Version n°1]
\begin{enumerate}
\item Camarades, merci de séparer le papier et le plastique. (Ou : triez le papier et le plastique séparément.)
\item Nous pouvons vendre les vieilles bouteilles au centre de recyclage.
\item Parce que protéger l'environnement est important, nous devons agir.
\item D'abord trier, ensuite mettre dans la caisse, enfin le campus est propre. (Style télégraphique d'affiche/slogan).
\item Cette affiche est très jolie, elle sera collée dans la salle de classe.
\end{enumerate}
\end{corrige}

\begin{exercice}[Version n°2 — Texte suivi]
Traduis en français le texte suivant :
\zh{我们学校下周有绿色周。校长说：“同学们，我们应该爱护环境。请大家把垃圾分类。把瓶子放进蓝箱子，把纸放进黄箱子。因为环境干净，所以我们健康。每个班可以做海报。最好的海报贴在食堂旁边。”}
\end{exercice}

\begin{corrige}[Version n°2]
« Notre école organise la Semaine verte la semaine prochaine. Le proviseur dit : "Camarades, nous devons chérir l'environnement. Merci à tous de trier les déchets. Mettez les bouteilles dans la caisse bleue, le papier dans la caisse jaune. Parce que l'environnement est propre, nous sommes en bonne santé. Chaque classe peut faire une affiche. La meilleure affiche sera collée à côté de la cantine." »
\end{corrige}

\courssub{Production guidée et grille d'évaluation}

\courssub{Objectifs de l'activité}

À l'issue de cette activité d'intégration, tu seras capable de :
\begin{itemize}
\item produire en binôme une \cle{affiche} en chinois de 6 à 8 phrases incitant au tri et au recyclage dans ton lycée ;
\item mobiliser correctement le \cle{vocabulaire du thème} (回收, 垃圾, 分类, 瓶子, 纸, 环境…) et écrire les caractères au propre, dans le bon ordre des traits ;
\item employer les structures modales \zh{应该} et \zh{可以}, la construction en \zh{把}, ainsi que les connecteurs \zh{首先}, \zh{然后}, \zh{最后} et la corrélation \zh{因为……所以……} ;
\item présenter oralement ton affiche en une minute et évaluer celles des autres groupes à l'aide de la grille d'évaluation de la leçon.
\end{itemize}

\courssub{Situation}

Le club environnement de ton lycée à Yaoundé organise la « Semaine verte ». La direction a constaté que les élèves jettent tout — bouteilles en plastique, papiers, restes de nourriture — dans la même poubelle de la cour. Le club décide donc d'afficher des messages en chinois, la langue que tu apprends, pour montrer que les élèves de Terminale A maîtrisent le sujet et pour sensibiliser toute l'école. Chaque binôme de la classe doit créer une affiche. Les meilleures seront collées près des poubelles et de la cantine.

Voici le message du président du club, affiché au tableau (rappel) :

\begin{textebox}[Message du club environnement]
同学们好！下个星期是我们学校的“绿色周”。\\
Tóngxuémen hǎo! Xià ge xīngqī shì wǒmen xuéxiào de “lǜsè zhōu”.\\
« Chers camarades, la semaine prochaine, c'est la “Semaine verte” de notre école. »\\[4pt]
我们应该保护学校的环境：请大家把垃圾分类，把瓶子放进蓝色的箱子，把纸放进黄色的箱子。\\
Wǒmen yīnggāi bǎohù xuéxiào de huánjìng: qǐng dàjiā bǎ lājī fēnlèi, bǎ píngzi fàng jìn lánsè de xiāngzi, bǎ zhǐ fàng jìn huángsè de xiāngzi.\\
« Nous devons protéger l'environnement de l'école : merci de trier les déchets, de mettre les bouteilles dans la caisse bleue et le papier dans la caisse jaune. »\\[4pt]
每个班可以做一张海报。最好的海报会贴在食堂旁边！\\
Měi ge bān kěyǐ zuò yì zhāng hǎibào. Zuì hǎo de hǎibào huì tiē zài shítáng pángbiān!\\
« Chaque classe peut faire une affiche. Les meilleures affiches seront collées à côté de la cantine ! »
\end{textebox}

\courssub{Consignes}

Travail en binôme, en trois étapes guidées (environ 45 minutes), puis présentation orale.

\begin{enumerate}
\item \textbf{Comprendre la commande.} Lis le message du club ci-dessus. Relève : (a) les deux verbes modaux utilisés ; (b) les trois phrases construites avec \zh{把} ; (c) le mot qui signifie « trier ». Note-les sur ton brouillon avec le pinyin.
\item \textbf{Préparer le lexique.} À l'aide du tableau de vocabulaire de la leçon, choisis 8 à 10 mots utiles pour ton affiche (回收, 垃圾, 分类, 瓶子, 纸, 塑料, 环境, 保护, 干净, 校园…). Écris chaque caractère trois fois au brouillon en respectant l'ordre des traits ; vérifie en particulier les caractères \zh{圾} (clé de la terre 土 à gauche), \zh{瓶} (clé de la tuile 瓦 à droite) et \zh{境} (ne pas confondre avec \zh{镜}, « miroir »).
\item \textbf{Rédiger l'affiche (6 à 8 phrases).} Ton affiche doit contenir obligatoirement :
\begin{itemize}
\item un \cle{titre accrocheur} (4 à 8 caractères) ;
\item trois consignes numérotées avec \zh{首先}、\zh{然后}、\zh{最后}, dont au moins deux phrases en \zh{把} et une avec \zh{应该} ou \zh{可以} ;
\item un \cle{slogan de conclusion} avec la corrélation \zh{因为……所以……}.
\end{itemize}
\item \textbf{Mettre au propre et illustrer.} Recopie l'affiche sur une feuille A4 : caractères grands, droits et lisibles (de haut en bas, de gauche à droite, l'horizontal avant le vertical). Ajoute un dessin simple (poubelles de couleurs, bouteille, feuille de papier).
\item \textbf{Présenter et voter.} Présente ton affiche à la classe en 1 minute : lis le titre, les trois consignes et le slogan, avec des tons soignés. La classe vote ensuite pour la meilleure affiche selon la grille d'évaluation (exactitude des caractères, structures exigées, connecteurs, clarté de la présentation). Les affiches gagnantes seront exposées dans la salle.
\end{enumerate}

\courssub{Ressources à mobiliser}

\begin{aretenir}[Structures et connecteurs indispensables — Synthèse]
\begin{itemize}
\item \cle{应该} + verbe : « devoir » (conseil moral). \zh{我们应该保护环境。} \emph{Wǒmen yīnggāi bǎohù huánjìng.} « Nous devons protéger l'environnement. »
\item \cle{可以} + verbe : « pouvoir / il est possible de ». \zh{我们可以回收瓶子。} \emph{Wǒmen kěyǐ huíshōu píngzi.} « Nous pouvons recycler les bouteilles. »
\item Construction en \cle{把} : Sujet + 把 + objet + verbe + \textbf{complément}. \zh{请把垃圾放进箱子。} \emph{Qǐng bǎ lājī fàng jìn xiāngzi.} « Mets les déchets dans la caisse, s'il te plaît. »
\item Connecteurs d'étapes : \zh{首先} (d'abord), \zh{然后} (ensuite), \zh{最后} (enfin).
\item Corrélation de cause : \zh{因为……所以……}. \zh{因为地球只有一个，所以我们要保护它。} \emph{Yīnwèi dìqiú zhǐyǒu yí ge, suǒyǐ wǒmen yào bǎohù tā.} « Parce qu'il n'y a qu'une seule Terre, nous devons la protéger. »
\end{itemize}
\end{aretenir}

\begin{attention}[Erreurs fréquentes à éviter — Check-list avant remise]
\begin{itemize}
\item \textbf{Pas de \zh{了} en impératif \zh{把}} : \zh{请把垃圾分类。} (Pas \zh{分类了}).
\item \textbf{Complément obligatoire après verbe \zh{把}} : \zh{把瓶子放进箱子里} (pas \zh{把瓶子放}).
\item \textbf{\zh{因为……所以……} ensemble} : correct en chinois, faux en français (*\emph{Parce que... donc...}).
\item \textbf{Caractères pièges} : \zh{已} (yǐ, déjà) $\neq$ \zh{己} (jǐ, soi) ; \zh{境} (env., clé 土) $\neq$ \zh{镜} (miroir, clé 金) ; \zh{圾} (clé 土) $\neq$ \zh{及} (partie droite seule).
\item \textbf{Classificateur \zh{张} pour \zh{海报}} : \zh{一张海报} (pas \zh{个}).
\end{itemize}
\end{attention}

\courssub{Proposition de production et corrigé commenté}

\begin{exoresolu}[Modèle d'affiche « Recyclons à notre lycée ! »]
Énoncé : rédige une affiche complète répondant à toutes les contraintes (titre, trois consignes avec 首先/然后/最后, structures 应该/可以/把, slogan avec 因为……所以……).
\tcblower
\textbf{Affiche modèle :}\\[4pt]
\zh{回收垃圾，美化校园！}\\
\emph{Huíshōu lājī, měihuà xiàoyuán!} « Recyclons les déchets, embellissons notre campus ! »\\[4pt]
\zh{同学们，我们的学校应该更干净、更美丽。}\\
\emph{Tóngxuémen, wǒmen de xuéxiào yīnggāi gèng gānjìng, gèng měilì.} « Camarades, notre école doit être plus propre et plus belle. »\\[4pt]
\zh{首先，请把垃圾分类：瓶子、纸和塑料不要放在一起。}\\
\emph{Shǒuxiān, qǐng bǎ lājī fēnlèi: píngzi, zhǐ hé sùliào bú yào fàng zài yìqǐ.} « D'abord, trie les déchets : ne mets pas les bouteilles, le papier et le plastique ensemble. »\\[4pt]
\zh{然后，请把瓶子放进蓝色的箱子，把纸放进黄色的箱子。}\\
\emph{Ránhòu, qǐng bǎ píngzi fàng jìn lánsè de xiāngzi, bǎ zhǐ fàng jìn huángsè de xiāngzi.} « Ensuite, mets les bouteilles dans la caisse bleue et le papier dans la caisse jaune. »\\[4pt]
\zh{最后，我们可以把旧报纸卖给回收站。}\\
\emph{Zuìhòu, wǒmen kěyǐ bǎ jiù bàozhǐ mài gěi huíshōu zhàn.} « Enfin, nous pouvons vendre les vieux journaux au centre de recyclage. »\\[4pt]
\zh{因为干净的校园让我们更健康，所以我们一起行动吧！}\\
\emph{Yīnwèi gānjìng de xiàoyuán ràng wǒmen gèng jiànkāng, suǒyǐ wǒmen yìqǐ xíngdòng ba!} « Parce qu'un campus propre nous rend en meilleure santé, agissons ensemble ! »\\[6pt]
\textbf{Commentaire des choix de langue :}
\begin{itemize}
\item Le titre en deux parties de quatre caractères (\zh{回收垃圾，美化校园}) imite le rythme des slogans chinois : court, symétrique, facile à mémoriser.
\item \zh{应该} introduit l'objectif général (un devoir collectif) ; \zh{可以} introduit une possibilité concrète (vendre les vieux journaux) : les deux modaux sont bien distingués.
\item Trois phrases en \zh{把} montrent la maîtrise de la structure : objet déplacé avant le verbe, complément de direction \zh{放进……} obligatoire.
\item Les connecteurs \zh{首先}、\zh{然后}、\zh{最后} organisent clairement les trois consignes, comme exigé.
\item Le slogan final utilise \zh{因为……所以……} avec la particule d'exhortation \zh{吧}, qui invite à l'action : effet de style typique d'une affiche.
\item Vocabulaire du thème entièrement mobilisé : 回收, 垃圾, 分类, 瓶子, 塑料, 校园, 干净, 环境 (implicite).
\end{itemize}
\end{exoresolu}

\courssub{Grille d'évaluation (sur 20 points)}

\begin{center}
\begin{tabularx}{\linewidth}{|X|c|X|}
\hline
\rowcolor{popblueL} \textbf{Critère} & \textbf{Points} & \textbf{Indicateurs de réussite} \\
\hline
Exactitude des caractères (écriture, traits, clés) & 4 & 0-1 erreur mineure ; lisibilité parfaite ; proportions respectées. \\
\hline
Structures grammaticales exigées (\zh{应该}, \zh{可以}, \zh{把} $\times$2 min.) & 5 & Toutes présentes, correctement formées (complément après \zh{把}, pas de \zh{了} intempestif). \\
\hline
Connecteurs (\zh{首先}, \zh{然后}, \zh{最后}, \zh{因为……所以……}) & 4 & Les 4 utilisés à bon escient, ponctuation chinoise correcte. \\
\hline
Richesse et justesse du vocabulaire (min. 8 mots du thème) & 3 & Mots variés, bien employés, pas de répétition excessive. \\
\hline
Cohérence du discours (titre, consignes, slogan) & 2 & Logique claire, ton adapté (affiche), slogan percutant. \\
\hline
Présentation orale (1 min, tons, fluidité, regard) & 2 & Lecture correcte, tons distincts, débit naturel, contact visuel. \\
\hline
\multicolumn{2}{|c|}{\textbf{Total}} & \textbf{/20} \\
\hline
\end{tabularx}
\end{center}

\courssub{Bilan de l'activité}

Coche ce que tu sais désormais faire :
\begin{itemize}
\item[$\square$] Je sais rédiger une affiche structurée en chinois : titre, consignes ordonnées, slogan.
\item[$\square$] J'utilise correctement \zh{应该} et \zh{可以} sans les confondre.
\item[$\square$] Je construis une phrase en \zh{把} avec un complément après le verbe.
\item[$\square$] J'emploie \zh{首先}、\zh{然后}、\zh{最后} et \zh{因为……所以……} pour lier mes phrases.
\item[$\square$] J'écris les caractères du thème (垃圾, 瓶子, 环境, 回收) sans erreur de traits ni de clés.
\item[$\square$] Je peux présenter mon affiche à l'oral en une minute, avec des tons corrects.
\item[$\square$] Je sais traduire un court texte du français vers le chinois (thème) et du chinois vers le français (version) sur ce thème.
\end{itemize}

\begin{savaistu}[Le tri en Chine : du obligatoire au high-tech]
Depuis 2019, Shanghai a rendu le tri obligatoire (\zh{强制分类}) avec quatre catégories codées par couleur : \zh{可回收物} (bleu, recyclables), \zh{有害垃圾} (rouge, dangereux : piles, médicaments), \zh{厨余垃圾} (vert, déchets de cuisine), \zh{其他垃圾} (gris, reste). Des caméras \zh{智能垃圾桶} (poubelles intelligentes) à reconnaissance faciale donnent des points \zh{积分} échangeables contre des biens. Au Cameroun, des start-ups comme \emph{Plastic Waste Recycling} à Douala transforment le plastique en pavés écologiques : un défi commun, des solutions locales !
\end{savaistu}

\begin{popfigure}
\begin{tikzpicture}[node distance=1.2cm, every node/.style={font=\small, align=center}]
\node[draw, rounded corners, fill=popgreenL, thick] (start) {Déchet jeté};
\node[draw, diamond, aspect=2, fill=poporangeL, thick, below of=start] (tri) {Tri à la source ?};
\node[draw, rounded corners, fill=popblueL, thick, left of=tri, xshift=-2.5cm, align=center] (oui){OUI : \zh{分类}\\(回收/厨余/有害/其他)};
\node[draw, rounded corners, fill=popredL, thick, right of=tri, xshift=2.5cm, align=center] (non){NON : \zh{混合垃圾}\\Enfouissement / Incinération};
\node[draw, rounded corners, fill=popgreenL, thick, below of=oui, yshift=-0.5cm, align=center] (recyclage){Recyclage / Valorisation\\(Compost, Énergie, Matières)};
\node[draw, rounded corners, fill=poppurpleL, thick, below of=recyclage, align=center] (fin){Environnement protégé\\(Cameroun \& Chine)};
\draw[->, thick, popgreen] (start) -- (tri);
\draw[->, thick, popgreen] (tri) -- node[above, sloped] {\zh{应该}} (oui);
\draw[->, thick, popred] (tri) -- node[above, sloped] {Non} (non);
\draw[->, thick, popgreen] (oui) -- (recyclage);
\draw[->, thick, popgreen] (recyclage) -- (fin);
\draw[->, dashed, thick, popmuted] (non) |- node[near start, right] {Pollution} (fin);
\end{tikzpicture}
\legende{Schéma du cycle de gestion des déchets : l'importance du tri (\zh{分类}) à la source.}
\end{popfigure}

\newpage

\coursec{Méthodes et exercices résolus}

\coursec{Leçon 11 — Expression écrite : le recyclage}

\courssub{Texte support \& Vocabulaire du thème}

\begin{textebox}[Modèle de paragraphe argumentatif — Le recyclage à l'école]
\zh{在雅温得，垃圾越来越多，污染很严重。因为塑料瓶和纸到处都是，所以我们学生应该行动。首先，我们在学校回收纸和塑料瓶。然后，我们把旧东西再利用，不乱扔。这样，我们就能保护环境，让城市更干净、更美丽。}
\\
\emph{Zài Yǎwēndé, lājī yuèláiyuè duō, wūrǎn hěn yánzhòng. Yīnwèi sùliàopíng hé zhǐ dàochù dōu shì, suǒyǐ wǒmen xuésheng yīnggāi xíngdòng. Shǒuxiān, wǒmen zài xuéxiào huíshōu zhǐ hé sùliàopíng. Ránhòu, wǒmen bǎ jiù dōngxi zài lìyòng, bù luàn rēng. Zhèyàng, wǒmen jiù néng bǎohù huánjìng, ràng chéngshì gèng gānjìng, gèng měilì.}
\\
« À Yaoundé, les déchets sont de plus en plus nombreux, la pollution est grave. Parce que les bouteilles en plastique et le papier sont partout, nous, les élèves, devons agir. D'abord, nous recyclons le papier et les bouteilles en plastique à l'école. Ensuite, nous réutilisons les vieux objets, nous ne les jetons pas n'importe où. Ainsi, nous pouvons protéger l'environnement et rendre la ville plus propre et plus belle. »
\end{textebox}

\begin{center}
\begin{tabularx}{\linewidth}{|X|X|X|X|}
\hline
\rowcolor{popblueL} Caractères & Pinyin & Nature & Traduction \\
\hline
垃圾 & lājī & n. & déchets, ordures \\
\hline
回收 & huíshōu & v. & recycler, récupérer \\
\hline
塑料瓶 & sùliàopíng & n. & bouteille en plastique \\
\hline
纸 & zhǐ & n. & papier \\
\hline
环境 & huánjìng & n. & environnement \\
\hline
保护 & bǎohù & v. & protéger \\
\hline
污染 & wūrǎn & v./n. & polluer / pollution \\
\hline
再利用 & zài lìyòng & v. & réutiliser, donner une seconde vie \\
\hline
越来越 & yuèláiyuè & adv. & de plus en plus \\
\hline
严重 & yánzhòng & adj. & grave, sérieux \\
\hline
到处 & dàochù & pr. & partout \\
\hline
行动 & xíngdòng & v. & agir, passer à l'action \\
\hline
首先 / 然后 & shǒuxiān / ránhòu & adv. & d'abord / ensuite \\
\hline
乱扔 & luàn rēng & v. & jeter n'importe comment \\
\hline
干净 & gānjìng & adj. & propre \\
\hline
美丽 & měilì & adj. & beau, joli \\
\hline
\end{tabularx}
\end{center}

\courssub{Fiche méthode 1 — Produire un texte simple sur le recyclage}

\begin{methode}[Rédiger un texte simple sur le recyclage en 5 étapes]
\begin{enumerate}
\item \cle{Analyser le sujet} : souligner le thème (le recyclage \zh{回收} \emph{huíshōu}), le type de texte demandé (paragraphe, message, affiche) et le nombre de caractères exigé (au Bac, souvent 60 à 100 caractères chinois, ponctuation non comprise).
\item \cle{Lister le vocabulaire utile avant d'écrire} : \zh{垃圾 lājī} (déchets), \zh{塑料瓶 sùliàopíng} (bouteille en plastique), \zh{纸 zhǐ} (papier), \zh{环境 huánjìng} (environnement), \zh{保护 bǎohù} (protéger), \zh{污染 wūrǎn} (pollution), \zh{再利用 zài lìyòng} (réutiliser), \zh{分类 fēnlèi} (trier/séparer), \zh{垃圾桶 lājītǒng} (poubelle).
\item \cle{Construire le plan en 3 temps} : (a) \textbf{Constat} — il y a beaucoup de déchets / la pollution est grave ; (b) \textbf{Action} — on peut trier, recycler, réutiliser ; (c) \textbf{Conclusion} — cela protège l'environnement / embellit la ville. Une idée = une phrase courte et maîtrisée.
\item \cle{Rédiger avec des structures maîtrisées (HSK 3-4)} :
\begin{itemize}
\item \zh{Subject + 可以/能/应该 + Verbe} (capacité/obligation) : \zh{我们应该回收。}
\item \zh{把 + Objet défini + Verbe + 结果/方向} (action sur objet) : \zh{把塑料瓶放进垃圾桶。}
\item \zh{越来越 + Adjectif} (évolution) : \zh{垃圾越来越少。}
\item \zh{因为……所以……} (cause-conséquence) : \zh{因为污染严重，所以我们要行动。}
\item \zh{首先……然后……最后……} (chronologie) : \zh{首先分类，然后回收。}
\end{itemize}
\item \cle{Relire en vérifiant 4 points critiques} : (1) Ordre des mots : Complément de lieu/temps \textbf{avant} le verbe ; (2) Classificateur obligatoire devant nom compté ; (3) Tons corrects sur le pinyin ; (4) Caractères complets (traits, radicaux, pas de confusion \zh{在}/\zh{再}, \zh{境}/\zh{镜}).
\end{enumerate}
\textbf{Réflexe d'examen :} N'écris que des phrases que tu sais écrire parfaitement. Une phrase simple correcte vaut plus qu'une phrase complexe fautive. Vise 3 à 5 phrases bien liées.
\end{methode}

\begin{exoresolu}[Sujet type Bac : paragraphe sur le recyclage]
\textbf{Énoncé.} Rédigez en chinois un texte de 60 à 80 caractères sur le sujet : « À Yaoundé, on produit beaucoup de déchets. Expliquez ce que les élèves peuvent faire pour recycler et protéger l'environnement. »
\tcblower
\textbf{Étape 1 — Analyse.} Type : paragraphe argumentatif simple. Personnes : les élèves (\zh{我们} / \zh{学生们}). Contenu attendu : constat + actions + conclusion.

\textbf{Étape 2 — Vocabulaire préparé.}
\begin{center}
\begin{tabularx}{\linewidth}{|X|X|X|X|}
\hline
\rowcolor{popblueL} Caractères & Pinyin & Nature & Traduction \\
\hline
垃圾 & lājī & n. & déchets \\
\hline
回收 & huíshōu & v. & recycler \\
\hline
塑料瓶 & sùliàopíng & n. & bouteille en plastique \\
\hline
保护环境 & bǎohù huánjìng & loc. v. & protéger l'environnement \\
\hline
越来越 & yuèláiyuè & adv. & de plus en plus \\
\hline
分类 & fēnlèi & v. & trier (les déchets) \\
\hline
\end{tabularx}
\end{center}

\textbf{Étape 3 — Plan.} (a) À Yaoundé, de plus en plus de déchets ; (b) Les élèves trient et recyclent papier et bouteilles ; (c) Ainsi on protège l'environnement.

\textbf{Étape 4 — Rédaction (réponse modèle).}

\zh{在雅温得，垃圾越来越多。我们学生应该分类回收纸和塑料瓶，也可以把旧东西再利用。这样，我们就能保护我们的环境。}

\emph{Zài Yǎwēndé, lājī yuèláiyuè duō. Wǒmen xuésheng yīnggāi fēnlèi huíshōu zhǐ hé sùliàopíng, yě kěyǐ bǎ jiù dōngxi zài lìyòng. Zhèyàng, wǒmen jiù néng bǎohù wǒmen de huánjìng.}

« À Yaoundé, les déchets sont de plus en plus nombreux. Nous, les élèves, devons trier et recycler le papier et les bouteilles en plastique, et nous pouvons aussi réutiliser les vieux objets. Ainsi, nous pouvons protéger notre environnement. » (environ 62 caractères hors ponctuation : conforme).

\textbf{Étape 5 — Justifications grammaticales.}
\begin{itemize}
\item \zh{垃圾越来越多} : structure \cle{越来越 + adjectif} = « de plus en plus [adj] » ; \textbf{pas de 很} devant l'adjectif dans cette structure.
\item \zh{应该分类回收} : auxiliaire modal \cle{应该} placé \textbf{avant} la série verbale, sans \zh{的} entre les deux. \zh{分类} (trier) précise le mode de recyclage.
\item \zh{把旧东西再利用} : construction \cle{把 + objet + verbe}, correcte car le verbe \zh{再利用} agit sur un objet défini (\zh{旧东西}) et implique une manipulation/résultat.
\item \zh{这样，…就…} : connecteur de conséquence « ainsi… alors… » ; \zh{就} marque la réalisation de la conséquence.
\end{itemize}
\textbf{Conclusion :} texte court, trois idées reliées par \zh{也} et \zh{这样}, structures HSK 3 maîtrisées (\zh{把}, \zh{越来越}, \zh{因为/所以} implicite) : c'est la stratégie gagnante au Bac.
\end{exoresolu}

\courssub{Fiche méthode 2 — Écrire correctement les caractères du thème}

\begin{methode}[Maîtriser traits, radicaux et caractères proches]
\begin{enumerate}
\item \cle{Identifier le radical (clé)} d'abord : il donne l'indice sémantique et dicte l'ordre général. 
\begin{itemize}
\item \zh{垃、圾} : clé \zh{土} (terre) à gauche → lien avec la matière/sol.
\item \zh{环} : clé \zh{王} (jade/roi) à gauche (forme simplifiée de \zh{玉}) → anneau de jade.
\item \zh{纸} : clé \zh{纟} (soie, forme cursive de \zh{糸}) à gauche → matière fibreuse.
\item \zh{瓶} : clé \zh{瓦} (tuile/poterie) à droite → objet en terre cuite/verre.
\item \zh{境} : clé \zh{土} (terre) à gauche → lieu, frontière, environnement.
\end{itemize}
\item \cle{Respecter l'ordre des traits} (règles universelles) : haut → bas ; gauche → droite ; horizontal avant vertical ; extérieur avant intérieur ; on ferme la boîte en dernier.
\item \cle{Compter les traits} pour vérifier (dictionnaire) : \zh{回} = 6, \zh{收} = 6, \zh{环} = 8, \zh{境} = 14, \zh{纸} = 7, \zh{瓶} = 10, \zh{再} = 6, \zh{在} = 6.
\item \cle{Distinguer les caractères proches (paires minimales)} par un mot-repère visuel ou sémantique :
\begin{center}
\begin{tabularx}{\linewidth}{|X|X|X|X|X|}
\hline
\rowcolor{popblueL} Cible & Pinyin & Sens & Confus possible & Différence clé \\
\hline
\zh{圾} & jí & déchets (\zh{垃圾}) & \zh{级} jí (niveau) & \zh{圾} a \zh{土} (terre) ; \zh{级} a \zh{纟} (soie/niveau) \\
\hline
\zh{瓶} & píng & bouteille & \zh{饼} bǐng (galette) & \zh{瓶} a \zh{瓦} (poterie) ; \zh{饼} a \zh{食} (manger) \\
\hline
\zh{境} & jìng & environnement & \zh{镜} jìng (miroir) & \zh{境} a \zh{土} (lieu) ; \zh{镜} a \zh{金} (métal/verre) \\
\hline
\zh{再} & zài & de nouveau & \zh{在} zài (être à) & \zh{再} : \zh{冉} au milieu ; \zh{在} : \zh{土} au milieu \\
\hline
\end{tabularx}
\end{center}
\item \cle{S'entraîner par séries} : écrire chaque caractère 5 fois en prononçant le pinyin à voix haute, puis le réécrire de mémoire dans une phrase complète.
\end{enumerate}
\end{methode}

\begin{exoresolu}[Dictée de caractères et correction d'erreurs classiques]
\textbf{Énoncé.} (a) Écrivez en caractères : \emph{huíshōu} (recycler), \emph{huánjìng} (environnement), \emph{zài lìyòng} (réutiliser), \emph{fēnlèi} (trier). (b) Corrigez la phrase fautive : \zh{我们在学校应该在利用纸。}
\tcblower
\textbf{(a) Écriture guidée (ordre des traits clé).}
\begin{itemize}
\item \emph{huíshōu} → \zh{回收}. 
\zh{回} (6 traits) : 1. Vertical gauche (丨), 2. Horizontal haut + vertical droit ((trait brisé)), 3. Horizontal milieu (一), 4. Vertical milieu (丨), 5. Horizontal bas intérieur (一), 6. Horizontal bas fermeture (一). 
\zh{收} (6 traits) : gauche \zh{丩} (3 traits), droite clé \zh{攵} (4 traits : 横、竖、撇、捺) — \textbf{attention :} \zh{攵} s'écrit en 4 traits, le trait final est une pointe (捺).
\item \emph{huánjìng} → \zh{环境}. 
\zh{环} (8 traits) : clé \zh{王} (4 traits : 一、一、丨、一), puis \zh{不} (4 traits). 
\zh{境} (14 traits) : clé \zh{土} (3 traits), puis \zh{竟} (11 traits : \zh{音} 9 traits + \zh{儿} 2 traits).
\item \emph{zài lìyòng} → \zh{再利用}. 
\zh{再} (6 traits) : \zh{冂} + \zh{冉} (structure haut-bas). 
\zh{利} (7 traits) : \zh{禾} (5 traits) + \zh{刂} (2 traits, vertical puis pointe). 
\zh{用} (5 traits) : haut \zh{卜}, milieu deux horizontaux, bas fermeture.
\item \emph{fēnlèi} → \zh{分类}. 
\zh{分} (4 traits) : \zh{八} + \zh{刀}. 
\zh{类} (18 traits) : clé \zh{米} (6 traits) + \zh{大} (3 traits) + \zh{页} (6 traits) — \textbf{note :} \zh{类} est dense, écrire grand.
\end{itemize}

\textbf{(b) Correction de la phrase.} \zh{我们在学校应该在利用纸。}
\begin{itemize}
\item \textbf{Faute 1 : Confusion \zh{在} / \zh{再}.} « Réutiliser » = \zh{再利用} (\zh{再} = de nouveau). \zh{在} indique le lieu (\zh{在学校}) ou le progressif (\zh{在做}).
\item \textbf{Faute 2 : Ordre des mots (place de l'auxiliaire).} L'auxiliaire modal \zh{应该} se place \textbf{avant} le complément de lieu long ou juste après le sujet. La place naturelle : Sujet + Lieu + Modal + Verbe.
\end{itemize}
\textbf{Phrase corrigée :}
\zh{我们在学校应该再利用纸。} \emph{Wǒmen zài xuéxiào yīnggāi zài lìyòng zhǐ.} « À l'école, nous devrions réutiliser le papier. »

\textbf{Conclusion :} La confusion \zh{在}/\zh{再} est l'homophone le plus pénalisé dans ce thème ; mémoriser le binôme \zh{再利用} comme un bloc unique l'évite définitivement.
\end{exoresolu}

\courssub{Fiche méthode 3 — Le thème : traduire du français vers le chinois}

\begin{methode}[Traduire une phrase française en chinois (Thème)]
\begin{enumerate}
\item \cle{Repérer le squelette SVO} : Sujet + Verbe + Objet. Le chinois suit l'ordre SVO comme le français : point d'appui majeur.
\item \cle{Placer les compléments de TEMPS et de LIEU AVANT le verbe} : « Nous recyclons le papier \emph{à l'école} » $\rightarrow$ \zh{我们\textbf{在学校}回收纸。} (Jamais après le verbe).
\item \cle{Traduire les modaux} : devoir $\rightarrow$ \zh{应该} / \zh{得} ; pouvoir $\rightarrow$ \zh{可以} / \zh{能} ; vouloir $\rightarrow$ \zh{想} / \zh{要}. Le modal se place \textbf{avant le verbe}, sans préposition.
\item \cle{Ajouter le CLASSIFICATEUR} dès qu'un nom est compté ou désigné : une bouteille $\rightarrow$ \zh{一个瓶子} / \zh{一个塑料瓶} ; cette boîte $\rightarrow$ \zh{这个盒子}. \textbf{Jamais} de \zh{一} + nom nu.
\item \cle{Gérer \zh{的} (particule de liaison)} : 
\begin{itemize}
\item Obligatoire : Possession (\zh{我们的环境}), Adjectif polysyllabique + Nom (\zh{干净的环境}), Subordonnée relative (\zh{我们回收的纸}).
\item Souvent omis : Adjectif monosyllabique très lié (\zh{好书}, \zh{大山}), Pronom + Nom de parenté/affection (\zh{我妈}, \zh{我们学校}).
\end{itemize}
\item \cle{Relire à voix haute} en vérifiant les tons des mots-clés : \emph{lājī} (1-1), \emph{huíshōu} (2-1), \emph{huánjìng} (2-4), \emph{yuèláiyuè} (4-2-4), \emph{sùliàopíng} (4-4-2).
\end{enumerate}
\end{methode}

\begin{exoresolu}[Thème guidé : trois phrases types Bac]
\textbf{Énoncé.} Traduisez en chinois : 
(1) « Les élèves recyclent le papier et le plastique à l'école. » 
(2) « Nous devons protéger l'environnement de notre ville. » 
(3) « Il y a de plus en plus de bouteilles en plastique dans les rues de Douala. »
\tcblower
\textbf{(1) Squelette :} \zh{学生们} (sujet) + \zh{回收} (verbe) + \zh{纸和塑料} (objet). Complément de lieu \zh{在学校} placé \textbf{avant} le verbe.

\zh{学生们在学校回收纸和塑料。} \emph{Xuéshengmen zài xuéxiào huíshōu zhǐ hé sùliào.}

\textbf{(2) « Devoir » =} \zh{应该} (recommandation morale), placé avant \zh{保护}. « De notre ville » = complément du nom avec \zh{的} : \zh{我们城市的环境}.

\zh{我们应该保护我们城市的环境。} \emph{Wǒmen yīnggāi bǎohù wǒmen chéngshì de huánjìng.}

\textbf{(3) « Il y a » =} \zh{有} ; « de plus en plus de » = \zh{越来越多(的)} ; Complément de lieu \zh{在杜阿拉的街上} avant le verbe \zh{有}.

\zh{在杜阿拉的街上，有越来越多的塑料瓶。} \emph{Zài Dù'ālā de jiē shang, yǒu yuèláiyuè duō de sùliàopíng.}

\textbf{Pièges évités :} 
\begin{itemize}
\item Ne pas écrire \zh{学生们回收纸在学校} (calque ordre français).
\item Ne pas oublier \zh{的} dans \zh{我们城市的环境} (nom + \zh{的} + nom).
\item Ne pas traduire « de plus en plus » par deux mots séparés : \zh{越来越} est un bloc inséparable.
\item \zh{塑料} seul = matière (plastique) ; \zh{塑料瓶} = objet (bouteille en plastique). Préférer le nom concret.
\end{itemize}
\end{exoresolu}

\courssub{Fiche méthode 4 — La version : traduire du chinois vers le français}

\begin{methode}[Traduire un texte chinois en français (Version)]
\begin{enumerate}
\item \cle{Lire le texte entier une première fois} sans traduire, pour dégager le sens global (Qui ? Quoi ? Où ? Quand ?).
\item \cle{Découper en segments de sens} (Propositions) : Sujet / Compléments (Temps, Lieu, Manière) / Verbe / Objet. \textbf{Règle d'or :} Tout ce qui précède le verbe en chinois (temps, lieu, manière) doit être replacé \textbf{après le verbe} (ou en tête de phrase) en français.
\item \cle{Traduire les mots-outils avec précision} : 
\begin{itemize}
\item \zh{了} (accompli/réalisé) $\rightarrow$ Passé composé / Imparfait selon contexte.
\item \zh{过} (expérience) $\rightarrow$ « avoir déjà » / Passé composé.
\item \zh{着} (durée/état) $\rightarrow$ « en train de » / Imparfait / Participe présent.
\item \zh{的} (liaison) $\rightarrow$ « de » / adjectif / relatif « que/qui ».
\end{itemize}
\item \cle{Restituer les classificateurs par l'article français} : \zh{一个瓶子} $\rightarrow$ « une bouteille » ; \zh{这些盒子} $\rightarrow$ « ces boîtes ». Ne pas traduire le classificateur mot à mot (« un individu bouteille »).
\item \cle{Rédiger en français naturel, soutenu} : Une version mot à mot est pénalisée. Vérifier accords, temps, connecteurs logiques (car, donc, mais, ainsi).
\end{enumerate}
\end{methode}

\begin{exoresolu}[Version : court texte sur une action scolaire]
\textbf{Énoncé.} Traduisez en français :

\zh{上个星期，我们学校的学生在雅温得的一个市场旁边回收了很多塑料瓶和纸。老师告诉我们，回收可以保护环境，也可以帮助我们的城市变得更干净。}
\tcblower
\textbf{Étape 1 — Sens global :} Des élèves ont fait une opération de recyclage près d'un marché ; le professeur explique l'intérêt.

\textbf{Étape 2 — Segmentation et analyse.}
\begin{itemize}
\item \zh{上个星期} (Complément de temps) : « La semaine dernière ».
\item \zh{我们学校的学生} (Sujet + \zh{的}) : « Les élèves de notre école ».
\item \zh{在雅温得的一个市场旁边} (Complément de lieu) : « À côté d'un marché de Yaoundé » (\zh{旁边} = à côté ; \zh{一个} + \zh{市场} : classificateur restitué par « un »).
\item \zh{回收了} (Verbe + \zh{了}) : \zh{了} = action accomplie $\rightarrow$ Passé composé « ont recyclé ».
\item \zh{很多塑料瓶和纸} (Objet) : « Beaucoup de bouteilles en plastique et de papier ».
\item \zh{老师告诉我们} : « Le professeur nous a dit / nous a expliqué ».
\item \zh{回收可以保护环境，也可以帮助……} : \zh{可以} = « peut / permet de » ; \zh{也可以} = « peut aussi » ; \zh{帮助……变得} = « aider … à devenir ».
\item \zh{更干净} : \zh{更} (comparatif) + \zh{干净} (propre) = « plus propre ».
\end{itemize}

\textbf{Étape 3 — Traduction rédigée (niveau attendu).}

« La semaine dernière, les élèves de notre école ont recyclé beaucoup de bouteilles en plastique et de papier à côté d'un marché de Yaoundé. Le professeur nous a expliqué que le recyclage permet de protéger l'environnement et peut aussi aider notre ville à devenir plus propre. »

\textbf{Points de vigilance justifiés :} 
\begin{itemize}
\item \zh{了} impose le passé (pas de présent narratif).
\item \zh{可以…也可以…} se rend élégamment par « permet de… et peut aussi… » (éviter « peut… peut aussi »).
\item \zh{变得更干净} : structure \zh{变得 + Adj} = « devenir [Adj] » ; \zh{更} = comparatif.
\item Une version littérale type « Recyclage peut protéger environnement » (sans articles, sans conjugaison) est sanctionnée en expression française.
\end{itemize}
\end{exoresolu}

\courssub{Fiche méthode 5 — Utiliser les connecteurs pour structurer le texte}

\begin{methode}[Les connecteurs utiles du texte argumentatif simple (Niveau Tle)]
\begin{enumerate}
\item \cle{Addition / Énumération} : \zh{也} (aussi, entre verbes/phrases), \zh{还有} (et aussi, en plus), \zh{和} / \zh{跟} (et, \textbf{uniquement entre noms/pronoms}).
\item \cle{Cause} : \zh{因为……所以……} (Parce que… donc…). \textbf{En chinois, les deux termes coexistent souvent dans la même phrase}, contrairement au français où « parce que » et « donc » font double emploi.
\item \cle{Conséquence / Conclusion} : \zh{所以} (donc, en milieu de phrase), \zh{这样} (ainsi, en tête de phrase), \zh{因此} (c'est pourquoi, plus formel).
\item \cle{Opposition / Concession} : \zh{但是} / \zh{可是} (mais), \zh{虽然……但是……} (bien que… pourtant…).
\item \cle{Progression / Étapes} : \zh{首先……然后……最后……} (D'abord… ensuite… enfin…).
\item \cle{Règle d'or} : Un connecteur clair par phrase suffit. Ne pas surcharger.
\end{enumerate}
\end{methode}

\begin{popfigure}
\begin{center}\tcbox[colback=popink,colframe=popink,coltext=white,arc=9pt,boxsep=4pt,left=6pt,right=6pt]{\bfseries\small Connecteurs Textuels}\end{center}
\vspace{-2pt}
\begin{tcbraster}[raster columns=3,raster equal height=rows,raster column skip=6pt,raster row skip=6pt,enhanced,blankest]
\begin{tcolorbox}[enhanced,colback=white,colframe=popgreen,boxrule=0.9pt,arc=3pt,title={Addition},fonttitle=\bfseries\scriptsize,coltitle=white,colbacktitle=popgreen,left=4pt,right=4pt,top=2pt,bottom=2pt,boxsep=1pt,toptitle=1.5pt,bottomtitle=1.5pt,halign=left]
\begin{itemize}[nosep,leftmargin=*,itemsep=1pt,font=\scriptsize]
\item \zh{也} (aussi)
\item \zh{还有} (en plus)
\item \zh{和/跟} (noms)
\end{itemize}
\end{tcolorbox}
\begin{tcolorbox}[enhanced,colback=white,colframe=poppurple,boxrule=0.9pt,arc=3pt,title={Cause},fonttitle=\bfseries\scriptsize,coltitle=white,colbacktitle=poppurple,left=4pt,right=4pt,top=2pt,bottom=2pt,boxsep=1pt,toptitle=1.5pt,bottomtitle=1.5pt,halign=left]
\begin{itemize}[nosep,leftmargin=*,itemsep=1pt,font=\scriptsize]
\item \zh{因为...所以...}
\end{itemize}
\end{tcolorbox}
\begin{tcolorbox}[enhanced,colback=white,colframe=poppink,boxrule=0.9pt,arc=3pt,title={Conséquence},fonttitle=\bfseries\scriptsize,coltitle=white,colbacktitle=poppink,left=4pt,right=4pt,top=2pt,bottom=2pt,boxsep=1pt,toptitle=1.5pt,bottomtitle=1.5pt,halign=left]
\begin{itemize}[nosep,leftmargin=*,itemsep=1pt,font=\scriptsize]
\item \zh{所以} (donc)
\item \zh{这样} (ainsi)
\item \zh{因此} (formel)
\end{itemize}
\end{tcolorbox}
\begin{tcolorbox}[enhanced,colback=white,colframe=popgold,boxrule=0.9pt,arc=3pt,title={Opposition},fonttitle=\bfseries\scriptsize,coltitle=white,colbacktitle=popgold,left=4pt,right=4pt,top=2pt,bottom=2pt,boxsep=1pt,toptitle=1.5pt,bottomtitle=1.5pt,halign=left]
\begin{itemize}[nosep,leftmargin=*,itemsep=1pt,font=\scriptsize]
\item \zh{但是/可是}
\item \zh{虽然...但...}
\end{itemize}
\end{tcolorbox}
\begin{tcolorbox}[enhanced,colback=white,colframe=popdark,boxrule=0.9pt,arc=3pt,title={Étapes},fonttitle=\bfseries\scriptsize,coltitle=white,colbacktitle=popdark,left=4pt,right=4pt,top=2pt,bottom=2pt,boxsep=1pt,toptitle=1.5pt,bottomtitle=1.5pt,halign=left]
\begin{itemize}[nosep,leftmargin=*,itemsep=1pt,font=\scriptsize]
\item \zh{首先}
\item \zh{然后}
\item \zh{最后}
\end{itemize}
\end{tcolorbox}
\end{tcbraster}

\legende{Carte mentale des connecteurs logiques pour l'expression écrite (Tle).}
\end{popfigure}

\begin{exoresolu}[Relier des phrases avec les connecteurs appropriés]
\textbf{Énoncé.} Reliez les phrases suivantes avec le connecteur le plus adapté (\zh{因为…所以…} / \zh{但是} / \zh{首先…然后…}) : 
(1) \zh{塑料污染环境。我们应该回收塑料。} 
(2) \zh{回收很重要。很多人还不回收。} 
(3) \zh{我们把纸放在一起。我们把塑料瓶放在一起。}
\tcblower
\textbf{(1) Relation de CAUSE :} La pollution est la cause, le recyclage la conséquence $\rightarrow$ \zh{因为…所以…}.

\zh{因为塑料污染环境，所以我们应该回收塑料。} \emph{Yīnwèi sùliào wūrǎn huánjìng, suǒyǐ wǒmen yīnggāi huíshōu sùliào.} « Parce que le plastique pollue l'environnement, nous devons le recycler. »

\textbf{(2) Relation d'OPPOSITION :} Importance vs comportement contraire $\rightarrow$ \zh{但是}.

\zh{回收很重要，但是很多人还不回收。} \emph{Huíshōu hěn zhòngyào, dànshì hěn duō rén hái bù huíshōu.} « Le recyclage est très important, mais beaucoup de gens ne recyclent pas encore. » (\zh{还} = « encore », placé \textbf{avant} la négation \zh{不}.)

\textbf{(3) Relation de SUCCESSION (Étapes) :} Deux actions successives $\rightarrow$ \zh{首先…然后…}.

\zh{首先我们把纸放在一起，然后我们把塑料瓶放在一起。} \emph{Shǒuxiān wǒmen bǎ zhǐ fàng zài yìqǐ, ránhòu wǒmen bǎ sùliàopíng fàng zài yìqǐ.} « D'abord nous regroupons le papier, ensuite nous regroupons les bouteilles en plastique. »

\textbf{Conclusion :} Ces trois schémas (\zh{因为…所以…}, \zh{但是}, \zh{首先…然后…}) suffisent à construire un paragraphe argumentatif complet et cohérent au niveau Terminale.
\end{exoresolu}

\courssub{Entraînement : Thème et Version (Exercices types Bac)}

\begin{exercice}[Thème — Français $\rightarrow$ Chinois]
Traduisez les phrases suivantes en chinois (caractères + pinyin). Respectez l'ordre des mots, les classificateurs et les tons.
\begin{enumerate}
\item Les habitants de Douala jettent beaucoup de déchets dans la rue.
\item Nous devons trier les déchets : le papier d'un côté, le plastique de l'autre.
\item Si tout le monde recycle, l'environnement deviendra de plus en plus beau.
\item Hier, j'ai réutilisé une vieille bouteille en plastique pour faire un pot de fleurs.
\item Le professeur nous a dit que protéger l'environnement, c'est protéger notre avenir.
\end{enumerate}
\end{exercice}

\begin{corrige}[Exercice Thème]
\begin{enumerate}
\item \zh{杜阿拉的居民在街上扔很多垃圾。} \emph{Dù'ālā de jūmín zài jiē shang rēng hěn duō lājī.} (\zh{居民} habitants ; \zh{扔} jeter ; lieu avant verbe).
\item \zh{我们应该分类垃圾：纸放一边，塑料放另一边。} \emph{Wǒmen yīnggāi fēnlèi lājī: zhǐ fàng yībiān, sùliào fàng lìngyībiān.} (\zh{分类} trier ; \zh{一边/另一边} d'un côté/de l'autre).
\item \zh{如果大家都回收，环境会越来越美丽。} \emph{Rúguǒ dàjiā dōu huíshōu, huánjìng huì yuèláiyuè měilì.} (\zh{如果…会…} si… alors… ; \zh{大家} tout le monde ; \zh{都} tous).
\item \zh{昨天，我再利用了一个旧塑料瓶，做花盆。} \emph{Zuótiān, wǒ zài lìyòng le yí gè jiù sùliàopíng, zuò huāpén.} (\zh{了} après verbe pour accompli ; \zh{一个} classif. ; \zh{做} faire/fabriquer).
\item \zh{老师告诉我们，保护环境就是保护我们的未来。} \emph{Lǎoshī gàosu wǒmen, bǎohù huánjìng jiùshì bǎohù wǒmen de wèilái.} (\zh{就是} c'est exactement / équivaut à ; \zh{未来} avenir).
\end{enumerate}
\end{corrige}

\begin{exercice}[Version — Chinois $\rightarrow$ Français]
Traduisez le texte suivant en français naturel et correct.
\zh{现在的雅温得，垃圾问题很严重。因为很多人不分类垃圾，把什么都扔在一个垃圾桶里。这污染了土壤和水。我们中学生应该从我做起，少用塑料袋，多用布袋子。同时，我们要把废纸、塑料瓶分开，送去回收站。只有大家一起行动，我们的家园才会变得更干净、更美好。}
\end{exercice}

\begin{corrige}[Exercice Version]
\textbf{Traduction proposée :}
« De nos jours à Yaoundé, le problème des déchets est grave. Parce que beaucoup de gens ne trient pas leurs déchets et jettent tout dans une seule poubelle, cela a pollué les sols et l'eau. Nous, collégiens/lycéens, devrions commencer par nous-mêmes : utiliser moins de sacs en plastique, utiliser plus de sacs en tissu. En même temps, nous devons séparer le papier usagé et les bouteilles en plastique et les apporter au centre de recyclage. Seulement si tout le monde agit ensemble, notre foyer/planète deviendra plus propre et plus beau. »

\textbf{Notes de traduction :}
\begin{itemize}
\item \zh{现在的雅温得} : « De nos jours à Yaoundé » / « L'actuel Yaoundé ».
\item \zh{从我做起} : locution « commencer par soi-même / par moi ».
\item \zh{少用…多用…} : structure « moins utiliser… plus utiliser… ».
\item \zh{废纸} : « papier usagé / déchets de papier ».
\item \zh{送去回收站} : \zh{送去} « apporter (quelque part) » ; \zh{回收站} « centre de tri / déchetterie ».
\item \zh{只有…才…} : structure « seulement si… alors… » (condition nécessaire).
\item \zh{家园} : « foyer / patrie / chez-nous » (plus affectif que \zh{家}).
\end{itemize}
\end{corrige}

\courssub{Production guidée et grille d'évaluation}

\begin{experience}[Atelier d'écriture : « Ma ville plus propre »]
\textbf{Consigne :} En binôme, rédigez un court texte (80-100 caractères) pour une affiche à coller dans le couloir du lycée. Sujet : « Comment les élèves de notre lycée peuvent réduire les déchets et recycler ? »
\textbf{Étapes :}
\begin{enumerate}
\item \textbf{Brainstorming (5 min) :} Liste 10 mots-clés (vocabulaire leçon + \zh{减少} réduire, \zh{布袋} sac tissu, \zh{水瓶} gourde, \zh{食堂} cantine).
\item \textbf{Plan commun (5 min) :} 1. Constat (cantine/cour). 2. 3 actions concrètes (tri, gourde, sac tissu). 3. Slogan final (\zh{保护环境，从我做起}).
\item \textbf{Rédaction individuelle (10 min) :} Chaque élève écrit sa version sur papier brouillon.
\item \textbf{Relecture croisée (5 min) :} Échangez les copies. Cochez : \ding{51} Ordre mots (Lieu/Temps avant V) \ding{51} Classificateurs \ding{51} \zh{在}/\zh{再} \ding{51} Connecteurs (\zh{因为/所以}, \zh{首先/然后}) \ding{51} Caractères lisibles.
\item \textbf{Affichage :} Les 3 meilleures affiches (vote classe) sont plastifiées et affichées.
\end{enumerate}
\end{experience}

\begin{aretenir}[Grille d'auto-évaluation / Évaluation enseignant (sur 20 pts)]
\begin{center}
\begin{tabularx}{\linewidth}{|X|X|X|}
\hline
\rowcolor{popblueL} \textbf{Critère} & \textbf{Indicateurs de réussite (4-5 pts)} & \textbf{Points} \\
\hline
\textbf{Contenu \& Structure} (5 pts) & Idées claires (Constat, Action, Conclusion), plan logique, connecteurs variés (\zh{因为/所以}, \zh{首先/然后}, \zh{这样}) & /5 \\
\hline
\textbf{Grammaire \& Syntaxe} (5 pts) & Ordre SVO + Compléments (T/L) avant V ; Modaux bien placés ; \zh{把} correct ; \zh{越来越} + Adj ; pas de calque français & /5 \\
\hline
\textbf{Vocabulaire \& Caractères} (5 pts) & Vocabulaire thème employé (\zh{回收, 分类, 再利用, 环境, 污染}) ; Caractères corrects (traits, radicaux) ; Pas confusion \zh{在}/\zh{再}, \zh{境}/\zh{镜}, \zh{瓶}/\zh{饼} & /5 \\
\hline
\textbf{Présentation \& Pinyin} (5 pts) & Écriture soignée, caractères proportionnés ; Pinyin avec tons exacts sur tout le texte ; Ponctuation chinoise pleine chasse & /5 \\
\hline
\end{tabularx}
\end{center}
\textbf{Seuil de réussite :} 12/20. \textbf{Objectif excellence :} 16+/20 (texte fluide, zéro faute de caractère, connecteurs naturels).
\end{aretenir}

\courssub{Éléments socioculturels : Cameroun — Chine}

\begin{savaistu}[Gestion des déchets : regards croisés Cameroun / Chine]
\textbf{Cameroun :} Défis majeurs à Yaoundé et Douala (décharges à ciel ouvert comme Nkolfoulou, collecte irrégulière, brûlage des ordures, pollution des cours d'eau). Initiatives : \emph{HYSACAM} (collecte), ONG locales de tri sélectif (plastique, métal), projet « Yaoundé Ville Propre ». Sensibilisation dans les écoles (clubs environnement).
\textbf{Chine :} Politique nationale \emph{垃圾分类} (tri obligatoire des déchets) déployée depuis 2019 (Shanghai pilote, puis 46 villes clés). 4 catégories : \zh{可回收物} (recyclables), \zh{有害垃圾} (dangereux), \zh{厨余垃圾} (organiques/cuisine), \zh{其他垃圾} (autres). Amendes pour non-respect. Applications mobiles pour guider le tri. Industrie du recyclage massive (leader mondial traitement déchets électroniques, plastique).
\textbf{Comparaison :} La Chine a imposé le tri par la loi (contrainte forte + technologie) ; le Cameroun mise davantage sur la sensibilisation et le secteur informel (chiffonniers \zh{拾荒者} jouent rôle clé). Les deux pays partagent l'urgence climatique et l'objectif « Ville propre ».
\textbf{Vocabulaire bonus :} \zh{垃圾分类} (tri sélectif), \zh{可回收物} (recyclables), \zh{厨余垃圾} (biodéchets), \zh{填埋场} (décharge/centre d'enfouissement), \zh{焚烧} (incinération), \zh{循环经济} (économie circulaire).
\end{savaistu}

\begin{attention}[Les 5 erreurs qui coûtent des points au Bac — Récapitulatif final]
\begin{enumerate}
\item \cle{Calquer l'ordre français (Lieu/Temps après verbe)} : \textbf{Faux :} \zh{我们回收纸在学校。} \textbf{Juste :} \zh{我们在学校回收纸。} (Règle : Compléments de Temps/Lieu $\rightarrow$ Avant Verbe).
\item \cle{Confondre les homophones \zh{在} et \zh{再}} : \zh{再利用} (réutiliser) s'écrit avec \zh{再} (de nouveau) ; \zh{在} sert pour le lieu (\zh{在学校}) ou le progressif (\zh{在回收}). Une faute sur ce mot dans ce thème = 0 pt souvent.
\item \cle{Oublier le CLASSIFICATEUR} : \textbf{Faux :} \zh{一瓶子、两纸。} \textbf{Juste :} \zh{一个瓶子、两张纸。} (Tout nom compté : Nombre + Classif + Nom). \zh{张} pour objets plats (papier, table, lit, carte).
\item \cle{Mal placer les MODAUX et ADVERBES} : \zh{应该, 可以, 也, 还, 都, 就} se placent \textbf{AVANT LE VERBE} (et après le sujet). \textbf{Faux :} \zh{我们回收也应该纸。} \textbf{Juste :} \zh{我们也应该回收纸。}
\item \cle{Négliger TONS et CARACTÈRES PROCHES} : \zh{境} (env.) $\neq$ \zh{镜} (miroir) ; \zh{瓶} (bouteille) $\neq$ \zh{饼} (galette) ; \zh{圾} $\neq$ \zh{级} ; \zh{再} $\neq$ \zh{在}. Pinyin sans tons = incomplet. \textbf{Relisez chaque caractère trait par trait avant de rendre la copie.}
\end{enumerate}
\end{attention}

\newpage

\coursec{Exercices}

\courssub{Je vérifie mes connaissances}

\begin{exercice}[QCM : le lexique du recyclage]
Pour chaque question, choisis la bonne réponse (a, b ou c).
\begin{enumerate}
\item \zh{回收} (huíshōu) signifie :
\begin{enumerate}
\item jeter \item recycler / récupérer \item brûler
\end{enumerate}
\item Comment dit-on « les déchets » en chinois ?
\begin{enumerate}
\item \zh{环境} \item \zh{垃圾} \item \zh{资源}
\end{enumerate}
\item \zh{塑料瓶} (sùliào píng) désigne :
\begin{enumerate}
\item une bouteille en verre \item une bouteille en plastique \item une boîte en carton
\end{enumerate}
\item Quel verbe signifie « protéger » ?
\begin{enumerate}
\item \zh{保护} \item \zh{节约} \item \zh{分类}
\end{enumerate}
\end{enumerate}
\end{exercice}

\begin{exercice}[Vrai ou faux justifié]
Dis si chaque affirmation est vraie (V) ou fausse (F), puis justifie en une phrase en français.
\begin{enumerate}
\item \zh{垃圾分类} (lājī fēnlèi) signifie « le tri des déchets ».
\item Dans la phrase \zh{我们把纸和塑料分开放。}, le mot \zh{把} sert à marquer le passé.
\item \zh{既……又……} (jì……yòu……) permet d'exprimer deux qualités ou deux actions simultanées (« à la fois... et... »).
\item Le caractère \zh{环} (huán, environnement) a pour radical la clé de la terre \zh{土}.
\end{enumerate}
\end{exercice}

\begin{exercice}[Vocabulaire : complète le tableau]
Complète le tableau suivant en écrivant les caractères chinois correspondant à chaque traduction (le pinyin est donné).
\begin{center}
\begin{tabularx}{\linewidth}{|X|X|X|X|}
\hline
\rowcolor{popblueL} Caractères & Pinyin & Nature & Traduction \\
\hline
 & huánjìng & nom & l'environnement \\
\hline
 & zīyuán & nom & les ressources \\
\hline
 & bōli & nom & le verre (matière) \\
\hline
 & jiéyuē & verbe & économiser \\
\hline
 & wūrǎn & nom / verbe & la pollution / polluer \\
\hline
\end{tabularx}
\end{center}
\end{exercice}

\begin{exercice}[Associe le pinyin aux caractères]
Relie chaque caractère ou mot à son pinyin.
\begin{enumerate}
\item \zh{回收} \hspace{1cm} a. fēnlèi
\item \zh{分类} \hspace{1cm} b. sùliào
\item \zh{塑料} \hspace{1cm} c. huíshōu
\item \zh{保护} \hspace{1cm} d. lājī
\item \zh{垃圾} \hspace{1cm} e. bǎohù
\end{enumerate}
\end{exercice}

\courssub{Je m'entraîne}

\begin{exercice}[Discrimination de caractères proches]
Complète chaque mot avec le bon caractère choisi parmi les paires proposées : \zh{圾}/\zh{级}, \zh{环}/\zh{坏}, \zh{收}/\zh{放}.
\begin{enumerate}
\item 垃\underline{\hspace{1cm}} (lājī, déchets)
\item \underline{\hspace{1cm}}境 (huánjìng, environnement)
\item \underline{\hspace{1cm}}到 (shōudào, recevoir)
\item 年\underline{\hspace{1cm}} (niánjí, niveau scolaire)
\item 破\underline{\hspace{1cm}} (pòhuài, détruire)
\item \underline{\hspace{1cm}}下 (fàngxià, poser)
\end{enumerate}
\end{exercice}

\begin{textebox}[Texte à trous]
\underline{\hspace{1.2cm}}我们只有一个地球，\underline{\hspace{1.2cm}}我们要保护环境。我们\underline{\hspace{1.2cm}}垃圾分类好，\underline{\hspace{1.2cm}}能节约资源，\underline{\hspace{1.2cm}}能减少污染。这是每个人\underline{\hspace{1.2cm}}责任。
\end{textebox}

\begin{exercice}[Texte à trous : mots-outils et connecteurs]
Complète le texte ci-dessus avec les mots de la liste : \zh{把}, \zh{既}, \zh{又}, \zh{因为}, \zh{所以}, \zh{的}. Puis traduis le texte complété en français.
\end{exercice}

\begin{exercice}[Remise en ordre et traduction]
Les quatre phrases suivantes sont dans le désordre. Remets-les dans l'ordre logique grâce aux connecteurs (\zh{首先} d'abord, \zh{然后} ensuite, \zh{最后} enfin), recopie le texte obtenu, puis traduis-le en français.
\begin{enumerate}
\item \zh{最后，我们把分好的垃圾送到回收站。}\\ \emph{(Zuìhòu, wǒmen bǎ fēn hǎo de lājī sòng dào huíshōuzhàn.)}
\item \zh{然后，我们把纸、塑料和玻璃分开放。}\\ \emph{(Ránhòu, wǒmen bǎ zhǐ, sùliào hé bōli fēnkāi fàng.)}
\item \zh{这样做既保护环境，又节约资源。}\\ \emph{(Zhèyàng zuò jì bǎohù huánjìng, yòu jiéyuē zīyuán.)}
\item \zh{首先，我们学习垃圾分类的知识。}\\ \emph{(Shǒuxiān, wǒmen xuéxí lājī fēnlèi de zhīshi.)}
\end{enumerate}
\end{exercice}

\begin{exercice}[Écriture des caractères du thème]
Recopie cinq fois chacun des mots suivants en respectant l'ordre des traits (de haut en bas, de gauche à droite, l'horizontal avant le vertical) : \zh{回收}, \zh{环境}, \zh{塑料}, \zh{分类}. Entoure ensuite le radical de chaque caractère et nomme-le en français (par exemple : clé du jade, clé de la terre, clé du couteau).
\end{exercice}

\courssub{Je me prépare au Bac}

\begin{textebox}[Le tri des déchets à l'école]
现在，很多学校都教学生垃圾分类。同学们把纸、塑料和玻璃分开放。这样既保护环境，又节约资源。有的学校还用旧纸做手工，非常有意思。\\
\emph{(Xiànzài, hěn duō xuéxiào dōu jiāo xuésheng lājī fēnlèi. Tóngxuémen bǎ zhǐ, sùliào hé bōli fēnkāi fàng. Zhèyàng jì bǎohù huánjìng, yòu jiéyuē zīyuán. Yǒu de xuéxiào hái yòng jiù zhǐ zuò shǒugōng, fēicháng yǒu yìsi.)}\\
« Aujourd'hui, beaucoup d'écoles enseignent le tri des déchets aux élèves. Les camarades séparent le papier, le plastique et le verre. Ainsi, on protège à la fois l'environnement et on économise les ressources. Certaines écoles utilisent aussi le vieux papier pour faire de l'artisanat, c'est très intéressant. »
\end{textebox}

\begin{exercice}[Compréhension de l'écrit : texte et questions]
Lis le texte ci-dessus, puis réponds aux questions en français.
\begin{enumerate}
\item Qu'est-ce que beaucoup d'écoles enseignent aux élèves ?
\item Cite les trois matières que les élèves séparent.
\item Quels sont les deux avantages du tri mentionnés dans le texte ?
\item \textbf{Question de langue :} relève dans le texte une phrase contenant la structure \zh{把} et explique sa construction (Sujet + \zh{把} + complément d'objet + verbe + élément).
\item \textbf{Question de langue :} que signifie la structure \zh{既……又……} ? Construis une phrase personnelle avec cette structure sur le thème de l'environnement.
\end{enumerate}
\end{exercice}

\begin{exercice}[Thème et version type Bac]
\textbf{A. Thème (français → chinois).} Traduis les phrases suivantes en chinois (caractères obligatoires, pinyin facultatif) :
\begin{enumerate}
\item « Nous devons trier les déchets et recycler les bouteilles en plastique. »
\item « Parce que nous n'avons qu'une Terre, nous devons la protéger. »
\item « Les élèves de notre lycée mettent le papier dans la boîte bleue. » (utilise la structure \zh{把})
\end{enumerate}
\textbf{B. Version (chinois → français).} Traduis en français le texte « Pour protéger l'environnement » ci-dessous.
\end{exercice}

\begin{textebox}[Pour protéger l'environnement]
为了保护环境，我们学校每个月都回收纸和塑料瓶。同学们觉得这个活动很有意义。\\
\emph{(Wèile bǎohù huánjìng, wǒmen xuéxiào měi ge yuè dōu huíshōu zhǐ hé sùliào píng. Tóngxuémen juéde zhè ge huódòng hěn yǒu yìyì.)}
\end{textebox}

\begin{exercice}[Expression écrite type Bac : rédaction]
\textbf{Sujet :} « Le recyclage dans ma ville » (\zh{我们城市的回收}).\\
Rédige un texte en chinois de \textbf{6 phrases minimum} (au moins 60 caractères) en respectant les consignes suivantes :
\begin{enumerate}
\item présente la situation du recyclage dans ta ville (par exemple Yaoundé ou Douala) ;
\item utilise \textbf{au moins deux connecteurs} parmi : \zh{首先}, \zh{然后}, \zh{最后}, \zh{因为……所以……}, \zh{既……又……} ;
\item utilise \textbf{au moins une structure} \zh{把} ;
\item emploie au moins quatre mots du lexique de la leçon : \zh{回收}, \zh{环境}, \zh{垃圾}, \zh{分类}, \zh{塑料}, \zh{保护}, \zh{节约}, \zh{资源} ;
\item termine par une phrase d'opinion personnelle (par exemple avec \zh{我觉得……}).
\end{enumerate}
\textbf{Grille d'évaluation (sur 10) :} correction des caractères (2 pts), grammaire et ordre des mots (3 pts), respect des consignes et connecteurs (3 pts), richesse du vocabulaire (2 pts).
\end{exercice}

\begin{methode}[Réussir la rédaction en 5 étapes]
\begin{enumerate}
\item \textbf{Plan :} situation (phrase 1-2) → actions (phrase 3-4 avec \zh{首先}, \zh{然后}) → résultat ou opinion (phrase 5-6 avec \zh{我觉得}).
\item \textbf{Connecteurs :} place le connecteur en tête de phrase, suivi d'une virgule chinoise \zh{，}.
\item \textbf{Structure \zh{把} :} vérifie que le verbe n'est jamais seul après l'objet : \zh{把垃圾分开放} (et non \zh{把垃圾放}).
\item \textbf{Relecture :} contrôle les tons des mots-clés (\zh{huíshōu}, \zh{huánjìng}, \zh{fēnlèi}), l'ordre Sujet-Verbe-Objet et les caractères proches (\zh{环}/\zh{坏}, \zh{圾}/\zh{级}).
\item \textbf{Comptage :} 60 caractères ≈ 6 phrases de 10 caractères ; les connecteurs et la ponctuation ne comptent pas comme des caractères de lexique.
\end{enumerate}
\end{methode}

\begin{exemplebox}[Modèle de rédaction corrigé]
\zh{我住在雅温得。我们城市的垃圾很多，所以回收非常重要。首先，我们学校教学生垃圾分类。然后，同学们把纸、塑料和玻璃分开放。最后，有人把分好的垃圾送到回收站。这样做既保护环境，又节约资源。我觉得每个人都应该参加这个活动。}\\
\emph{(Wǒ zhù zài Yǎwēndé. Wǒmen chéngshì de lājī hěn duō, suǒyǐ huíshōu fēicháng zhòngyào. Shǒuxiān, wǒmen xuéxiào jiāo xuésheng lājī fēnlèi. Ránhòu, tóngxuémen bǎ zhǐ, sùliào hé bōli fēnkāi fàng. Zuìhòu, yǒu rén bǎ fēn hǎo de lājī sòng dào huíshōuzhàn. Zhèyàng zuò jì bǎohù huánjìng, yòu jiéyuē zīyuán. Wǒ juéde měi ge rén dōu yīnggāi cānjiā zhè ge huódòng.)}\\
« J'habite à Yaoundé. Il y a beaucoup de déchets dans notre ville, donc le recyclage est très important. D'abord, notre école enseigne le tri des déchets aux élèves. Ensuite, les camarades séparent le papier, le plastique et le verre. Enfin, des gens apportent les déchets triés au centre de recyclage. Ainsi, on protège à la fois l'environnement et on économise les ressources. Je pense que chacun devrait participer à cette activité. »
\end{exemplebox}

\begin{attention}[Erreurs fréquentes des francophones]
\begin{itemize}
\item Ne dis pas \zh{把垃圾放} : après la structure \zh{把}, le verbe doit être accompagné d'un élément (\zh{分开放}, \zh{放好}, \zh{送到}).
\item \zh{因为……所以……} : en chinois, on peut utiliser les deux ensemble dans la même phrase, contrairement au français (« parce que... donc... » est incorrect en français mais normal en chinois).
\item Attention aux tons : \zh{回收} huíshōu (2\up{e} + 1\up{er} ton), \zh{环境} huánjìng (2\up{e} + 4\up{e} ton), \zh{塑料} sùliào (4\up{e} + 4\up{e} ton).
\item Ne confonds pas \zh{既} (jì, « à la fois ») et \zh{即} ; et n'oublie pas que \zh{又} (yòu) seul signifie « de nouveau » : c'est la paire \zh{既……又……} qui exprime la simultanéité.
\item \zh{的} est obligatoire dans \zh{每个人的责任} : sans \zh{的}, la phrase est incorrecte.
\end{itemize}
\end{attention}

\newpage

\coursec{Corrigés des exercices}

\coursec{Corrigés des exercices — Leçon 11 : Expression écrite : le recyclage}

\begin{corrige}[Exercice 1 — QCM : le lexique du recyclage]
1. \textbf{b} — \zh{回收} (huíshōu) = recycler, récupérer. « Jeter » se dit \zh{扔} (rēng) ; « brûler » se dit \zh{烧} (shāo).\\
2. \textbf{b} — \zh{垃圾} (lājī) = les déchets. \zh{环境} (huánjìng) = l'environnement ; \zh{资源} (zīyuán) = les ressources.\\
3. \textbf{b} — \zh{塑料瓶} (sùliào píng) = bouteille en plastique ; une bouteille en verre serait \zh{玻璃瓶} (bōli píng).\\
4. \textbf{a} — \zh{保护} (bǎohù) = protéger. \zh{节约} (jiéyuē) = économiser ; \zh{分类} (fēnlèi) = trier, classer.\\
\emph{Point de vigilance :} mémorise ces mots par paires utiles (collocations) : \zh{保护环境} (protéger l'environnement), \zh{节约资源} (économiser les ressources), \zh{垃圾分类} (le tri des déchets), \zh{回收塑料瓶} (recycler les bouteilles en plastique).
\end{corrige}

\begin{corrige}[Exercice 2 — Vrai ou faux justifié]
1. \textbf{V} — \zh{垃圾} = déchets, \zh{分类} = trier, classer ; le composé signifie bien « le tri des déchets ».\\
2. \textbf{F} — \zh{把} (bǎ) ne marque pas le passé : c'est une préposition qui place le complément d'objet \emph{avant} le verbe (Structure : Sujet + \zh{把} + Objet + Verbe + Complément). L'aspect accompli se marque avec la particule \zh{了} (le) après le verbe.\\
3. \textbf{V} — \zh{既……又……} (jì……yòu……) coordonne deux qualités ou deux actions de même nature : « à la fois... et... ». Exemple : \zh{这样做既保护环境，又节约资源。} (Zhèyàng zuò jì bǎohù huánjìng, yòu jiéyuē zīyuán.) « Ainsi, on protège à la fois l'environnement et on économise les ressources. »\\
4. \textbf{F} — Le radical (clé) de \zh{环} (huán) est la clé du jade \zh{王} (à gauche), non la clé de la terre \zh{土}. Attention à ne pas confondre \zh{环} (huán, anneau, environnement) et \zh{坏} (huài, mauvais, cassé), qui lui contient bien la clé de la terre \zh{土}.
\end{corrige}

\begin{corrige}[Exercice 3 — Vocabulaire : complète le tableau]
\begin{center}
\begin{tabularx}{\linewidth}{|X|X|X|X|}
\hline
\rowcolor{popblueL} Caractères & Pinyin & Nature & Traduction \\
\hline
\zh{环境} & huánjìng & nom & l'environnement \\
\hline
\zh{资源} & zīyuán & nom & les ressources \\
\hline
\zh{玻璃} & bōli & nom & le verre (matière) \\
\hline
\zh{节约} & jiéyuē & verbe & économiser \\
\hline
\zh{污染} & wūrǎn & nom / verbe & la pollution / polluer \\
\hline
\end{tabularx}
\end{center}
\emph{Point de vigilance :} \zh{玻璃} s'écrit avec deux caractères portant tous deux la clé du jade \zh{王} ; \zh{污染} contient la clé de l'eau \zh{氵} dans chaque caractère — logique, puisque l'eau est souvent le milieu pollué.
\end{corrige}

\begin{corrige}[Exercice 4 — Associe le pinyin aux caractères]
1 — c (\zh{回收} huíshōu) ; 2 — a (\zh{分类} fēnlèi) ; 3 — b (\zh{塑料} sùliào) ; 4 — e (\zh{保护} bǎohù) ; 5 — d (\zh{垃圾} lājī).\\
\emph{Astuce :} dans \zh{塑料}, le premier caractère \zh{塑} (sù, 4ᵉ ton) contient la clé de la terre \zh{土} en bas (le plastique est un matériau « terre » synthétique) ; dans \zh{保护}, les deux caractères portent le 3ᵉ ton à l'écrit, mais à l'oral le premier 3ᵉ ton se prononce comme un 2ᵉ ton (règle du 3ᵉ ton + 3ᵉ ton) : \emph{báohù}.
\end{corrige}

\begin{corrige}[Exercice 5 — Discrimination de caractères proches]
1. \zh{垃圾} (lājī) — \zh{圾} a la clé de la terre \zh{土} (déchets).\\
2. \zh{环境} (huánjìng) — \zh{环} a la clé du jade \zh{王} (anneau, cercle).\\
3. \zh{收到} (shōudào, recevoir) — ne pas confondre avec \zh{放下} (fàngxià, poser).\\
4. \zh{年级} (niánjí, niveau/classe) — \zh{级} a la clé de la soie \zh{纟} (rang, échelon).\\
5. \zh{破坏} (pòhuài, détruire) — \zh{坏} a la clé de la terre \zh{土} (casser, gâter).\\
6. \zh{放下} (fàngxià, poser / lâcher).\\
\emph{Retiens :} \zh{收} (shōu, recevoir, ramasser, clé \zh{攵}) et \zh{放} (fàng, poser, lâcher, clé \zh{方} à gauche) sont antonymes. De même, \zh{圾} (clé \zh{土}, ordures) et \zh{级} (clé \zh{纟}, rang) ne sont pas interchangeables malgré une prononciation proche (jī vs jí).
\end{corrige}

\begin{corrige}[Exercice 6 — Texte à trous : mots-outils et connecteurs]
Texte complété :\\
\zh{因为我们只有一个地球，所以我们要保护环境。我们把垃圾分类好，既能节约资源，又能减少污染。这是每个人的责任。}\\
\emph{(Yīnwèi wǒmen zhǐ yǒu yí ge dìqiú, suǒyǐ wǒmen yào bǎohù huánjìng. Wǒmen bǎ lājī fēnlèi hǎo, jì néng jiéyuē zīyuán, yòu néng jiǎnshǎo wūrǎn. Zhè shì měi ge rén de zérèn.)}\\
Traduction : « Parce que nous n'avons qu'une seule Terre, nous devons protéger l'environnement. Nous trions bien les déchets : cela permet à la fois d'économiser les ressources et de réduire la pollution. C'est la responsabilité de chacun. »\\
\emph{Points de langue :}
\begin{itemize}
\item \zh{因为……所以……} (yīnwèi… suǒyǐ…) introduit la cause puis la conséquence (les deux peuvent coexister en chinois, contrairement au français « parce que / donc » redondant).
\item \zh{把} (bǎ) place l'objet \zh{垃圾} avant le verbe \zh{分类} ; \zh{好} (hǎo) est un complément de résultat signifiant « correctement, bien fait ».
\item \zh{既……又……} (jì… yòu…) coordonne deux verbes (\zh{节约} et \zh{减少}) partageant le même sujet implicite.
\item \zh{的} (de) marque l'appartenance dans \zh{每个人的责任} « la responsabilité de chacun ».
\end{itemize}
\end{corrige}

\begin{corrige}[Exercice 7 — Remise en ordre et traduction]
L'ordre logique suit la chronologie des connecteurs : \zh{首先} (d'abord) → \zh{然后} (ensuite) → \zh{最后} (enfin) → conclusion \zh{这样做} (en faisant ainsi).\\
Texte reconstitué :\\
\zh{首先，我们学习垃圾分类的知识。然后，我们把纸、塑料和玻璃分开放。最后，我们把分好的垃圾送到回收站。这样做既保护环境，又节约资源。}\\
\emph{(Shǒuxiān, wǒmen xuéxí lājī fēnlèi de zhīshi. Ránhòu, wǒmen bǎ zhǐ, sùliào hé bōli fēnkāi fàng. Zuìhòu, wǒmen bǎ fēn hǎo de lājī sòng dào huíshōuzhàn. Zhèyàng zuò jì bǎohù huánjìng, yòu jiéyuē zīyuán.)}\\
Traduction : « D'abord, nous apprenons les connaissances sur le tri des déchets. Ensuite, nous mettons le papier, le plastique et le verre séparément. Enfin, nous apportons les déchets triés au centre de recyclage. Ainsi, on protège à la fois l'environnement et on économise les ressources. »\\
\emph{Remarque :} la phrase commençant par \zh{这样做} « en faisant ainsi / c'est ainsi que » est une phrase de conclusion : elle se place donc après la séquence d'actions. Les connecteurs \zh{首先} → \zh{然后} → \zh{最后} forment la trame standard d'un texte procédural (récit d'opérations successives) : réutilise-les dans tes propres rédactions.
\end{corrige}

\begin{corrige}[Exercice 8 — Écriture des caractères du thème : radicaux et ordre des traits]
Radicaux (clés) à identifier et ordre d'écriture principal :
\begin{itemize}
\item \zh{回} (huí) : clé de l'enceinte \zh{囗} (wéi). Ordre : trait vertical gauche → haut et droite du cadre (un seul trait \zh{乛}) → \zh{口} intérieur → trait horizontal de fermeture (bas).
\item \zh{收} (shōu) : clé du couteau-frappe \zh{攵} (pū, à droite). Ordre : partie gauche (\zh{又} + trait) → partie droite \zh{攵} (point, horizontal, point, trait oblique).
\item \zh{环} (huán) : clé du jade \zh{王} (wáng, à gauche). \zh{境} (jìng) : clé de la terre \zh{土} (tǔ, à gauche). Règle gauche-droite : partie gauche d'abord.
\item \zh{塑} (sù) : clé de la terre \zh{土} (tǔ, en bas). \zh{料} (liào) : clé du riz \zh{米} (mǐ, à gauche). Règle haut-bas : partie haute d'abord ; gauche-droite : gauche d'abord.
\item \zh{分} (fēn) : clé du couteau \zh{刀} (dāo, en bas). \zh{类} (lèi) : clé du riz \zh{米} (mǐ, en haut). Règle haut-bas : partie haute d'abord.
\end{itemize}
\emph{Méthode générale (règles de base) :} de haut en bas (上→下), de gauche à droite (左→右), l'horizontal avant le vertical (横→竖), on ferme les cadres en dernier (框最后闭口). Pour les caractères composés, on écrit généralement le composant principal (souvent la clé) en premier s'il est en haut ou à gauche.
\end{corrige}

\begin{corrige}[Exercice 9 — Compréhension de l'écrit : texte et questions]
1. Beaucoup d'écoles enseignent aux élèves le tri des déchets (\zh{垃圾分类} lājī fēnlèi).\\
2. Les élèves séparent le papier (\zh{纸} zhǐ), le plastique (\zh{塑料} sùliào) et le verre (\zh{玻璃} bōli).\\
3. Les deux avantages cités sont : protéger l'environnement (\zh{保护环境} bǎohù huánjìng) et économiser les ressources (\zh{节约资源} jiéyuē zīyuán).\\
4. Phrase du texte : \zh{同学们把纸、塑料和玻璃分开放。}\\
   Analyse structurelle : Sujet \zh{同学们} + \zh{把} + Objet \zh{纸、塑料和玻璃} + Verbe \zh{分开放} (verbe \zh{放} précédé du complément de manière \zh{分开} « séparément »). \textbf{Règle cruciale :} dans la structure \zh{把}, le verbe ne doit \textbf{jamais} rester seul ; il doit être suivi d'un complément (résultat, direction, durée, manière, particule \zh{了}, etc.).\\
5. \zh{既……又……} (jì……yòu……) signifie « à la fois... et... » / « aussi bien... que... ». Les deux éléments coordonnés doivent être de même nature grammaticale (deux verbes, deux adjectifs, deux noms).\\
   Exemple personnel : \zh{回收既减少垃圾，又保护资源。} (Huíshōu jì jiǎnshǎo lājī, yòu bǎohù zīyuán.) « Recycler réduit à la fois les déchets et protège les ressources. »
\end{corrige}

\begin{corrige}[Exercice 10 — Thème et version type Bac]
\textbf{A. Thème (Français → Chinois).}
\begin{enumerate}
\item \zh{我们应该垃圾分类，回收塑料瓶。} \emph{(Wǒmen yīnggāi lājī fēnlèi, huíshōu sùliào píng.)} — Variante acceptée avec \zh{把} : \zh{我们应该把垃圾分类好，还要回收塑料瓶。} (Wǒmen yīnggāi bǎ lājī fēnlèi hǎo, hái yào huíshōu sùliào píng.)
\item \zh{因为我们只有一个地球，所以我们要保护它。} \emph{(Yīnwèi wǒmen zhǐ yǒu yí ge dìqiú, suǒyǐ wǒmen yào bǎohù tā.)} — \zh{只} (zhǐ) = « seulement / ne ... que » ; le pronom \zh{它} (tā) reprend \zh{地球} (objet inanimé).
\item \zh{我们中学的学生把纸放进蓝色的箱子里。} \emph{(Wǒmen zhōngxué de xuésheng bǎ zhǐ fàng jìn lánsè de xiāngzi lǐ.)} — Structure \zh{把} : Sujet + \zh{把} + Objet \zh{纸} + Verbe \zh{放} + Complément de direction \zh{进……里} (jìn… lǐ, « à l'intérieur de »). \zh{蓝色的} (lánse de) = bleu (couleur de la poubelle à papier en Chine).
\end{enumerate}

\textbf{B. Version (Chinois → Français).}
Texte source (implicite) : \zh{为了保护环境，我们学校每个月回收纸和塑料瓶。同学们觉得这个活动很有意义。}\\
Traduction attendue : « Pour protéger l'environnement, notre école recycle chaque mois du papier et des bouteilles en plastique. Les camarades trouvent que cette activité est très utile (pleine de sens). »\\
\emph{Points de langue :} \zh{为了} (wèile) + verbe = « pour + infinitif » (but) ; \zh{每个月} (měi ge yuè) = « chaque mois » (pas de \zh{的} après \zh{每}) ; \zh{觉得} (juéde) = « trouver, penser que » (avis subjectif) ; \zh{有意义} (yǒu yìyì) = « avoir du sens, être significatif / utile ».\\
\emph{Barème indicatif (sur 10) :} Thème 3 × 2 pts (lexique 1 pt + grammaire/structure 1 pt par phrase) ; Version 4 pts (fidélité du sens 2 pts, qualité du français 2 pts).
\end{corrige}

\begin{corrige}[Exercice 11 — Expression écrite type Bac : rédaction guidée]
\textbf{Proposition de rédaction modèle (environ 90 caractères, 7 phrases) :}\\
\zh{我住在雅温得。我们城市的垃圾很多，所以回收非常重要。首先，我们学校教学生垃圾分类。然后，同学们把纸、塑料和玻璃分开放。最后，有人把分好的垃圾送到回收站。这样做既保护环境，又节约资源。我觉得每个人都应该参加这个活动。}\\
\emph{(Wǒ zhù zài Yǎwēndé. Wǒmen chéngshì de lājī hěn duō, suǒyǐ huíshōu fēicháng zhòngyào. Shǒuxiān, wǒmen xuéxiào jiāo xuésheng lājī fēnlèi. Ránhòu, tóngxuémen bǎ zhǐ, sùliào hé bōli fēnkāi fàng. Zuìhòu, yǒu rén bǎ fēn hǎo de lājī sòng dào huíshōuzhàn. Zhèyàng zuò jì bǎohù huánjìng, yòu jiéyuē zīyuán. Wǒ juéde měi ge rén dōu yīnggāi cānjiā zhè ge huódòng.)}\\
Traduction : « J'habite à Yaoundé. Il y a beaucoup de déchets dans notre ville, donc le recyclage est très important. D'abord, notre école enseigne le tri des déchets aux élèves. Ensuite, les camarades séparent le papier, le plastique et le verre. Enfin, des gens apportent les déchets triés au centre de recyclage. Ainsi, on protège à la fois l'environnement et on économise les ressources. Je pense que chacun devrait participer à cette activité. »\\

\textbf{Vérification des consignes (auto-évaluation) :}
\begin{itemize}
\item \textbf{Longueur / Structure :} 7 phrases, ~90 caractères, paragraph unique.
\item \textbf{Connecteurs logiques :} \zh{所以} (conséquence), \zh{首先} / \zh{然后} / \zh{最后} (chronologie), \zh{既……又……} (double avantage).
\item \textbf{Structures grammaticales ciblées :} Deux structures \zh{把} correctes (\zh{把纸、塑料和玻璃分开放}, \zh{把分好的垃圾送到回收站}) ; verbe jamais seul après \zh{把}.
\item \textbf{Lexique du thème :} \zh{垃圾}, \zh{回收}, \zh{分类}, \zh{塑料}, \zh{玻璃}, \zh{保护}, \zh{环境}, \zh{节约}, \zh{资源}, \zh{回收站}.
\item \textbf{Opinion personnelle :} \zh{我觉得每个人都应该参加这个活动} (structure \zh{觉得} + proposition ; \zh{应该} obligation morale).
\end{itemize}

\textbf{Barème indicatif (sur 10) :}
\begin{itemize}
\item Correction des caractères (forme, traits, radicaux) : 2 pts.
\item Grammaire et ordre des mots (structure \zh{把}, compléments, mots-outils \zh{的/地/得}, négation) : 3 pts.
\item Respect des consignes (connecteurs, structures imposées, nombre de phrases) : 3 pts.
\item Richesse et justesse du vocabulaire (variété, collocations) : 2 pts.
\end{itemize}

\textbf{Points de vigilance (pièges fréquents) :}
\begin{itemize}
\item \textbf{Verbe nu après \zh{把} :} \zh{把垃圾分类} $\times$ → \zh{把垃圾分类好} / \zh{把垃圾分开} ✓.
\item \textbf{Ponctuation :} virgule chinoise \zh{，} après \zh{首先}, \zh{然后}, \zh{最后}, \zh{所以} ; point final \zh{。}.
\item \textbf{Caractères proches :} ne pas confondre \zh{环} (huán, clé \zh{王}) / \zh{坏} (huài, clé \zh{土}) ; \zh{圾} (jī, clé \zh{土}) / \zh{级} (jí, clé \zh{纟}).
\item \textbf{Particule \zh{的} :} obligatoire dans \zh{我们城市\zh{的}垃圾}, \zh{蓝色\zh{的}箱子}, \zh{分好\zh{的}垃圾}, \zh{这个\zh{的}活动} (ou \zh{这项活动}).
\end{itemize}
\end{corrige}

\newpage

\coursec{Fiche bilan}

\begin{popfigure}
\begin{center}
\begin{center}\tcbox[colback=popblue,colframe=popblue,coltext=white,arc=9pt,boxsep=4pt,left=6pt,right=6pt]{\bfseries\small Leçon 11 Expression écrite le recyclage \zh{回收}}\end{center}
\vspace{-2pt}
\begin{tcbraster}[raster columns=3,raster equal height=rows,raster column skip=6pt,raster row skip=6pt,enhanced,blankest]
\begin{tcolorbox}[enhanced,colback=poporange,colframe=poporange,coltext=white,boxrule=0.9pt,arc=3pt,left=4pt,right=4pt,top=3pt,bottom=3pt,boxsep=1pt,halign=left]\bfseries\scriptsize Lexique clé \zh{垃圾分类} \zh{保护环境}\end{tcolorbox}
\begin{tcolorbox}[enhanced,colback=popgreen,colframe=popgreen,coltext=white,boxrule=0.9pt,arc=3pt,left=4pt,right=4pt,top=3pt,bottom=3pt,boxsep=1pt,halign=left]\bfseries\scriptsize Structure du texte introduction corps conclusion\end{tcolorbox}
\begin{tcolorbox}[enhanced,colback=poppurple,colframe=poppurple,coltext=white,boxrule=0.9pt,arc=3pt,left=4pt,right=4pt,top=3pt,bottom=3pt,boxsep=1pt,halign=left]\bfseries\scriptsize Connecteurs \zh{首先} \zh{然后} \zh{最后} \zh{因为}\end{tcolorbox}
\begin{tcolorbox}[enhanced,colback=popgold,colframe=popgold,coltext=white,boxrule=0.9pt,arc=3pt,left=4pt,right=4pt,top=3pt,bottom=3pt,boxsep=1pt,halign=left]\bfseries\scriptsize Grammaire passif \zh{被} \zh{把} + objet verbes modaux\end{tcolorbox}
\begin{tcolorbox}[enhanced,colback=poppink,colframe=poppink,coltext=white,boxrule=0.9pt,arc=3pt,left=4pt,right=4pt,top=3pt,bottom=3pt,boxsep=1pt,halign=left]\bfseries\scriptsize Caractères radicaux ordre des traits\end{tcolorbox}
\begin{tcolorbox}[enhanced,colback=popblue!60,colframe=popblue!60,coltext=white,boxrule=0.9pt,arc=3pt,left=4pt,right=4pt,top=3pt,bottom=3pt,boxsep=1pt,halign=left]\bfseries\scriptsize Traduction thème version\end{tcolorbox}
\end{tcbraster}

\end{center}
\legende{Carte mentale de la leçon : écrire un texte simple sur le recyclage}
\end{popfigure}

\begin{definition}[Lexique clé du thème : à connaître par cœur]
\begin{center}
\begin{tabularx}{\linewidth}{|X|X|X|X|}
\hline
\rowcolor{popblueL} Caractères & Pinyin & Nature & Traduction \\
\hline
回收 & huíshōu & verbe / nom & recycler, recyclage \\
\hline
垃圾 & lājī & nom & déchets, ordures \\
\hline
分类 & fēnlèi & verbe / nom & trier, classer ; tri, classification \\
\hline
环境 & huánjìng & nom & environnement \\
\hline
保护 & bǎohù & verbe & protéger \\
\hline
污染 & wūrǎn & nom / verbe & pollution, polluer \\
\hline
瓶子 & píngzi & nom & bouteille \\
\hline
纸 & zhǐ & nom & papier \\
\hline
塑料 & sùliào & nom & plastique \\
\hline
节约 & jiéyuē & verbe & économiser \\
\hline
减少 & jiǎnshǎo & verbe & réduire, diminuer \\
\hline
资源 & zīyuán & nom & ressources \\
\hline
严重 & yánzhòng & adjectif & grave, sérieux \\
\hline
一起 & yìqǐ & adverbe & ensemble \\
\hline
\end{tabularx}
\end{center}
\emph{Remarque de prononciation :} \zh{垃圾} se lit \emph{lājī} (et non « lèsè » : cette lecture existe mais n'est pas la norme du putonghua enseignée au lycée). Le classificateur des bouteilles est \zh{个} ou \zh{只} : \zh{一个瓶子} (\emph{yí ge píngzi}, « une bouteille »).
\end{definition}

\begin{propriete}[Grammaire à connaître par cœur]
\begin{center}
\begin{tabularx}{\linewidth}{|X|X|X|}
\hline
\rowcolor{popblueL} Structure & Exemple & Traduction \\
\hline
Sujet + 把 + objet + verbe + complément & 我们把垃圾进行分类。Wǒmen bǎ lājī jìnxíng fēnlèi. & Nous trions les déchets. \\
\hline
Objet + 被 + agent + verbe (+ 了) & 瓶子被回收了。Píngzi bèi huíshōu le. & Les bouteilles ont été recyclées. \\
\hline
应该 / 可以 + verbe & 我们应该保护环境。Wǒmen yīnggāi bǎohù huánjìng. & Nous devons protéger l'environnement. \\
\hline
越来越 + adjectif & 污染越来越严重。Wūrǎn yuè lái yuè yánzhòng. & La pollution est de plus en plus grave. \\
\hline
为了 + but，+ proposition & 为了保护环境，我们要回收垃圾。Wèile bǎohù huánjìng, wǒmen yào huíshōu lājī. & Pour protéger l'environnement, nous devons recycler les déchets. \\
\hline
\end{tabularx}
\end{center}
\end{propriete}

\begin{attention}[Pièges fréquents des francophones]
\begin{itemize}
  \item Avec \zh{把}, le verbe ne peut presque jamais rester « nu » : il faut un complément, un résultat ou \zh{了}. On évite donc \zh{我们把垃圾分类。} seul ; on dit \zh{我们把垃圾进行分类。} ou \zh{我们把垃圾分好了。} (\emph{Wǒmen bǎ lājī fēn hǎo le.}, « Nous avons bien trié les déchets. »).
  \item \zh{被} marque le passif, souvent avec une nuance défavorable ; à l'écrit, on peut omettre l'agent : \zh{旧纸被回收了。} (\emph{Jiù zhǐ bèi huíshōu le.}, « Le vieux papier a été recyclé. »).
  \item \zh{了} se place après le verbe pour une action accomplie, jamais après un adjectif seul : on dit \zh{污染很严重。} et non \emph{污染很严重了} pour décrire un état (avec \zh{了}, le sens devient « est devenue grave »).
  \item Les compléments de temps et de lieu se placent avant le verbe : \zh{我们每天在学校分类垃圾。} (\emph{Wǒmen měitiān zài xuéxiào fēnlèi lājī.}, « Chaque jour, nous trions les déchets à l'école. »).
\end{itemize}
\end{attention}

\begin{definition}[Connecteurs du texte argumentatif simple]
\begin{itemize}
  \item \zh{首先} (shǒuxiān) : d'abord ; \zh{然后} (ránhòu) : ensuite ; \zh{最后} (zuìhòu) : enfin.
  \item \zh{因为……所以……} (yīnwèi... suǒyǐ...) : parce que... donc... (en chinois, on peut garder les deux dans la même phrase, contrairement au français).
  \item \zh{但是} (dànshì) : mais ; \zh{而且} (érqiě) : de plus ; \zh{如果} (rúguǒ) : si ; \zh{为了} (wèile) : pour, afin de.
\end{itemize}
\emph{Exemple :} \zh{因为塑料污染很严重，所以我们应该少用塑料袋。} (\emph{Yīnwèi sùliào wūrǎn hěn yánzhòng, suǒyǐ wǒmen yīnggāi shǎo yòng sùliào dài.}) « Parce que la pollution plastique est très grave, nous devons utiliser moins de sacs plastiques. »
\end{definition}

\begin{methode}[Méthodes express pour l'expression écrite et la traduction]
\begin{enumerate}
  \item \textbf{Plan en 3 parties} : introduction (le problème : \zh{现在污染很严重。}), corps (les actions : \zh{我们应该……首先……然后……}), conclusion (l'appel : \zh{让我们一起保护环境！}).
  \item \textbf{Thème (français $\rightarrow$ chinois)} : repérer le sujet ; choisir la structure (\zh{把} si l'objet est mis en avant, \zh{被} si le français est au passif) ; vérifier l'ordre Sujet + Temps/Lieu + Verbe ; ajouter \zh{了} si l'action est accomplie.
  \item \textbf{Version (chinois $\rightarrow$ français)} : traduire d'abord les groupes (sujet / verbe / compléments), puis reformuler en français naturel ; attention au passif \zh{被} = « être + participe passé ».
  \item \textbf{Caractères} : écrire de haut en bas, de gauche à droite, l'horizontal avant le vertical, l'extérieur avant l'intérieur (puis fermer). Repérer les clés : \zh{土} (terre) dans \zh{垃} et \zh{圾}, \zh{纟} (soie) dans \zh{纸}, \zh{氵} (eau) dans \zh{污} et \zh{染}, \zh{囗} (enceinte) dans \zh{回}.
\end{enumerate}
\end{methode}

\begin{attention}[Caractères proches à ne pas confondre]
\begin{itemize}
  \item \zh{纸} (zhǐ, papier, clé \zh{纟}) $\neq$ \zh{低} (dī, bas, clé \zh{亻}) : la clé change tout le sens.
  \item \zh{圾} (jī, dans \zh{垃圾}) $\neq$ \zh{级} (jí, niveau, clé \zh{纟}) : mêmes sons proches, clés différentes.
  \item \zh{回} (huí, retourner) s'écrit avec la grande enceinte \zh{囗} autour du petit \zh{口} : on écrit l'extérieur, puis l'intérieur, puis on ferme le trait du bas.
  \item \zh{染} (rǎn, dans \zh{污染}) contient \zh{氵} + \zh{九} + \zh{木} : ne pas écrire \zh{丸} à la place de \zh{九}.
\end{itemize}
\end{attention}

\begin{aretenir}[Checklist avant le Bac]
\begin{itemize}
  \item[$\square$] Je connais par cœur le lexique du recyclage (\zh{回收}, \zh{垃圾}, \zh{分类}, \zh{环境}, \zh{保护}, \zh{污染}).
  \item[$\square$] Je sais construire une phrase avec \zh{把} : Sujet + \zh{把} + objet + verbe + complément.
  \item[$\square$] Je sais transformer une phrase active en passive avec \zh{被}.
  \item[$\square$] J'utilise correctement les verbes modaux \zh{应该}, \zh{可以}, \zh{要} avant le verbe.
  \item[$\square$] Je place les compléments de temps et de lieu avant le verbe.
  \item[$\square$] Je sais organiser mon texte avec \zh{首先}, \zh{然后}, \zh{最后}.
  \item[$\square$] Je sais exprimer la cause avec \zh{因为……所以……} et le but avec \zh{为了}.
  \item[$\square$] Je mets \zh{了} après le verbe pour une action accomplie, jamais après un adjectif seul.
  \item[$\square$] J'écris les caractères du thème dans le bon ordre des traits, sans confondre \zh{纸} et \zh{低}, \zh{圾} et \zh{级}.
  \item[$\square$] En version, je repère le passif \zh{被} et je le traduis par « être + participe passé ».
  \item[$\square$] Je relis mon texte : tons du pinyin, ponctuation chinoise pleine chasse （， 。 ！）, aucun caractère inventé.
  \item[$\square$] Je sais conclure par un appel à l'action : \zh{让我们一起保护环境！} (\emph{Ràng wǒmen yìqǐ bǎohù huánjìng!}, « Protégeons ensemble l'environnement ! »).
\end{itemize}
\end{aretenir}

\findecours

\end{document}
